% ---- hand chunk preamble
\documentclass[11pt]{article}
\usepackage[T1]{fontenc}
\usepackage[utf8]{inputenc}
\usepackage{lmodern}
\usepackage{amsmath,amssymb,amsthm,mathtools}
\usepackage{mathrsfs}
\usepackage[margin=1in,headheight=14pt]{geometry}
\usepackage{microtype}
\usepackage{booktabs,longtable,array,tabularx,enumitem}
\usepackage{xcolor,fancyhdr,listings}
\usepackage[most]{tcolorbox}
\usepackage{aliascnt}
\usepackage{hyperref}
\usepackage{xurl}
\usepackage[nameinlink,noabbrev]{cleveref}
\definecolor{navy}{RGB}{25,54,88}
\definecolor{linkblue}{RGB}{26,77,127}
\definecolor{Ink}{HTML}{183247}
\definecolor{Accent}{HTML}{126E75}
\definecolor{Pale}{HTML}{EDF5F5}
\definecolor{Gray}{HTML}{52616B}
\hypersetup{colorlinks=true,linkcolor=linkblue,citecolor=linkblue,urlcolor=linkblue,bookmarksdepth=2,
  pdftitle={Lexicographic Orders of Well-Orderings of the Surreal Numbers},
  pdfauthor={Research report merged from four manuscripts prepared for Vladimir Reshetnikov (sources 08, 09, 11 and 12 of batch 80)},
  pdfsubject={Surreal collection, foundations and computation: surreal-well-orders}}
\setcounter{tocdepth}{1}
\urlstyle{same}
\setlength{\emergencystretch}{3em}
\setlength{\headheight}{14pt}
\pagestyle{fancy}
\fancyhf{}
\fancyhead[L]{\small\textsc{Well-orders of the surreal numbers}}
\fancyhead[R]{\small\thepage}
\renewcommand{\headrulewidth}{0.3pt}
\setlist{topsep=5pt,itemsep=3pt,parsep=0pt}
\newtheorem{theorem}{Theorem}[section]
\newaliascnt{lemma}{theorem}
\newtheorem{lemma}[lemma]{Lemma}
\aliascntresetthe{lemma}
\newaliascnt{proposition}{theorem}
\newtheorem{proposition}[proposition]{Proposition}
\aliascntresetthe{proposition}
\newaliascnt{corollary}{theorem}
\newtheorem{corollary}[corollary]{Corollary}
\aliascntresetthe{corollary}
\theoremstyle{definition}
\newaliascnt{definition}{theorem}
\newtheorem{definition}[definition]{Definition}
\aliascntresetthe{definition}
\newaliascnt{example}{theorem}
\newtheorem{example}[example]{Example}
\aliascntresetthe{example}
\newaliascnt{question}{theorem}
\newtheorem{question}[question]{Research question}
\aliascntresetthe{question}
\theoremstyle{remark}
\newaliascnt{remark}{theorem}
\newtheorem{remark}[remark]{Remark}
\aliascntresetthe{remark}
\newaliascnt{note}{theorem}
\newtheorem{note}[note]{Note}
\aliascntresetthe{note}
\crefname{lemma}{lemma}{lemmas}
\Crefname{lemma}{Lemma}{Lemmas}
\crefname{proposition}{proposition}{propositions}
\Crefname{proposition}{Proposition}{Propositions}
\crefname{corollary}{corollary}{corollaries}
\Crefname{corollary}{Corollary}{Corollaries}
\crefname{definition}{definition}{definitions}
\Crefname{definition}{Definition}{Definitions}
\crefname{example}{example}{examples}
\Crefname{example}{Example}{Examples}
\crefname{question}{research question}{research questions}
\Crefname{question}{Research question}{Research questions}
\crefname{remark}{remark}{remarks}
\Crefname{remark}{Remark}{Remarks}
\crefname{note}{note}{notes}
\Crefname{note}{Note}{Notes}
\crefname{part}{part}{parts}
\Crefname{part}{Part}{Parts}
\newtcolorbox{resultbox}{colback=Pale,colframe=Accent,boxrule=.6pt,
 arc=1mm,left=3mm,right=3mm,top=2mm,bottom=2mm,breakable}
\newcolumntype{Y}{>{\raggedright\arraybackslash}X}
% ---- one notation for all four sources (Section 1.4 lists every renaming)
\newcommand{\No}{\mathbf{No}}
\newcommand{\Ord}{\mathrm{Ord}}
\newcommand{\WO}{\mathcal W}
\newcommand{\W}{\mathcal W}
\newcommand{\Words}{\mathcal I}
\newcommand{\I}{\mathcal I}
\newcommand{\WordsAll}{\mathcal T}
\newcommand{\Hsl}{\mathfrak W_{\mathrm{sl}}}
\newcommand{\Hall}{\mathfrak W_{\mathrm{all}}}
\newcommand{\Jinj}{\mathfrak J}
\newcommand{\Pbd}{\mathcal P_{\mathrm{bd}}}
\newcommand{\Psupp}{\mathcal P_{\mathrm{s}}}
\newcommand{\B}{\mathbb B}
\newcommand{\BB}{\mathbb B}
\newcommand{\BOrd}{\mathbb B_{\Ord}}
\newcommand{\two}{\{0,1\}}
\newcommand{\lex}{\mathrm{lex}}
\newcommand{\ran}{\operatorname{ran}}
\newcommand{\dom}{\operatorname{dom}}
\newcommand{\otp}{\operatorname{otp}}
\newcommand{\cf}{\operatorname{cf}}
\newcommand{\rk}{\operatorname{rank}}
\newcommand{\bl}{\ell_2}
\newcommand{\bd}{\operatorname{bd}}
\newcommand{\lh}{\operatorname{lh}}
\newcommand{\pred}{\operatorname{Pred}}
\newcommand{\supp}{\operatorname{supp}}
\newcommand{\Cut}{\operatorname{Cut}}
\newcommand{\Bij}{\operatorname{Bij}}
\newcommand{\Sym}{\operatorname{Sym}}
\newcommand{\id}{\mathrm{id}}
\newcommand{\op}{\mathrm{op}}
\newcommand{\End}{\mathtt{end}}
\newcommand{\emb}{\hookrightarrow}
\newcommand{\equi}{\mathrel{\simeq_{\!e}}}
\newcommand{\equim}{\mathrel{\simeq_{\!e}}}
\newcommand{\cat}{\mathbin{{}^\frown}}
\newcommand{\concat}{\mathbin{{}^\frown}}
\newcommand{\restr}{\mathbin{\upharpoonright}}
\newcommand{\res}{\mathbin{\upharpoonright}}
\newcommand{\GB}{\mathsf{GB}}
\newcommand{\GBC}{\mathsf{GBC}}
\newcommand{\ZFC}{\mathsf{ZFC}}
\newcommand{\ZF}{\mathsf{ZF}}
\newcommand{\KM}{\mathsf{KM}}
\newcommand{\ETR}{\mathsf{ETR}}
\newcommand{\GC}{\mathsf{GC}}
\newcommand{\card}[1]{\lvert #1\rvert}
\newcommand{\code}[1]{\texttt{\small #1}}
\newcommand{\merge}{\textsf{[merge]}}
\lstset{basicstyle=\small\ttfamily,columns=fullflexible,keepspaces=true,
  breaklines=true,frame=single,rulecolor=\color{black!20},
  backgroundcolor=\color{black!2},showstringspaces=false}

\title{\color{navy}\bfseries Lexicographic Orders of Well-Orderings\\
of the Surreal Numbers\\[5pt]
\large Set approximations, the singular birthday transition, set-length words,\\
class orders, and higher-order structure}
\author{Research report merged from four independent manuscripts\\
prepared for Vladimir Reshetnikov (sources 08, 09, 11 and 12 of batch 80)}
\date{Sources dated 2 October 2026 (source 11: 3 October 2026, UTC); merged 2 October 2026}

\begin{document}
\maketitle
\thispagestyle{empty}
\begin{abstract}

% ---- source 11, lines 82-113
The phrase \emph{all well-orders of the surreal numbers} has several
inequivalent meanings. Well-orders of a fixed birthday fragment form a set;
well-orders of arbitrary set-sized subsets form an ordinary class;
well-orders of the whole surreal line are themselves classes, and their
totality belongs to a higher logical level. We develop these settings and
distinguish set-like from unrestricted class well-orders.

For an infinite set-sized alphabet \(X\), the existing ProveIt theory gives
\(\card{\WO(X)}=2^{\card X}\) and the exact ordinal embedding spectrum
\(\beta<(\card X)^+\). We apply it to surreal birthday fragments. For
arbitrary set-length surreal words, we give an explicit sign-sequence
embedding into \(\No\), with birthday bounds, and identify the
one-sided cuts created by limit-length words.

For the whole surreal line, we define lexicographic comparison directly on
labeled class well-orders, without assuming ordinal enumerations or a class
of class order types. Every such well-order has a canonical initial segment
of type \(\Ord\); its remaining tail determines a convex fiber. In the
set-like fragment, bounded-support permutations of one global enumeration
form an ordinary class isomorphic to \(\No\), and interpolate every
uniformly set-indexed strict cut of the entire fragment. The fragment
contains every class linear order, admits binary class-word coding, and
cannot have an exhaustive uniform class of set-codes. The unrestricted
order has adjacent pairs, all arising from finite residual alphabets.
An omitted-range construction exhibits a proper-class cut filled by a
non-set-like well-order. In a full inaccessible-universe interpretation,
the unrestricted order does not embed into its set-like fragment.

The proofs use \(\ZF\), \(\ZFC\), or \(\GBC\), as specified, with
a further uniform nonembedding theorem in \(\KM+\GC\). We give
universe-relative and proof-assistant formulations, and propose further
research on the higher-order structure.

% ---- hand chunk abstract_merge@1

\merge\ This report is being merged from four manuscripts written
independently on 2 October 2026 on the same question and delivered together
in batch~80. This edition contains source~11 (the base, whose abstract is
printed above) and source~12; sources~09 and~08 follow in the second write.
Each shared theorem is printed once; source~12's versions of source~11's
results are recorded in numbered notes, and its genuinely different arguments
are kept as marked routes. Source~12 adds an exact dichotomy at birthday
cutoffs: for an infinite cardinal $\kappa$ and $\mu=2^{<\kappa}$, the
minimal-length enumeration order $\WO_\mu(\No_{<\kappa})$ is equimorphic with
the binary cube $\B_{\kappa\cdot\kappa}$ (ordinal product) exactly when
$\kappa$ is a singular strong limit, and with $\B_\mu$ otherwise, with
$\kappa=\beth_\omega$ as an unconditional example. It also adds a second
direct comparison of all class well-orders, by their largest common labelled
prefix, and an exact external identification of both global orders in a
full-class model at an inaccessible cardinal. Source~11's choiceless proof
answers source~12's question on the weak-choice strength, and source~12's
dichotomy answers in part a question of the repository's report on the reals.
The set-sized general-alphabet layer re-proves that report and is printed
with pointers to it, never as new.
\end{abstract}

% ---- hand chunk status_box
\begin{resultbox}

% ---- source 11, lines 116-121
\paragraph{Status.}
This is a mathematical research report with proofs, not a claim of
machine verification or bibliographic priority. Existing general
set-alphabet results are credited to the repository. The further
constructions are proved here and should be evaluated by their arguments.
The repository audit was targeted; no Lean build or remote change was made.

% ---- hand chunk status_box_end@1

\smallskip
\merge\ \textbf{Status of the merged report.} AI-assisted and unrefereed;
not formalized. Source~12 calls itself AI-assisted on its title page;
source~11 does not use the word (author lines: 11 ``Research report'', 12
``Research report prepared for Vladimir Reshetnikov''). Neither claims
historical priority, peer review or machine verification. Their finite checks
verify finite instances only. No statement of this report is formalized; the
formal status of the inputs it uses is in \cref{swo:sec:formal}.
\end{resultbox}

\clearpage
\tableofcontents
\clearpage

% ---- hand chunk about@1
\section{About this report}\label{swo:sec:about}

\subsection{Four manuscripts, one report}\label{swo:sub:sources}

\merge\ All four manuscripts of this report arrived in commit \code{4e270aa46}
and were placed in \code{ccc046989} (batch~80, cluster K3). This edition
merges source~11 (the base; pin \code{ae7c1e6aa}, 2 October 2026, 18:03:53
PDT = 3 October 2026, 01:03:53 UTC, which is the date its title page
gives; 2,807 lines, 35-page PDF, archive \code{surreal\_well\_orders (1).zip};
labels \code{swo:}, files \code{11-raw-orders-*}) and source~12 (pin
\code{63bb8b1d9}, 18:01:32 PDT; 1,689 lines, 24-page PDF, archive
\code{surreal\_well\_orders.zip}; labels \code{swo:st:}, files
\code{12-singular-*}). Sources~09 and~08 (labels \code{swo:ec:} and
\code{swo:sk:}, files \code{09-core-*} and \code{08-skeleton-*}) are merged in
the second write. The two archives sharing the name \code{surreal\_well\_orders}
are not editions of one another: word 8-gram overlap 0.7\,\%, no shared
member file, different pins and result sets.

Source~11 is the base. Its global-choice equivalence is proved over $\GB$
without set choice (\cref{swo:cf:choice}); it covers set alphabets,
set-length words, and set-like and unrestricted class well-orders; and it
attributes the general set-alphabet theorems to the repository's report on
the reals,
\path{SetTheory/Cardinals/docs/reports/ordinals-and-order-types/lexicographic-well-orderings-of-reals/}
(cited below as the reals report, labels \code{lwo:}).

\subsection{How the texts are printed}\label{swo:sub:printing}

\merge\ Source~11's text is printed in full and in its own order. Source~12's
introduction and foundational section are printed whole; a result that a
printed result already proves by the same argument is recorded in a numbered
\emph{Note} with the source's statement, number and differences; a genuinely
different proof is printed in full as a marked route; a result proved only by
source~12 is printed in full where it belongs. Every label survives with its
prefix; unlabelled numbered statements carry \code{swo:}\emph{sub}\code{:n}%
\emph{number}. \Cref{swo:app:crosswalk} maps every numbered statement to its
place here. Text written for the merge is marked \merge.

% ---- hand chunk notation@1
\subsection{One notation for the sources}\label{swo:sec:notation}

\merge\ The sources' symbols collide. The following convention is fixed
before any of them is used; every renaming is listed.

\begingroup\small
\begin{longtable}{@{}>{\raggedright\arraybackslash}p{0.24\textwidth}>{\raggedright\arraybackslash}p{0.70\textwidth}@{}}
\toprule
Here & Meaning, and the sources' symbols\\
\midrule\endhead
$\GB$, $\GBC$ &
$\GB$ has a $\ZF$ set part, elementary class comprehension and class
replacement, and \emph{no} choice principle (\code{found:sub:gbconvention});
$\GBC=\GB+\GC$. Source~12 analyses choice ``over GB plus set-level AC''
(\cref{swo:st:thm:choice}), written $\GB+\mathrm{AC}$ in the merge text;
source~11 works over $\GB$ itself.\\[3pt]
$\bd(x)$ &
Birthday. Source~11's $\ell(x)$ and source~12's $b(s)$ are renamed.\\[3pt]
$\No_{<\kappa}$ &
$\{x\in\No:\bd(x)<\kappa\}$. Renamed: source~11's $X_\kappa$, source~12's
$S_\kappa$ (the reals report's $S_{<\kappa}$, not $S_{<\kappa^+}$).\\[3pt]
$\bl(L)$ &
Binary coding length (\code{lwo:def:binary-length}); source~12 prints the
same symbol and calls it the ``binary coding rank''. $\rk(x)$ is the von
Neumann rank (source~11).\\[3pt]
$\WO(X)$, $\WO_\alpha(X)$ &
Exhaustive enumerations of a set $X$, and those of length $\alpha$; source~11's
$\mathcal P_\lambda(X)$ is renamed $\WO_\lambda(X)$.\\[3pt]
$\Hsl(\No)$, $\Hall(\No)$ &
Set-like and all class well-orders of $\No$, predicates on classes;
source~12's $\mathrm{SL}(\No)=\mathscr E(\No)$ and $\WO_{\mathrm{cl}}(\No)$.\\[3pt]
$\B_\gamma$, $\BOrd$ &
Binary cubes; class binary words on $\Ord$ (source~12's $\mathfrak B$).\\[3pt]
$\equi$ &
Equimorphism; source~11's $\approx_{\mathrm{emb}}$.\\[3pt]
$e_0$, $B$ & The baseline enumeration of sources~11 and~12.\\
\bottomrule
\end{longtable}
\endgroup

% ---- hand chunk claims@1
\subsection{What the report claims, and what it does not}\label{swo:sec:claims}

\merge\ The report claims the written proofs printed below, each in the
foundation its statement names. None is claimed to be historically first.
The set-sized general-alphabet layer re-proves the reals report and is
printed with pointers to it (\cref{swo:sub:lwo-pointer}); universality and
homogeneity of $\No$ are classical (Ehrlich \cite{Ehrlich2012}; Hamkins
\cite{HamkinsUniversal}).
\begin{itemize}
\item \emph{Source~11} (its Status paragraph, its delivery README's ``Scope
and verification'' and its audit): no machine verification, no Lean build, no
new Lean file; no bibliographic priority; research questions are proposed
directions; the audit was targeted; the finite checks do not verify
transfinite proofs, cardinal arithmetic or class comprehension. Its README
states: ``The mathematical proofs underwent a separate proof review.'' That
is the delivery's own statement; no such review is recorded in the
repository, and this report does not rely on it.
\item \emph{Source~12} (title page, its Section~1.2, and
\code{12-singular-RESEARCH\_STATUS.md}, ``What is not claimed''):
AI-assisted, unrefereed; priority not claimed; no catalogued open problem
claimed settled; no new Lean or Rocq verification; the Python checks are not
substitutes for the proofs; the bounded theorem is not a construction of a
full class model at a singular cardinal; the collection of class well-orders
is not treated as a set or class at the same level; no universal semantic
decoder is assumed; its research questions are proposals; its source search
was not sufficient to establish originality.
\item \emph{Review status.} The repository's Hilbert's-tenth research tree
holds a written-proof review of the four archives of this cluster
(\path{Computability/HilbertTenthProblem/Papers/research-wip/native-stream-queue/review_batch80_surreal.md},
commit \code{c338541e3}); it found their principal arguments sound under the
stated foundations, with one scoped correction to source~08, which joins in
the second write. That review is pinned to the delivered archives and does
not cover the text written for this merge; it is not an entry of the surreal
collection's review record, where this report's review is pending.
\end{itemize}

% ---- hand chunk formal@1
\subsection{Formal status}\label{swo:sec:formal}

\merge\ No statement of this report is formalized, and no \code{swo:} label
has a Lean mapping in the collection's formalization ledger
(\path{Algebra/SurrealNumbers/docs/FORMALIZATION.md}). Placement in the
surreal collection confers no formal status. Inputs the sources use are
proved in Lean in universe-relative form: properness and bounded birthdays
(\code{found:prop:proper}; \code{birthdays\_bounded}, \code{small\_bounded},
\code{not\_small} in \path{SignSequence.lean}), small-cut filling
(\code{small\_cut\_fillers}, \code{exists\_separator\_of\_small\_sets}) and
the inconsistency of unrestricted cuts (\code{noUnrestrictedCuts}). The
declarations cited by sources~11 and~12 exist at the current revision; none
states a result of this report.

% ---- hand chunk choices@1
\subsection{Where the merge had to choose}\label{swo:sub:choices}

\merge
\begin{itemize}
\item \emph{Base and order.} Source~11; its section order is kept, and
source~12's material follows the base result it duplicates or extends.
\item \emph{Strongest statement printed as the statement.} The choice
equivalence is source~11's (over $\GB$; source~12 assumes set choice). The
diagonal theorem is printed twice: source~11's (\cref{swo:gl:diagonal}) and
source~12's stronger \cref{swo:st:thm:diagonal}. The birthday-cutoff
statement is source~12's dichotomy.
\item \emph{Different routes kept in full:} source~12's set-sized spectra
through binary coding length, its avoidance lemma, its proof of the choice
equivalence over $\GB+\mathrm{AC}$, its largest common labelled prefix, its
residual adjacency criterion, its binary-class concatenation, its prefix
extension, and its birthday-bounded universality.
\item \emph{Corrections.} The renamings of \cref{swo:sec:notation}; one
mistyped reference in source~11; source~12's weak-choice question printed as
answered (\cref{swo:st:n14.5}).
\end{itemize}

% ---- hand chunk answers@1
\subsection{Where one source answers another}\label{swo:sec:answers}

\merge\ Dated notes, 2 October 2026.
\begin{enumerate}[label=(\alph*)]
\item \emph{Source~12's question on the weak-choice strength}
(\cref{swo:st:n14.5}, and its \cref{swo:st:n4.2}): answered by source~11.
Over $\GB$ without set choice, a class well-order of $\No$ exists exactly when
global choice holds (\cref{swo:cf:choice}); see \cref{swo:rem:answered-choice}.
\item \emph{The reals report's question ``Other ground orders''}
(\code{lwo:sec:research}, which names ``set-sized suborders of surreal
numbers''): answered in part by \cref{swo:st:thm:transition} for
$X=\No_{<\kappa}$ and the minimal stratum: $\WO_\mu(\No_{<\kappa})\equi\B_\mu$,
except at singular strong-limit $\kappa$, where it is
$\equi\B_{\kappa\cdot\kappa}$ and not $\equi\B_\kappa$. Open for longer
strata and other alphabets (\cref{swo:st:n14.1}).
\item \emph{Hamkins's question whether universality of the surreal line is a
weak global choice principle} \cite{HamkinsUnivChoice}: not settled here, as
source~12 says after its \cref{swo:st:n14.5}.
\end{enumerate}


% ---- source 11, lines 127-189
\section{The objects being studied}\label{swo:sec:objects}

There are two orders in this problem. Each object under consideration is
itself a well-order of some collection of surreal numbers. We then compare
these objects lexicographically using the \emph{numerical} order on their
surreal labels. This second order need not be a well-order. Indeed, the
orders below contain abundant descending sequences.

Write \(L\emb M\) for an order embedding and \(L\equi M\) for embeddings
in both directions. Equimorphism does not imply isomorphism.

\begin{center}
\small
\begin{tabular}{@{}>{\raggedright\arraybackslash}p{0.18\textwidth}>{\raggedright\arraybackslash}p{0.41\textwidth}>{\raggedright\arraybackslash}p{0.33\textwidth}@{}}
\toprule
Object & Elements & Logical size and interpretation\\
\midrule
\(\WO(X)\) & Ordinal-indexed bijections onto a fixed set \(X\subseteq\No\)
& A set; ordinary \(\ZFC\) order theory\\[4pt]
\(\Words(\No)\) & Injective set-length ordinal words in \(\No\)
& A class of sets; well-orders of arbitrary set-sized subsets\\[4pt]
\(\Hsl(\No)\) & Set-like class well-orders of all \(\No\), equivalently
class bijections \(\Ord\to\No\)
& A higher-order totality; use class variables or uniform families\\[4pt]
\(\Hall(\No)\) & All class well-orders of all \(\No\), including
non-set-like ones
& A higher-order totality; direct labeled comparison is needed\\
\bottomrule
\end{tabular}
\end{center}

The fraktur notation is deliberate. It does not declare a class with
proper classes as members. Statements about \(\Hsl\) and \(\Hall\)
are formulas with class variables, or statements relative to a specified
larger universe. An embedding of an ordinary class order into one of these
totalities is a uniform family: one class relation encodes the object
assigned to each set-valued index.

\paragraph{Principal answers.}
For a fixed infinite alphabet \(X\), the exact ordinal ceiling is
\((\card X)^+\). Passing to all set-length surreal words removes that
ceiling: there is an explicit embedding of the entire class into
\(\No\), and the singleton words embed \(\No\) in the other direction.
The two orders are nevertheless not isomorphic, because limit-length
words have set-sized left cofinality.

Under global choice, the set-like global fragment contains a uniformly
indexed copy of every class linear order. This includes every set ordinal,
\(\Ord\), reverse \(\Ord\), \(\Ord+\Ord\), and the whole numerical
surreal order. Yet the entire fragment cannot be exhausted by one uniform
class of set-codes. Class universality of a subfamily is compatible with
a larger totality of class objects.

The unrestricted global order is different again. It has adjacent pairs
and convex residual fibers. Every raw well-order begins with a canonical
\(\Ord\)-sequence; if that sequence omits labels, those labels form a
later well-ordered tail. This decomposition permits a precise study of
\emph{all} well-orders without treating them all as set-like.
It also supplies a nonembedding obstruction: the unrestricted order has
too many disjoint two-point fibers to embed into the set-like order in
the full containing-universe model. A related uniform definability
obstruction is provable in \(\KM+\GC\).


% ---- source 11, lines 190-249
\section{Relationship to ProveIt and the literature}\label{swo:sec:sources}

We inspected the repository at the immutable revision
\[
 \texttt{ae7c1e6aae3dc58c9c638b35646865ad6c10387d}.
\]
The closest precursor is the report on lexicographic well-orderings of
the reals \cite{ProveItReals}, in the subdirectory
\path{lexicographic-well-orderings-of-reals/}
of the repository's ordinal-order reports.
Its general alphabet theorems already establish the cardinality,
universality, and ordinal cutoff for every infinite linearly ordered
set \(X\). Their application to birthday fragments is a transfer of
existing results, not a new generic alphabet theorem.

That report also proves, with \(\mathfrak c=2^{\aleph_0}\),
\[
 \No_{<\theta}\emb\WO(\mathbb R)\quad(\theta<\mathfrak c^+),
 \qquad
 \No_{<\mathfrak c^+}\not\emb\WO(\mathbb R).
\]
The negative assertion follows from the all-plus chain of length
\(\mathfrak c^+\). It warns against taking an incoherent union of
stagewise embeddings.

The repository's foundational report \cite{ProveItFoundations} supplies
the relevant size discipline: individual surreals are set-coded, bounded
fragments are sets, arbitrary class cuts need not be fillable, and
quantifying over class functions does not create a class of those
functions. Its local set-valued recursion arguments are directly useful.
A nearby vector-space report uses the same distinction to avoid an
unnecessary \(\ETR\) hypothesis \cite{ProveItVector}.

The sign-sequence presentation is standard \cite{Gonshor,BagayokoHoeven}.
Absolute universality and homogeneity of the surreal order are part of
the established theory discussed by Ehrlich \cite{Ehrlich2012}.
We supply the short recursion arguments needed here.

\paragraph{The Kanovei--Shelah connection.}
There is a methodological connection, but the objects are different.
Kanovei and Shelah's actual index consists of maps
\[
 a:\mathfrak c\longrightarrow\mathcal P(\mathbb N)
\]
whose ranges are ultrafilters. They compare these maps lexicographically
using the binary order on \(\mathcal P(\mathbb N)\)
\cite[Section 1]{KanoveiShelah}. The maps need not be injective and do
not enumerate the entire real line. Their construction illustrates how
a definable linear order on many codes can replace a definably selected
well-order. The present problem shares that idea while introducing a
different level of class-size complexity.

\paragraph{Contribution and attribution.}
The central developments here are explicit surreal-word compression,
direct comparison of unrestricted class well-orders, their canonical
initial-\(\Ord\) decomposition, the bounded-support interpolation
theorem, and the relation between omitted ranges and non-set-like tails.
No claim about real-specific countability, separability, or gap spectra
is transferred to surreal alphabets without proof.


% ---- hand chunk intros_head@1
\section{The other source's introduction}\label{swo:sec:intros}

\merge\ Source~12 opens with its own statement of the problem, of its
relation to ProveIt, and of the Kanovei--Shelah connection. Its introduction
is printed whole, followed by notes where the merge qualifies it. Every
result it announces is printed, once, in the body.

% ---- source 12, lines 149-231
\subsection{Source 12: scope, predecessors, and the results at a glance}
\label{swo:st:sec:scope}

The phrase ``well-orders of surreal numbers'' hides several distinct
objects. The usual numerical order of $\No$ is not a well-order. One may
instead put an unrelated well-order on the entire surreal carrier, or
consider numerical well-orders on subclasses, or work with an ordinary
set of bounded-birthday surreals. This article treats all three viewpoints,
but the principal object is the lexicographic comparison of
\emph{exhaustive well-orders of a fixed carrier}.

The supplied ProveIt repository contains a substantial predecessor,
\emph{Lexicographic Orders of the Well-Orderings of the Reals}, merging
three reports and distinguishing the full enumeration order from its
minimal-length slice \cite{ProveItReals}. The inspected README explicitly
marks that report as unrefereed and not formalized. Its real-number
spectra, binary coding, finite-tail geometry, and warnings about
isomorphism versus equimorphism motivate the present work. The repository
also contains a Lean development of simplest elements of convex sets of
sign sequences \cite{ProveItSimplicity}. These are useful foundations, not a
formal verification of the new theorems below.

We pinned the audit to repository commit
\begin{center}
\small\texttt{63bb8b1d99359e7188abf0a33fb152e121858150}.
\end{center}
The audit was targeted, not an exhaustive review of the repository.
All essential arguments used here are written out, so no unreviewed
research claim from the repository is required as a black box.

\subsubsection{Three different orders}

For an infinite set-sized linear order $X$ of cardinality $\mu$, define
\[
 \WO_\alpha(X)=\{e:\alpha\longrightarrow X:e\text{ is bijective}\},
 \qquad
 \WO(X)=\bigcup_{\mu\leq\alpha<\mu^+}\WO_\alpha(X).
\]
They are ordered by their first unequal entry, using the given order of
$X$, not the well-order being represented.

For the proper class $\No$, let $\Hsl(\No)$ denote, as a \emph{predicate on
classes}, the class bijections $E:\Ord\to\No$. These represent exactly the
set-like global well-orders. Let $\Hall(\No)$ denote the predicate
of being an arbitrary class well-order relation on $\No$.
Neither displayed notation asserts the existence of a class having these
class-sized objects as members.

The comparison on $\Hsl(\No)$ uses a least ordinal of disagreement.
The comparison on $\Hall(\No)$ uses a largest common
\emph{labelled} initial segment. The latter can be a proper class, and
this is essential rather than a cosmetic extension.

\subsubsection{What is proved, and what is not being claimed}

The principal developments are the elementary common-prefix theorem
(\cref{swo:st:thm:fulllex}), the global-choice equivalence
(\cref{swo:st:thm:choice}), the non-enumerability theorem
(\cref{swo:st:thm:diagonal}), and the singular transition
(\cref{swo:st:thm:transition}). The full and minimal bounded spectra are proved
in \cref{swo:st:sec:bounded,swo:st:sec:transition}. The difference between the dense
set-like order and finite-tail geometry in the full order is proved in
\cref{swo:st:sec:setlike,swo:st:sec:fullgeometry}.

These are mathematical claims with supplied proofs. We do not claim
that they settle a previously catalogued open problem, that their exact
formulations are historically new, or that they have been independently
refereed. The proposed research questions at the end are precisely that:
questions suggested by the results, not an assertion of a complete
literature census.

\subsubsection{The Kanovei--Shelah connection}

Kanovei and Shelah order maps
$a:\mathfrak c\to\mathcal P(\mathbb N)$ whose ranges are ultrafilters,
using a fixed lexicographic order on $\mathcal P(\mathbb N)$
\cite[\S1]{KanoveiShelah}. Their maps need not be bijections, and their index is not
the literal collection of all well-orders of $\mathbb R$, still less all
class well-orders of $\No$. The useful common idea is to compare
representations at their first difference instead of choosing one
canonical representative. Their index construction remains set-sized.
Moving to all global surreal well-orders creates an additional size level
that their construction does not remove.

% ---- hand chunk intro12_note

\merge\ \emph{On source~12's introduction.} Source~12 read the reals
report's README and the Lean file \path{SignSequenceSimplicity.lean}, not the
reals report's article; it therefore re-proved the set-sized patterns it
needed. They are printed here as a marked route with pointers to the reals
report (\cref{swo:sub:lwo-pointer,swo:st:sec:bounded}). The reals report's
article has not changed since source~12's pin. Source~12's principal
developments are printed as \cref{swo:st:thm:fulllex} (a second route to
\cref{swo:cf:lex}), \cref{swo:st:thm:choice} (a second route to
\cref{swo:cf:choice}, over $\GB+\mathrm{AC}$), \cref{swo:st:thm:diagonal}
(printed beside \cref{swo:gl:diagonal}) and \cref{swo:st:thm:transition}
(new).

% ---- source 11, lines 250-308
\section{Surreal order facts and foundational conventions}\label{swo:sec:prelim}

\subsection{Signs, birthdays, and set cuts}

A surreal number is represented by a sign sequence
\[
 s:\alpha\longrightarrow\{-,+\}
\]
of set-ordinal length \(\alpha\), its birthday \(\bd(s)\). Numerical
comparison at the first disagreement uses
\[
 -\ <\ \text{termination}\ <\ +.
\]
This differs from dictionary order with every proper prefix smaller
than every extension. A sign sequence can be on either side of an
extension, depending on its next sign.

For an ordinal \(\theta\), let
\[
 \No_{<\theta}=\bigcup_{\alpha<\theta}\{-,+\}^{\alpha}.
\]
It is a set. The full \(\No\) is proper-class-sized, because all-plus
sequences embed all ordinals. Every set of surreals has bounded birthdays
by Replacement.

\begin{definition}\label{swo:n3.1}
A class linear order \(L\) has the \emph{set-cut property}, called
\(\eta_{\Ord}\) here, if for all sets \(A,B\subseteq L\) with \(A<B\)
there is \(x\) with \(A<x<B\). Empty sides are allowed. For a cardinal
\(\kappa\), the set-sized notation \(\eta_\kappa\) requires the same
property when \(\card{A\cup B}<\kappa\).
\end{definition}

The surreal cut construction supplies this property for \(\No\)
\cite[Section 2]{BagayokoHoeven}. A separator can be chosen canonically
as the unique simplest separator. This statement concerns \emph{sets}
of options; applying it to the whole class against the empty right side
would be inconsistent.

\begin{lemma}[Deleting a set]\label{swo:lem:delete}
If \(S\subseteq\No\) is a set, then \(\No\setminus S\) has the
set-cut property. Every nonempty open interval of \(\No\) is a
proper class.
\end{lemma}
\begin{proof}
Let \(A<B\) be sets in \(\No\setminus S\). First choose \(x\) with
\(A<x<B\), and put
\[
 S_-=\{s\in S:s\leq x\},\qquad S_+=\{s\in S:x<s\}.
\]
Then \(A\cup S_-<B\cup S_+\). A separator of this new set cut
avoids both parts of \(S\). The argument allows empty sides.
If an interval \((a,b)\) were a set \(S\), applying the result to
\(\{a\}<\{b\}\) in \(\No\setminus S\) would be impossible.
\end{proof}

Choosing both separators by simplicity makes the deletion argument
uniform in the set parameters without global choice.


% ---- hand chunk avoid12_head

\merge\ \emph{Third route to the first clause of \cref{swo:lem:delete}
(source~12, Lemma~6.2).} Source~12 avoids a forbidden set by using that open
intervals of $\No$ are proper classes.

% ---- source 12, lines 580-589
\begin{lemma}[Avoiding a forbidden set]\label[lemma]{swo:st:lem:avoid}
If $A<B$ are sets of surreals and $D$ is a set, then some separator
$z$ of $A,B$ lies outside $D$.
\end{lemma}
\begin{proof}
Choose $u$ between $A$ and $B$, and then $v$ between $\{u\}$ and $B$.
Every point of $(u,v)$ separates $A,B$, and that interval is a proper
class. It cannot be contained in $D$. With a fixed global well-order,
choose the first acceptable point, making the operation uniform.
\end{proof}

% ---- source 11, lines 309-329
\subsection{Which recursion principle is used?}

Here \(\GB\) has a \(\ZF\) set part, elementary class comprehension,
and class replacement; choice is stated separately.
We use \(\GBC=\GB+\GC\), which includes set choice.
Cardinality calculations and cardinal-valued cofinality statements
throughout the article are read in \(\ZFC\), or in the set part of
\(\GBC\). The explicitly definable sign code and the direct
set-cut constructions retain their stated choice-free strength.

Ordinary recursion along \(\Ord\), with a \emph{set} output at every
stage and a rule using the set of earlier values and fixed class
parameters, is available in \(\GB\). On each set ordinal perform the
usual recursion with those parameters. Uniqueness makes the local
solutions agree; elementary comprehension assembles their graph.

This is different from arbitrary elementary transfinite recursion,
\(\ETR\), whose stages may themselves be classes. The latter supports
iterated truth predicates and is stronger than \(\GBC\)
\cite{GitmanHamkins}. The prefix searches, partial permutations, and
back-and-forth maps used here have set-sized stages.

% ---- hand chunk found12_head

\merge\ The other sources fix their foundational conventions in their own
sections, printed whole next. They agree with source~11's conventions, with
one difference of base theory: when it analyses choice, source~12 works over
$\GB+\mathrm{AC}$ (\cref{swo:sec:notation}).

% ---- source 12, lines 233-315
\subsection{Source 12: the surreal carrier and the foundational conventions}
\label{swo:st:sec:foundations}

\subsubsection{Sign sequences}

We use the sign-sequence presentation
\[
 \No=\bigcup_{\alpha\in\Ord}\{-,+\}^{\alpha}.
\]
Each individual sequence is a set. Its length $\bd(s)=\lh(s)$ is its
birthday. At the first difference, signs and termination are compared by
\[
 -\ <\ \End\ <\ +.
\]
This is the usual numerical order, not the ordinary ``shorter comes
first'' dictionary convention. Simplicity is the separate relation of
being a proper sign prefix. The standard surreal construction and its
simplicity structure are treated by Gonshor and Ehrlich
\cite{Gonshor,Ehrlich2012}.

We shall use the following standard feature, also obtainable directly
from the construction: whenever $A,B$ are \emph{sets} of surreals with
$A<B$, there is a surreal $z$ with $A<z<B$. There is a unique such
separator of minimum birthday. Moreover, if every member of $A\cup B$
has birthday less than an ordinal $\gamma$, a separator can be constructed
with birthday at most $\gamma$. This is interpolation, not Dedekind
completeness: no least upper bound is being asserted.

Every set of surreals has a birthday bound, by Replacement. Every
nonempty numerical interval contains a proper class of surreals. To see
the latter without arithmetic, start with $a<b$ and recursively choose
$x_\alpha$ between $\{a\}\cup\{x_\xi:\xi<\alpha\}$ and $\{b\}$.
The previous values form a set at each stage. Thus one obtains distinct
points for all ordinals $\alpha$.

\subsubsection{Well-orders and set-likeness}

A class well-order on a class $X$ is a strict linear class relation $R$
for which every nonempty subclass has an $R$-least element. In GBC this
agrees with the usual formulations using nonempty set subsets or absence
of descending $\omega$-sequences; the relationship is discussed in
\cite[\S3]{HamkinsWoodin}. We use strict relations throughout.

It is \emph{set-like} if
\[
 \pred_R(x)=\{y\in X:yRx\}
\]
is a set for every $x$. A proper-class set-like well-order is isomorphic
to $\Ord$: map $x$ to the ordinal type of its set of predecessors. The
range is an initial segment of $\Ord$, and cannot be a set ordinal when
$X$ is proper. This construction uses only set-sized order types.

Not all class well-orders are set-like. The ordered sums
$\Ord+1$, $\Ord+\Ord$, and $\Ord+\Ord+1$ are examples. Their notation
means ordered sums of class orders, not new von Neumann ordinals.

\subsubsection{Foundational levels}

Our set-sized results are theorems of ZFC. Our global results are in GBC,
meaning Gödel--Bernays set theory with global choice. When analyzing the
necessity of choice we work over GB plus ordinary, set-level AC.
The full collection of global well-orders is handled by formulas with
class variables, not by pretending that a proper class is an element of
another GBC class.

Kelley--Morse theory allows stronger comprehension, but its classes still
have sets, not proper classes, as members. Moving from GBC to KM does not
by itself create a powerclass object of all subclasses of $\No$.
One can instead use a larger set-theoretic universe, a set model with a
specified collection of classes, or an explicitly higher-order theory.
These are distinct interpretations and will be kept distinct.

\subsubsection{Why an ordinal-length recursion here need not invoke ETR}

A recursion on $\Ord$ that produces one \emph{set} at each stage and uses
the set of previous values is available by the ordinary recursion theorem
with class parameters in GB. In contrast, ETR permits a genuinely
class-valued recursive stage, represented by a class of pairs of a stage
and an arbitrary set. Even a truth hierarchy brings stronger principles
into play \cite[\S2]{HamkinsWoodin}. Our forth constructions, block lengths, and
canonical enumeration of a class with set-sized initial segments are
set-valued recursions. We do not invoke a universal satisfaction predicate
or unrestricted class-valued recursion.

% ---- source 11, lines 331-358
\subsection{Universality and uniqueness of the saturated class order}

\begin{proposition}\label{swo:prop:universal}
In \(\ZFC\), every set linear order embeds into \(\No\). In
\(\GBC\), every class linear order embeds uniformly into \(\No\).
Any two proper-class \(\eta_{\Ord}\) orders are isomorphic in
\(\GBC\).
\end{proposition}
\begin{proof}
Enumerate the source in a cardinal-length well-order, or in type
\(\Ord\) for a proper class, using global choice. At stage \(\alpha\),
the earlier images form a set. Separate the images of earlier points
below the new point from those above it, and choose a separator.
This extends an order embedding. Each partial map is a set.

For uniqueness, globally enumerate both class orders and alternately
insert the first unused point from each enumeration. The appropriate
set cut in the other order is fillable. Global choice selects a
separator when none is distinguished. The union of the partial
isomorphisms is a surjective class order isomorphism.
\end{proof}

The same construction, started from a specified set-sized partial
isomorphism, gives set homogeneity. These are the order-theoretic
universality and homogeneity features of \(\No\)
\cite{Ehrlich2012}. The assertion that every \emph{class} order embeds
is made with global choice; it is not used as an unproved converse
to global choice.

% ---- hand chunk univ12_head

\merge\ \emph{Source~12's version (its Lemmas~6.5 and~6.6).} It exposes the
size of the forth recursion, with a birthday bound for set orders, and proves
the class isomorphism $\No\setminus A\cong\No$ that source~11 uses in
\cref{swo:gl:cylinders}.

% ---- source 12, lines 655-687
\begin{lemma}[Surreal universality, with the size of the recursion exposed]
\label[lemma]{swo:st:lem:universalNo}
In GBC, every class linear order admits a class order-embedding into
$\No$. For a set order enumerated in type $\lambda$, one can arrange
that the image of its $\xi$th enumerated element has birthday at most
$\xi$, for each $\xi<\lambda$.
\end{lemma}
\begin{proof}
Use a global-choice enumeration of the carrier, either by an ordinal or
by $\Ord$. When the next element is considered, its previously embedded
predecessors and successors form two sets $A<B$. Send it to the simplest
surreal separator. It is distinct from all previous images and preserves
order. Every stage has set-sized history. Induction gives the birthday
bound, because the earlier birthdays are less than the current stage.
This is the standard forth argument behind surreal universality
\cite{Ehrlich2012,HamkinsUniversal}.
\end{proof}

\begin{lemma}[Deleting a set does not change the numerical class order]
\label[lemma]{swo:st:lem:delete}
In GBC, $\No\setminus A$ is order-isomorphic to $\No$ for every set $A$.
Every nonempty open numerical interval is also order-isomorphic to $\No$.
\end{lemma}
\begin{proof}
Each target has interpolation for arbitrary separated sets, by
\cref{swo:st:lem:avoid}, adding interval bounds when needed. Each is a proper
class and has an $\Ord$-enumeration by global choice. Carry out a
back-and-forth between the two enumerations. At each ordinal stage,
extend the partial isomorphism to the next unused domain point and then
to the next unused target point. The previous partial map is a set, and
the required cut can be filled while avoiding its set-sized range.
At the end the class map is onto and order-preserving.
\end{proof}

% ---- hand chunk part1
\part{Sets of surreals}\label{swo:part:sets}

% ---- source 11, lines 362-531
\section{Well-order codes on a fixed set alphabet}\label{swo:sw:fixed}

Let \(X\) be an infinite linearly ordered set, and put
\(\lambda=\card{X}\).  We work in \(\ZFC\) throughout this section.
A well-order on the whole carrier \(X\) has a unique increasing
enumeration \(f:\alpha\longrightarrow X\), where increasing refers to
that well-order, not to the native order of \(X\).  Its domain satisfies
\(\card{\alpha}=\lambda\), and hence \(\alpha<\lambda^+\).
Write \(\WO(X)\) for these bijective enumeration codes, and introduce
two related sets:
\[
 \mathcal I(X)=
 \bigcup_{\alpha<\lambda^+}\{f:\alpha\hookrightarrow X\},
 \qquad
 \WO_\lambda(X)=\{f:\lambda\longrightarrow X:
                                      f\text{ is bijective}\}.
\]
Thus \(\mathcal I(X)\) codes well-orders of arbitrary subsets of \(X\),
including the empty subset, whereas \(\WO(X)\) always uses the whole
carrier.  All three sets inherit the lexicographic order defined
earlier.  A proper prefix is smaller than its extension in
\(\mathcal I(X)\).  Two distinct members of \(\WO(X)\) cannot be proper
prefixes of one another: the shorter enumeration would already have
used every element of \(X\).

The repository's real-well-ordering report already proves the
exhaustive and cardinal-length statements for arbitrary infinite
alphabets \cite{ProveItReals}.  We record those results together
with the injective-word variant, and include a common proof
to keep the size distinctions explicit.  Neither density nor
completeness of the native order is needed.

\begin{theorem}[Cardinality and exact ordinal spectrum]\label{swo:sw:setspectrum}
Each of \(\WO_\lambda(X)\), \(\WO(X)\), and \(\mathcal I(X)\)
has cardinality \(2^\lambda\), contains a copy of every linear order
of cardinality at most \(\lambda\), and contains neither
\(\lambda^+\) nor \((\lambda^+)^*\).
Consequently the ordinals embeddable in any of them are exactly
the ordinals \(\alpha<\lambda^+\).
\end{theorem}

\begin{proof}
There are at most \(\lambda^+\) permitted domains.  For each domain
there are at most \(\lambda^\lambda=2^\lambda\) functions, so
\(\card{\mathcal I(X)}\leq\lambda^+\cdot2^\lambda=2^\lambda\).
For the reverse inequality, partition \(X\) into pairs
\(\{a_\xi,b_\xi\}\), \(\xi<\lambda\), with \(a_\xi<b_\xi\).
The lexicographic order on \(\lambda\times\{0,1\}\), taking
\(\xi\) first, has order type \(2\cdot\lambda=\lambda\).
For \(\varepsilon\in2^\lambda\), list the \(\xi\)-th pair as
\((a_\xi,b_\xi)\) if \(\varepsilon(\xi)=0\), and as
\((b_\xi,a_\xi)\) otherwise.  The resulting map
\[
                 2^\lambda\longrightarrow\WO_\lambda(X)
\]
is an order embedding when \(2^\lambda\) has lexicographic order.
This proves all three cardinality assertions.

Let \(L\) be any linear order of cardinality at most \(\lambda\),
and fix an injection \(\iota:L\hookrightarrow\lambda\).
Associate to \(t\in L\) the binary function \(c_t\) whose value at
\(\iota(s)\) is \(1\) precisely when \(s<t\), and whose remaining
coordinates are \(0\).  If \(t<u\), then \(c_t\leq c_u\)
coordinatewise, and the inequality is strict at \(\iota(t)\).
Their first difference therefore has the correct direction.
The preceding pair construction embeds \(L\) into
\(\WO_\lambda(X)\).

It remains to exclude longer monotone sequences.  Suppose that
\((f_\xi)_{\xi<\lambda^+}\) is strictly increasing or strictly
decreasing in \(\mathcal I(X)\).
We recursively construct eventual common prefixes \(p_\alpha\)
of length \(\alpha\), for \(\alpha<\lambda^+\).
At a limit stage, take the union of the earlier prefixes.
Regularity of \(\lambda^+\) ensures that a tail of the original
sequence extends that union.  At a successor step, at most one
member of this tail equals \(p_\alpha\); discard it and all earlier
indices.  The values at coordinate \(\alpha\) are now monotone.
A monotone sequence indexed by a tail of \(\lambda^+\), taking
at most \(\lambda\) values, is eventually constant: the first
occurrence indices of all its values have a supremum below
\(\lambda^+\).  Append this eventual value to \(p_\alpha\).
Every constructed prefix is injective because it occurs in
members of the original sequence.  The union of the prefixes
through \(\lambda^+\) would therefore be an injection
\(\lambda^+\hookrightarrow X\), a contradiction.
Finally, an ordinal at least \(\lambda^+\) contains
\(\lambda^+\), whereas every smaller ordinal has cardinality
at most \(\lambda\).
\end{proof}

The distinction between the models also affects their elementary
order topology.  The order \(\WO_\lambda(X)\) is dense.
Indeed, let \(f<_{\lex}g\), with first difference at \(\xi\).
If the tail of \(f\) after \(\xi\) has an ascent, interchanging
the two entries produces a point strictly between \(f\) and \(g\).
If the tail of \(g\) has a descent, the corresponding interchange
works in the other direction.  If neither exists, the first
remaining tail is strictly decreasing and the second strictly
increasing.  Their common set of entries, after removing the two
entries at \(\xi\), is infinite and would be both well-ordered
and reverse-well-ordered.  This is impossible.

By contrast, \(\WO(X)\) has adjacent pairs.  Choose \(a<b\) in \(X\),
enumerate \(X\setminus\{a,b\}\) as a prefix \(p\), and compare
\(p\cat\langle a,b\rangle\) with
\(p\cat\langle b,a\rangle\).  Any intervening enumeration would have
to share \(p\), after which only these two orders remain.
Thus a cardinal-length convention changes the resulting linear
order even though it does not change its cardinality or ordinal
embedding spectrum.

\begin{proposition}[Exactly where jumps occur]\label{swo:sw:adjacency}
Suppose \(f<_{\lex}g\) in \(\WO(X)\), with common prefix \(p\)
up to their first difference.  Put
\(Y=X\setminus\ran(p)\), \(a=f(\dom(p))\), and
\(b=g(\dom(p))\), so \(a<b\).
The pair is adjacent precisely when \(Y\) is finite,
\(a,b\) are consecutive in its native order, the tail of
\(f\) after \(a\) lists \(Y\setminus\{a\}\) decreasingly,
and the tail of \(g\) after \(b\) lists
\(Y\setminus\{b\}\) increasingly.
\end{proposition}

\begin{proof}
If an unused element lies strictly between \(a\) and \(b\),
start an enumeration with that element after \(p\) and
complete it; this lies between \(f\) and \(g\).
Likewise, an ascent after \(a\), or a descent after \(b\),
can be interchanged to create an intervening enumeration.
Adjacency thus forces the two remaining tails to be
decreasing and increasing, respectively.
Their common native suborder \(Y\setminus\{a,b\}\) is then
both well-ordered and reverse-well-ordered.
It must be finite: an infinite well-order contains its
first \(\omega\) elements, a subset with no greatest element.
This proves necessity.
Conversely, an intervening enumeration must share \(p\).
There is no possible next entry strictly between \(a\)
and \(b\).  Among enumerations beginning with \(a\), the
displayed decreasing continuation is the largest; among
those beginning with \(b\), the increasing continuation
is the smallest.  Neither branch admits an intervening point.
\end{proof}

This description also identifies the finite components generated
by adjacency.  Write the domain of an enumeration uniquely as
\(\theta=\delta+n\), with \(\delta\) a limit ordinal and
\(n<\omega\).  Its component consists of the \(n!\) permutations
of the final \(n\) entries, keeping the prefix of length
\(\delta\) fixed.  In particular, words of limit length have
no immediate neighbor.  A jump cannot cross from one such
component to another, because its common-prefix remainder
would be infinite.

Endpoints can be classified just as concretely.
The least element of \(\WO(X)\), if one exists, must always
choose the least unused native element: otherwise interchanging
that entry with a later smaller one decreases the enumeration.
Hence a least element exists exactly when the native order
on \(X\) is a well-order, and is its increasing enumeration.
The greatest-element statement is the reverse assertion.
For \(\WO_\lambda(X)\), the same argument shows that
a least element exists exactly when the native order type is
the initial ordinal \(\lambda\), not merely an arbitrary
ordinal of cardinality \(\lambda\).
For example, an alphabet of native type \(\lambda+1\)
has a least full enumeration, but no least
\(\lambda\)-length permutation.


% ---- hand chunk lwo_pointer@1

\subsection{Pointers to the reals report}\label{swo:sub:lwo-pointer}

\merge\ The theorems of this section hold for every infinite linearly ordered
set and are theorems of the repository's reals report, which proves them in
that generality. Source~11 says so (\cref{swo:sec:sources}); its repository
audit (\path{11-raw-orders-repository_audit.md}) adds that their application
to birthday fragments ``should be attributed as applications of existing
repository results, not marketed as new general theorems''. Source~12
re-proves them; they are printed with these pointers, and no novelty is
claimed for any of them.

\begingroup\small
\begin{longtable}{@{}>{\raggedright\arraybackslash}p{0.30\textwidth}>{\raggedright\arraybackslash}p{0.30\textwidth}>{\raggedright\arraybackslash}p{0.32\textwidth}@{}}
\toprule
Statement & Here & Reals report\\
\midrule\endhead
Prefix-freeness & before \cref{swo:sw:setspectrum}; \cref{swo:st:lem:prefixfree} &
\code{lwo:lem:prefix-free}\\[2pt]
Cardinality; pair swaps & \cref{swo:sw:setspectrum};
\cref{swo:st:lem:pairs,swo:st:thm:setspectrum} &
\code{lwo:prop:cardinality}, \code{lwo:lem:all-small-cubes}\\[2pt]
Cut coding & \cref{swo:sw:setspectrum}; \cref{swo:st:lem:cutcode} &
\code{lwo:lem:cut-code}, \code{lwo:thm:ordinal-spectrum}(a)\\[2pt]
No $\lambda^+$-sequence & \cref{swo:sw:setspectrum}; \cref{swo:st:lem:nochain} &
\code{lwo:thm:no-long}, \code{lwo:thm:ordinal-spectrum}\\[2pt]
Strict cubes & \cref{swo:st:lem:strictcube} & \code{lwo:thm:cube-strict}\\[2pt]
Coding length $\lambda^+$ & \cref{swo:st:thm:setspectrum} &
\code{lwo:thm:full-length}\\[2pt]
Extrema; adjacency; $n!$ blocks & paragraphs after \cref{swo:sw:adjacency} &
\code{lwo:lem:extrema}, \code{lwo:thm:adjacency}, \code{lwo:thm:blocks}\\
\bottomrule
\end{longtable}
\endgroup

The arguments are the reals report's: pair swaps, cut codes, prefix
stabilization and cylinder fusion. Source~12's Section~9 is printed in full
because the singular transition of \cref{swo:st:sec:transition} uses its
lemmas.

% ---- source 12, lines 863-1016
\subsection{Source 12's route: a general set-sized spectrum theorem}
\label{swo:st:sec:bounded}

We temporarily forget the surreal structure. Let $X$ be any infinite
linearly ordered set, and write $|X|=\mu$, an initial cardinal.
The spectra below depend only on the cardinality, not on density,
separability, or field operations of $X$.

\begin{lemma}[Prefix-freeness of exhaustive enumerations]\label[lemma]{swo:st:lem:prefixfree}
Two distinct exhaustive injective ordinal enumerations of $X$ are never
proper prefixes of one another. Hence first-difference lexicography is
unambiguous on $\WO(X)$.
\end{lemma}
\begin{proof}
An exhaustive prefix has already used every element of $X$. An injective
proper extension would have to use a new element, which is impossible.
\end{proof}

For any ordinal $\gamma$, let $\BB_\gamma=(2^\gamma,<_{\lex})$.
For linear orders $L,M$, write $L\emb M$ for an order-embedding and
$L\equim M$ for embeddings in both directions. Define
\[
 \bl(L)=\min\{\gamma:L\emb\BB_\gamma\}
\]
when such a binary coding ordinal exists. This is an ordinal-valued rank,
not a cardinal invariant.

\begin{lemma}[Strict binary hierarchy]\label[lemma]{swo:st:lem:strictcube}
For ordinals $\alpha,\beta$,
\[
 \BB_\alpha\emb\BB_\beta\quad\Longleftrightarrow\quad\alpha\leq\beta.
\]
In particular $\bl(\BB_\alpha)=\alpha$.
\end{lemma}
\begin{proof}
Zero-padding proves the forward construction when $\alpha\leq\beta$.
It suffices to rule out an embedding
$f:\BB_{\beta+1}\to\BB_\beta$.

Recursively construct an input prefix $s_\xi$ of length $\xi$ and an
output prefix $t_\xi$ of length $\xi$, so that every input extending
$s_\xi$ is sent to an output extending $t_\xi$. At a successor stage,
split the input cylinder into its two next-bit halves and split the
output cylinder likewise. The two input halves have strictly ordered
images. They cannot both have images meeting both output halves:
an image point in the upper output half from the lower input half would
exceed an image point in the lower output half from the upper input half.
Thus one input half has image contained in a single output half.
Choose those halves to continue.

At a limit take unions; the corresponding cylinders are nonempty and
the containment persists. At stage $\beta$, two inputs remain, differing
in their last bit, but their images have the same full $\beta$-long
output. This contradicts injectivity. For $\alpha>\beta$, padding embeds
$\BB_{\beta+1}$ into $\BB_\alpha$, giving the general impossibility.
\end{proof}

\begin{lemma}[Characteristic-cut coding]\label[lemma]{swo:st:lem:cutcode}
Every linear order of cardinality at most $\mu$ embeds in $\BB_\mu$.
\end{lemma}
\begin{proof}
Enumerate its carrier by distinct coordinates $d_\xi$, padding unused
coordinates if necessary. Send $x$ to the bits
$\mathbf1[d_\xi<x]$. For $x<y$, every differing coordinate changes from
$0$ to $1$, and the coordinate representing $x$ does change. Thus the
least difference gives the correct order.
\end{proof}

\begin{lemma}[Pair orientation coding]\label[lemma]{swo:st:lem:pairs}
$\BB_\mu\emb\WO_\mu(X)$. More generally,
$\BB_\gamma\emb\WO(X)$ for every $\gamma<\mu^+$.
\end{lemma}
\begin{proof}
Partition $X$ into $\mu$ unordered pairs, sort each pair by $X$'s order,
and orient the pair according to its bit. The ordinal sum of $\mu$
two-element blocks has type $2\cdot\mu=\mu$.

For a general $\gamma<\mu^+$, choose $|\gamma|$ disjoint pairs and retain
a disjoint reserve of size $\mu$. Such a choice is possible even when
$|\gamma|=\mu$. List the pairs in type $\gamma$, then well-order all
remaining labels in type $\mu$. The resulting enumeration has length
$2\cdot\gamma+\mu<\mu^+$. Its comparison is decided by the first
oriented pair at which the inputs differ.
\end{proof}

\begin{lemma}[No $\mu^+$-long monotone chain]\label[lemma]{swo:st:lem:nochain}
Neither $\mu^+$ nor its reverse embeds into $\WO(X)$.
\end{lemma}
\begin{proof}
Suppose $(e_i)_{i<\mu^+}$ is strictly increasing. We recursively stabilize
its entries on final segments of the index set. Assume that a final
segment agrees on a prefix $p$ of length $\delta<\mu^+$. No enumeration
in a further infinite final segment can end there, by
\cref{swo:st:lem:prefixfree}. Their $\delta$th entries form a nondecreasing
$\mu^+$-sequence in a set of size $\mu$.

Such a sequence is eventually constant: choose one first occurrence for
each value in its range; their at most $\mu$ indices have supremum below
$\mu^+$ by regularity. Past that supremum no larger value can first
appear. At a limit stage $\delta<\mu^+$, take the supremum of the earlier
stabilization indices, still below $\mu^+$, and the union of the earlier
prefixes.

This would produce an injective sequence of labels from $X$ of length
$\mu^+$: every set-sized earlier part is shared by a final segment of
injective enumerations. That contradicts $|X|=\mu$. Replace
nondecreasing by nonincreasing for the reverse order.
\end{proof}

\begin{theorem}[Cardinality and exact spectra]\label{swo:st:thm:setspectrum}
For every infinite linearly ordered set $X$ of cardinality $\mu$,
\begin{align*}
 |\WO(X)|=|\WO_\mu(X)|&=2^\mu,\\
 \beta\emb\WO(X)\quad&\Longleftrightarrow\quad\beta<\mu^+,\\
 \beta\emb\WO_\mu(X)\quad&\Longleftrightarrow\quad\beta<\mu^+,\\
 \BB_\gamma\emb\WO(X)\quad&\Longleftrightarrow\quad\gamma<\mu^+,\\
 \bl(\WO(X))&=\mu^+.
\end{align*}
The ordinal statements also hold for reversed ordinals, and every
linear order of cardinality at most $\mu$ embeds even into
$\WO_\mu(X)$.
\end{theorem}
\begin{proof}
Pair coding gives $2^\mu$ elements, while all well-orders are relations
on $X$, giving the upper bound $2^{|X\times X|}=2^\mu$.
The lower ordinal bounds and universality follow from
\cref{swo:st:lem:cutcode,swo:st:lem:pairs}; \cref{swo:st:lem:nochain} gives the upper bounds.
The same characteristic coding applies to reversed ordinals.

For the cube spectrum use \cref{swo:st:lem:pairs}. If $\gamma\geq\mu^+$,
$\BB_\gamma$ contains an increasing chain of type $\mu^+$: use the
sequences that are $1$ before a moving cut and $0$ afterward. Thus it
cannot embed into $\WO(X)$.

Embed $X$ into $\BB_\mu$ using \cref{swo:st:lem:cutcode}. Encode an exhaustive
enumeration of length $\alpha<\mu^+$ by concatenating its $\mu$-long
entry codes, then pad with zeroes to length $\mu^+$. The unpadded length
$\mu\cdot\alpha$ has cardinality at most $\mu$, so is less than $\mu^+$.
Distinct enumerations disagree before either ends, and their encoded
first disagreement preserves the order. This proves
$\WO(X)\emb\BB_{\mu^+}$.

If it embedded in $\BB_\rho$ for $\rho<\mu^+$, then the copy of
$\BB_{\rho+1}$ supplied by \cref{swo:st:lem:pairs} would contradict
\cref{swo:st:lem:strictcube}. Hence the rank is exactly $\mu^+$.
\end{proof}

\begin{remark}[What the rank does not say]\label{swo:st:n9.7}
Although $\bl(\WO(X))=\mu^+$, the cube $\BB_{\mu^+}$ does \emph{not}
embed back into $\WO(X)$. An order can first admit a binary encoding at
an ordinal without containing the full cube of that ordinal.
Nor is the absence of $\mu^+$ and its reverse claimed sufficient for an
arbitrary order to embed in $\WO(X)$.
\end{remark}

% ---- source 11, lines 532-615
\section{Birthday cutoffs and their set-sized orders}\label{swo:sw:cutoffs}

Use the sign presentation of surreal numbers: \(x\in\No\) is a
function from an ordinal \(\bd(x)\) to \(\{-,+\}\).
At the first difference, the comparison is
\(-<\text{termination}<+\).
This presentation and the fundamental set-cut theorem are treated
in Gonshor and in the modern account of Bagayoko and van der Hoeven
\cite{Gonshor,BagayokoHoeven}.
For an infinite cardinal \(\kappa\), define
\[
                 \No_{<\kappa}
                   =\{x\in\No:\bd(x)<\kappa\}.
\]
This is a set.  Counting sign sequences gives
\[
       \card{\No_{<\kappa}}
          =\sum_{\alpha<\kappa}2^{\card{\alpha}}
          =2^{<\kappa}.
\]
Here \(2^{<\kappa}\) denotes the supremum of \(2^\mu\) over
cardinals \(\mu<\kappa\); it is at least \(\kappa\).

\begin{corollary}\label{swo:sw:cutoffspectrum}
Put \(\lambda=2^{<\kappa}\).
Then \(\WO(\No_{<\kappa})\) has cardinality \(2^\lambda\), is
universal for linear orders of cardinality at most \(\lambda\),
and admits exactly the ordinals below \(\lambda^+\).
The same conclusions hold for its cardinal-length permutation
suborder and for the injective-word order \(\mathcal I(\No_{<\kappa})\).
\end{corollary}

The parameter controlling this spectrum is the cardinality
\(\lambda\) of the alphabet.  It need not equal the birthday
cutoff \(\kappa\).
For \(\kappa=\omega\), the alphabet consists of the dyadic
rationals.  Its well-ordering order has cardinality
\(2^{\aleph_0}\), contains every countable linear order, and
does not contain \(\omega_1\).
For \(\kappa=\omega_1\), the alphabet consists of the surreals
with countable sign expansions and has cardinality
\(\mathfrak c=2^{\aleph_0}\).  Its well-ordering order consequently
has cardinality \(2^{\mathfrak c}\) and ordinal spectrum
\(\alpha<\mathfrak c^+\), without any use of the continuum
hypothesis.
When \(2^{<\kappa}=\kappa\), the spectrum simplifies to
\(\alpha<\kappa^+\).

\begin{lemma}\label{swo:sw:stageeta}
If \(\kappa\) is regular infinite, \(\No_{<\kappa}\) fills every cut
\(A<B\) with \(\card{A\cup B}<\kappa\).  Such a cut can still
be filled after excluding any prescribed subset of
cardinality less than \(\kappa\).
\end{lemma}

\begin{proof}
The birthdays of the members of \(A\cup B\) have supremum below
\(\kappa\), by regularity.  The simplest surreal between them
has birthday at most the supremum of their birthdays plus one,
and hence still belongs to \(\No_{<\kappa}\).
For the exclusion statement, let \(S\) be the forbidden set,
and let \(S_0\) consist of its members that already lie strictly
between \(A\) and \(B\).  Fill the cut
\(A\cup S_0<B\).  Its filler lies in the original interval but
strictly above every forbidden point in that interval.
\end{proof}

We call this property \(\eta_\kappa\), always allowing either side
of a cut to be empty.  The cardinal parameter in this notation
is explicit; conventions indexing eta orders by ordinal subscripts
vary in the literature.  This elementary form of saturation is
all that the stage arguments below require.

Regularity is necessary for this particular cutoff conclusion.
If \(\kappa\) is singular, a cofinal sequence of ordinals below
\(\kappa\), of length \(\cf(\kappa)<\kappa\), is unbounded
in \(\No_{<\kappa}\).  Indeed, a surreal of birthday \(\beta\)
is at most the ordinal \(\beta\), the all-positive sign
expansion of that length.  Every candidate upper bound is
therefore exceeded by some member of the cofinal sequence.
The cut consisting of this sequence on the left and the
empty set on the right has no filler in \(\No_{<\kappa}\).
This failure concerns the bounded stage; it does not weaken
the set-cut theorem for the full surreal class.

% ---- hand chunk transition_head

\merge\ The next section is source~12's Section~10. Its dichotomy is new
against the repository; it answers in part the reals report's question
``Other ground orders'' (\cref{swo:sec:answers}). It uses source~12's
self-delimiting code (\cref{swo:st:lem:code}, printed in
\cref{swo:part:setlike}) and its lemmas of
\cref{swo:st:sec:bounded}. Here $S_\kappa$ of source~12 is printed
$\No_{<\kappa}$, and its ``binary coding rank'' is the binary coding length
$\bl$ of the reals report.

% ---- source 12, lines 1018-1205
\section{Birthday cutoffs and the singular transition (source 12)}
\label{swo:st:sec:transition}

Let $\kappa$ be an infinite initial cardinal and put
\[
 \No_{<\kappa}=\{s\in\No:\bd(s)<\kappa\},\qquad
 \mu=|\No_{<\kappa}|=2^{<\kappa}
       =\sup_{\lambda<\kappa}2^{|\lambda|}.
\]
The cardinality identity follows by counting all levels. The supremum is
at least $\kappa$, and summing $\kappa$ levels of size at most that
supremum does not increase it. The reverse inequality is supplied by the
levels themselves. We do not assume $\No_{<\kappa}$ is a field; none of the
following proofs uses arithmetic closure.

\subsection{A scheduler that preserves the first varying coordinate}

\begin{lemma}[Row scheduler]\label[lemma]{swo:st:lem:scheduler}
Suppose a set order $X$ of infinite cardinality $\kappa$ has disjoint
subsets $(X_i)_{i<\kappa}$, each containing a copy of $X$ and each of
cardinality $\kappa$. Suppose their complement is either empty or has
cardinality $\kappa$. Then
\[
 X^\kappa_{\lex}\emb\WO_\kappa(X).
\]
\end{lemma}
\begin{proof}
Fix embeddings $\phi_i:X\to X_i$ and a baseline enumeration of each
$X_i$ in length $\kappa$. For $x\in X$, transpose the label $\phi_i(x)$
with the first label of that baseline, giving an enumeration $r_{i,x}$
of $X_i$ whose first value is $\phi_i(x)$. Its dependence on $x$ is fixed
and uniform. If the complement is nonempty, add a fixed filler row
indexed before all variable rows. The row index order still has type
$1+\kappa=\kappa$.

Well-order row-stage pairs $(i,j)\in\kappa\times\kappa$ first by
$\max(i,j)$ and then lexicographically within a level. This schedule has
type $\kappa$: the predecessors of any pair have cardinality less than
$\kappa$, since both coordinates there are bounded below $\kappa$, while
the entire set of pairs has size $\kappa$. The first appearance of row
$i$ is $(i,0)$, and these first appearances occur in increasing row order.

For an input $(x_i)_{i<\kappa}$, output the $j$th entry of $r_{i,x_i}$
when pair $(i,j)$ is scheduled. The output is an exhaustive bijection
in type $\kappa$. If two inputs first differ in row $i$, every earlier
row agrees, no later row has appeared before $(i,0)$, and the first
output discrepancy is $\phi_i(x_i)$ versus $\phi_i(y_i)$. Its sign is
exactly the sign of the input comparison.
\end{proof}

\begin{lemma}[Surreal cone decomposition]\label[lemma]{swo:st:lem:cones}
If $|\No_{<\kappa}|=\kappa$, then
\[
 \WO_\kappa(\No_{<\kappa})\equim (\No_{<\kappa})^\kappa_{\lex}.
\]
\end{lemma}
\begin{proof}
One embedding is the inclusion of bijective sequences among all
sequences. For the other, partition the sign sequences with a minus into
the cones with prefixes
\[
 p_i=+^i\concat -,\qquad i<\kappa.
\]
They are disjoint, and prefixing by $p_i$ gives a numerical order
isomorphism from $\No_{<\kappa}$ onto that cone. For $i,\alpha<\kappa$,
$i+1+\alpha<\kappa$; this uses that $\kappa$ is an initial cardinal and
is valid even when it is singular. Removing the initial prefix also
leaves a sign sequence of length less than $\kappa$.
The complement consists of all-plus sequences, one for each length
below $\kappa$, so has cardinality $\kappa$. Apply
\cref{swo:st:lem:scheduler}.
\end{proof}

\begin{theorem}[Exact minimal-slice dichotomy]\label{swo:st:thm:transition}
For $\kappa$ and $\mu=2^{<\kappa}$ as above,
\[
 \WO_\mu(\No_{<\kappa})\equim
 \begin{cases}
 \BB_{\kappa\cdot\kappa},
    &\mu=\kappa\text{ and }\cf\kappa<\kappa,\\[2pt]
 \BB_\mu,&\text{otherwise}.
 \end{cases}
\]
Consequently its binary coding rank is respectively
$\kappa\cdot\kappa$ or $\mu$. The exceptional case is exactly the case
of a singular strong-limit cardinal $\kappa$.
\end{theorem}
\begin{proof}
We separate the three mechanisms.

\emph{Case 1: $\mu>\kappa$.} By \cref{swo:st:lem:code}, each member of
$\No_{<\kappa}$ has a prefix-free numerical binary code of length less than
$\kappa$. Pad each such code to length $\kappa$. Concatenating the entry
codes of a $\mu$-long enumeration produces a binary word of length
$\kappa\cdot\mu=\mu$. The equality holds because $\mu$ is an initial
cardinal strictly larger than $\kappa$: every proper partial product
$\kappa\cdot\xi$, $\xi<\mu$, has cardinality below $\mu$.
Thus $\WO_\mu(\No_{<\kappa})\emb\BB_\mu$. Pair coding gives the reverse
embedding.

\emph{Case 2: $\mu=\kappa$ and $\kappa$ is regular.} Use the unpadded,
self-delimiting code $c(s)$. In a sequence of $\kappa$ entries, the sum
of code lengths before any stage $\xi<\kappa$ is less than $\kappa$,
by regularity. The full sum is $\kappa$, since every length is positive.
Consequently concatenation embeds even
$(\No_{<\kappa})^\kappa_{\lex}$ into $\BB_\kappa$. Pair coding again supplies
the reverse embedding for $\WO_\kappa(\No_{<\kappa})$.

\emph{Case 3: $\mu=\kappa$ and $\delta=\cf\kappa<\kappa$.}
The padded single-entry code embeds $\No_{<\kappa}$ into $\BB_\kappa$,
whence
\[
 \WO_\kappa(\No_{<\kappa})\emb
 (\No_{<\kappa})^\kappa_{\lex}\emb\BB_{\kappa\cdot\kappa}.
\]
Choose increasing nonzero ordinals $(\lambda_j)_{j<\delta}$ cofinal in
$\kappa$. Their ordinal sum is $\kappa$. As $\delta\cdot\kappa=\kappa$,
partition the $\kappa$ input coordinates into $\kappa$ consecutive blocks
of length $\delta$. In each block, restrict its $j$th coordinate to sign
sequences of exactly length $\lambda_j$. The lexicographic product of
those coordinates is the binary cube of the sum of their lengths.
Thus, over all the blocks, this restricted product is
\[
 \BB_{\sum_{\xi<\kappa}\sum_{j<\delta}\lambda_j}
   =\BB_{\kappa\cdot\kappa}.
\]
It embeds into $(\No_{<\kappa})^\kappa_{\lex}$, which embeds into
$\WO_\kappa(\No_{<\kappa})$ by \cref{swo:st:lem:cones}. This gives the claimed
equimorphism. The rank assertions follow from \cref{swo:st:lem:strictcube}.

Finally, a singular strong limit satisfies $2^{<\kappa}=\kappa$.
Conversely, suppose $\kappa$ is singular and $2^{<\kappa}=\kappa$.
If $2^\lambda=\kappa$ for some $\lambda<\kappa$, choose
$\rho<\kappa$ with $\rho\geq\lambda,\cf\kappa$. Then
$\kappa\leq2^\rho\leq2^{<\kappa}=\kappa$, so $2^\rho=\kappa$,
contradicting König's inequality $\cf(2^\rho)>\rho$.
Thus $2^\lambda<\kappa$ for every $\lambda<\kappa$.
\end{proof}

\begin{corollary}[A complete cube test for the minimal slice]\label{swo:st:n10.4}
Let $\rho=\kappa\cdot\kappa$ in the exceptional case of
\cref{swo:st:thm:transition}, and let $\rho=\mu$ otherwise. Then
\[
 \BB_\gamma\emb\WO_\mu(\No_{<\kappa})
 \quad\Longleftrightarrow\quad\gamma\leq\rho.
\]
In every case its ordinal embedding spectrum is still precisely
$\{\beta:\beta<\mu^+\}$.
\end{corollary}

\begin{example}[An unconditional singular example]\label{swo:st:n10.5}
Let $\kappa=\beth_\omega=\sup_{n<\omega}\beth_n$. It is a singular
strong-limit cardinal of cofinality $\omega$, independently of GCH.
Thus
\[
 \bl\bigl(\WO_\kappa(\No_{<\kappa})\bigr)=\kappa\cdot\kappa>
 \kappa,\qquad
 \bl\bigl(\WO(\No_{<\kappa})\bigr)=\kappa^+.
\]
All three displayed ordinals must be distinguished. The first two have
the same cardinality, but are different coding lengths; the successor
cardinal is larger still.
\end{example}

\begin{example}[Regular cutoffs]\label{swo:st:n10.6}
For $\kappa=\omega$, $\No_{<\kappa}$ is the order of finite-birthday surreals,
and the minimal slice is equimorphic with $\BB_\omega$.
For $\kappa=\omega_1$, its cardinality is $2^{\aleph_0}$, and the
minimal slice is equimorphic with $\BB_{2^{\aleph_0}}$, whether or not
CH holds. At a strongly inaccessible $\kappa$, it is equimorphic with
$\BB_\kappa$.
\end{example}

\subsection{Why the singular phenomenon is genuinely ordinal}

A set of $\kappa$ coordinates does not license arbitrary reordering of
those coordinates in a lexicographic product. In the singular case, a
short cofinal block of coordinate lengths can already consume ordinal
length $\kappa$; repeating that event $\kappa$ times consumes
$\kappa\cdot\kappa$. Replacing this ordinal sum by its cardinality would
incorrectly collapse the rank back to $\kappa$. The strict-cube lemma
shows that the lost information is mathematically detectable.

The row scheduler addresses a separate possible error: an arbitrary
bijection $\kappa\times\kappa\to\kappa$ need not preserve the first
varying row. Bounding by the maximum coordinate ensures that all earlier
observable output is independent of later varying rows. This is why the
product-to-permutation embedding is valid even at singular $\kappa$.

% ---- hand chunk transition_after

\merge\ The last sentence of \cref{swo:st:n10.4} is \cref{swo:sw:cutoffspectrum}
for the minimal slice. The dichotomy computes binary coding length only;
source~11's question on the exact cut spectrum and point characters of
$\WO_\lambda(\No_{<\kappa})$ without $2^{<\kappa}=\kappa$ (\cref{swo:n16.7})
stays open.

% ---- hand chunk part2
\part{Set-length words}\label{swo:part:words}

% ---- source 11, lines 617-737
\section{Set-length surreal words embed explicitly in the surreal line}
\label{swo:sw:encoding}

Let \(\WordsAll\) be the class of all set-length words
\(f:\alpha\to\No\), and let \(\Words\) be its injective subclass.
The latter codes well-orders of set-sized subsets of \(\No\).
Both are ordered lexicographically with proper prefixes smaller.
Their members are sets, so these are ordinary definable classes
in \(\GB\), not collections whose members are proper classes.
They should not be confused with all global well-orders of
the entire surreal class.

In particular, a word cannot exhaust \(\No\): its range is
a set by Replacement.  Passing from fixed alphabets to
\(\Words\) consequently changes which carrier is being
well-ordered from one point of the lexicographic order to
another.  This is a mathematically natural extension, but
it answers a different question from ordering all class
well-orders on a single proper-class carrier.
Making this distinction before forming any collection of
codes is what prevents a hidden ascent from classes to
collections of classes.

The following construction uses neither a global enumeration of
\(\No\) nor Global Choice.  It uses only ordinal concatenation
of sign expansions, whose general properties are discussed in
\cite[Section 3.2]{BagayokoHoeven}.

\begin{theorem}[Explicit sign coding]\label{swo:sw:code}
There is an explicitly defined order embedding
\(\mathcal E:\WordsAll\longrightarrow\No\).
It maps the empty word to \(0\), maps every nonempty word to
a positive surreal, and restricts to an order embedding
of \(\Words\).
Singleton words give embeddings in the opposite direction.
\end{theorem}

\begin{proof}
For a sign expansion \(s\), let \(D(s)\) duplicate each sign:
replace \(-\) by \((-,-)\), and \(+\) by \((+,+)\).
For \(x\in\No\), put
\[
               C(x)=(+)\cat D(s_x)\cat(+,-).
\]
The length here is the ordinal
\(1+2\cdot\bd(x)+2\).
If \(s_x,s_y\) first differ at a sign, their doubled blocks
have the same comparison.  If \(s_x\) is a proper prefix of
\(s_y\), the terminal block \((+,-)\) of \(C(x)\) is compared
with the next doubled block of \(C(y)\).  We have
\[
                  (-,-)<(+,-)<(+,+).
\]
This is precisely the comparison of a missing sign with a
negative or positive continuation.  Thus \(C\) preserves
the surreal order.  Moreover, two distinct values of \(C\)
cannot be prefixes of one another: the same first differing
block supplies an actual differing sign before either code ends.

For \(f:\alpha\to\No\), concatenate these blocks in domain order:
\[
        \mathcal E(f)=\mathop{\cat}_{i<\alpha}C(f(i)).
\]
Replacement and ordinal summation make this a set-length sign
expansion.  If two words first differ at \(i\), their preceding
blocks coincide and their blocks at \(i\) have a first differing
sign.  Prefix freedom prevents that difference from being
displaced into a subsequent block.  The comparison therefore
agrees with the comparison of the original entries.
If \(f\) is a proper prefix of \(g\), the expansion
\(\mathcal E(g)\) continues \(\mathcal E(f)\) with the initial
\(+\) of the next block, so \(\mathcal E(f)<\mathcal E(g)\).
This also proves the assertions about empty and nonempty words.
Finally, \(x\mapsto\langle x\rangle\) embeds \(\No\) into
both word classes.
\end{proof}

Accordingly, \(\WordsAll\), \(\Words\), and \(\No\) are pairwise
equimorphic as class linear orders.  This conclusion asserts
embeddings both ways, not isomorphism.  Section~\ref{swo:sw:characters}
exhibits an obstruction to isomorphism that persists after
the empty word is removed.

\begin{proposition}[Birthday bounds]\label{swo:sw:codebounds}
The coding has the following bounds.
\begin{enumerate}
\item If \(\kappa\) is regular infinite, \(\alpha<\kappa\), and
      \(f(i)\in\No_{<\kappa}\) for every \(i<\alpha\), then
      \(\mathcal E(f)\in\No_{<\kappa}\).
\item For any infinite cardinal \(\kappa\), put
      \(\lambda=2^{<\kappa}\).  If \(f\) is an arbitrary
      injective word over \(\No_{<\kappa}\), then
      \(\mathcal E(f)\in\No_{<\lambda^+}\).
\end{enumerate}
\end{proposition}

\begin{proof}
The birthday of the image is
\[
       \bd(\mathcal E(f))
           =\sum_{i<\alpha}\bigl(1+2\cdot\bd(f(i))+2\bigr),
\]
with ordinal addition and multiplication.
For the first assertion, this sums fewer than \(\kappa\)
ordinals below \(\kappa\), and regularity keeps the sum
below \(\kappa\).
For the second, injectivity gives \(\card{\alpha}\leq\lambda\);
each block has cardinality at most \(\lambda\).
The whole sign expansion has cardinality at most \(\lambda\),
so its ordinal length is below \(\lambda^+\).
\end{proof}

The first bound need not hold at a singular cutoff.
Choose distinct ordinals \(\gamma_i<\kappa\),
\(i<\cf(\kappa)\), cofinal in \(\kappa\), and regard them as
surreal entries.  The word has length less than \(\kappa\),
but the lengths of its code blocks are unbounded in \(\kappa\),
so its encoded birthday is at least \(\kappa\).
Also, injectivity is essential to the unrestricted second bound:
constant words over a one-element alphabet can have every
ordinal length.

% ---- source 11, lines 739-836
\section{Local characters and set suprema of word orders}
\label{swo:sw:characters}

In this section \(\mathcal W\) denotes either \(\WordsAll\)
or \(\Words\).
Deleting the set of entries of an existing prefix causes no
difficulty when choosing a new coordinate: apply
Lemma~\ref{swo:lem:delete}.
We say that a segment has no set cofinal subset rather than
assigning it a cardinal cofinality that does not exist.

\begin{theorem}[Exact local character]\label{swo:sw:character}
Let \(f\in\mathcal W\) have domain \(\alpha\).
\begin{enumerate}
\item If \(\alpha\) is a nonzero limit ordinal, then
      \(\cf(\mathcal W_{<f})=\cf(\alpha)\).
\item If \(\alpha\) is a successor ordinal, then
      \(\mathcal W_{<f}\) has no set cofinal subset.
\item For every \(\alpha\), the segment
      \(\mathcal W_{>f}\) has no set coinitial subset.
\end{enumerate}
The empty word is the minimum of \(\mathcal W\).
\end{theorem}

\begin{proof}
Suppose first that \(\alpha\) is a nonzero limit.
Its proper prefixes \(f\restriction\beta\), \(\beta<\alpha\),
are cofinal below \(f\).  A word \(g<f\) is either a proper
prefix, or differs from \(f\) first at some coordinate \(\xi\);
in the latter case it lies below every prefix long enough to
include coordinate \(\xi\).  A cofinal sequence of indices
in \(\alpha\) therefore supplies a cofinal subset of size
\(\cf(\alpha)\).
Conversely, to each \(g<f\) assign a prefix length strictly
large enough that the corresponding prefix exceeds \(g\).
Fewer than \(\cf(\alpha)\) such lengths are bounded below
\(\alpha\), so the corresponding words cannot be cofinal.

Next let \(\alpha=\beta+1\), write \(p=f\restriction\beta\),
and take a set \(A<f\).  For each member \(a\) properly
extending \(p\), its next coordinate satisfies
\(a(\beta)<f(\beta)\).  Choose an unused surreal \(d\)
strictly above all these coordinates and strictly below
\(f(\beta)\).
Then \(h=p\cat\langle d\rangle\) is below \(f\).
It exceeds every member of \(A\) extending \(p\);
the remaining members are at most \(p\), so they also lie
below \(h\).  Thus \(A\) is not cofinal.

Finally, take a set \(B>f\).
Collect \(b(\alpha)\) for those \(b\in B\) properly extending
\(f\), and choose an unused \(d\) below all these coordinates.
The word \(h=f\cat\langle d\rangle\) lies above \(f\).
It lies below the members of \(B\) extending \(f\).
Every other member of \(B\) first differs earlier with a
larger coordinate, and consequently exceeds every extension
of \(f\), including \(h\).  Hence \(f<h<B\).
\end{proof}

\begin{corollary}\label{swo:sw:notiso}
Neither \(\WordsAll\) nor \(\Words\), with or without its
empty word, is order-isomorphic to \(\No\).
Each has nontrivial set suprema, but a set with no minimum
cannot have an infimum in either word order.
\end{corollary}

\begin{proof}
Take an injective word \(f\) of length \(\omega\).
Its nonempty finite prefixes increase to \(f\), and \(f\)
is their supremum by the preceding theorem.
No surreal can be the supremum of a set lying strictly
below it: the surreal set-cut property would insert a
smaller upper bound.  This distinguishes the orders even
after deleting their minimum.
If \(f\) were the infimum of a set \(B\) not containing \(f\),
part (3) would produce \(h\) with \(f<h<B\), contradicting
the definition of infimum.
\end{proof}

These local obstructions coexist with very rich intervals.
Given \(f<_{\lex}g\) with a first differing coordinate, vary
that coordinate through a surreal interval strictly between
the two values and disjoint from the set of preceding entries.
If \(f\) is a proper prefix of \(g\), vary a new coordinate
below the next entry of \(g\) instead.
The surreal set-cut property provides such an interval after
excluding the preceding entries.
Every surreal interval \((u,v)\) is order-isomorphic to
\(\No\), for example by
\[
 t\longmapsto u+\frac{v-u}{2}
                      \left(1+\frac{t}{1+|t|}\right).
\]
Thus every open word interval contains a copy of \(\No\),
and Theorem~\ref{swo:sw:code} then places a copy of the entire
word order inside it.
Interval universality therefore does not imply that set cuts
are filled.

% ---- source 11, lines 838-944
\section{The complete obstruction to filling a set cut}
\label{swo:sw:wordcuts}

The preceding prefix example is the only nontrivial obstruction.
Here and below, \(A<B\) means that every member of \(A\) is
strictly below every member of \(B\); either set may be empty.

\begin{theorem}[Word-cut classification]\label{swo:sw:cutclassification}
A set cut \(A<B\) in \(\mathcal W\) has no interpolant
if and only if one of the following holds:
\begin{enumerate}
\item \(B\) contains the empty word;
\item \(B\) has a minimum \(b\) of nonzero limit length,
      and \(A\) is cofinal in \(\mathcal W_{<b}\).
\end{enumerate}
If neither obstruction occurs, an interpolant can be chosen
to have nonzero successor length.
\end{theorem}

\begin{proof}
Either displayed obstruction plainly prevents an interpolant.
For the converse, perform the following common-prefix search.
Start with \(p\) empty.  Maintain these invariants:
members of \(A\) not extending \(p\) lie below \(p\);
members of \(B\) not extending \(p\) lie above every extension
of \(p\).  Here extending allows equality.

If \(p\in B\), stop.  Otherwise, with \(\delta=\dom(p)\),
form the sets of next coordinates
\[
 A(p)=\{a(\delta):a\in A,\ p\text{ is a proper prefix of }a\},
 \quad
 B(p)=\{b(\delta):b\in B,\ p\text{ is a proper prefix of }b\}.
\]
Every member of \(A(p)\) is at most every member of \(B(p)\),
since the corresponding words satisfy \(a<b\).
If these coordinate sets are disjoint, all these inequalities
are strict.  Choose a surreal \(c\) strictly between them,
outside \(\ran(p)\).  Then \(p\cat\langle c\rangle\)
interpolates the original cut by the invariants, and the
construction is complete.

If the two coordinate sets meet, their intersection consists
of a unique element \(c\): two distinct common elements would
violate the preceding inequalities.
Replace \(p\) by \(p\cat\langle c\rangle\).
The invariants persist because the discarded \(A\)-branches
have smaller next coordinates and the discarded \(B\)-branches
have larger ones.
In the injective-word case, \(c\) is unused, since it is
witnessed by an injective word extending the previous prefix.
At a limit stage take the union of the prefixes.
It is a set word, is injective when required, and satisfies
the same invariants.  A member of \(B\) cannot have become a
proper prefix of this union: the search would already have
stopped upon reaching that member.

This procedure requires only set recursion.
Choose an ordinal strictly above the domains of all words
in \(A\cup B\).  The search cannot continue taking common
next coordinates up to that length, because there would be
no words supplying those coordinates.  It therefore terminates.
The remaining possibility is termination at \(p\in B\).
The invariants show that \(p=\min B\).
If its length is zero, the first obstruction holds.

Termination cannot occur at a successor length
\(\delta+1\).
The last appended coordinate was chosen from the intersection,
so some member of \(A\) extends this successor prefix.
Such a word is at least \(p\), contrary to \(A<B\).
Hence a nonzero terminal length \(\alpha\) is a limit.
For every \(\beta<\alpha\), a member of \(A\) witnessed the
common coordinate at stage \(\beta\); it extends
\(p\restriction(\beta+1)\) and lies below \(p\).
These witnesses exceed prefixes of arbitrarily large length.
Since the proper prefixes are cofinal below \(p\), \(A\)
is cofinal below \(p\).  This is exactly the second obstruction.
\end{proof}

\begin{corollary}[Successor words]\label{swo:sw:successors}
The suborder of \(\mathcal W\) consisting of words of nonzero
successor length fills every set cut.
In \(\GBC\), it is order-isomorphic to \(\No\).
\end{corollary}

\begin{proof}
Neither obstruction in Theorem~\ref{swo:sw:cutclassification}
can occur when the upper set consists of successor words.
The construction returns another successor word, proving
the set-cut property.
Under Global Choice, enumerate each of the two proper classes
by a set-like class well-order and perform back-and-forth.
At every ordinal stage the previously matched elements form
a set; their images specify a set cut, and an unused filler
extends the partial order isomorphism.
The union is an isomorphism.  This is the class analogue of
the usual eta-order back-and-forth, with its choice assumption
made explicit; compare the class framework in
\cite{Ehrlich2012}.
\end{proof}

The distinction is precise.  Deleting the empty word alone
does not remove the prefix-supremum obstruction.  Restricting
to successor lengths does remove it.  The explicit embeddings
of Theorem~\ref{swo:sw:code} require no Global Choice, while the
back-and-forth identification just given uses it.

% ---- source 11, lines 946-1008
\section{Saturation at a regular stage}\label{swo:sw:stagesaturation}

For regular infinite \(\kappa\), restrict to successor-length
injective words of length below \(\kappa\), with entries in
\(\No_{<\kappa}\).
This order is \(\eta_\kappa\).
Indeed, a cut with fewer than \(\kappa\) words has all its
word lengths bounded below \(\kappa\).
The common-prefix search in the preceding proof consequently
stops below \(\kappa\), and Lemma~\ref{swo:sw:stageeta} supplies
the last coordinate while excluding the fewer than \(\kappa\)
entries already used.
Proposition~\ref{swo:sw:codebounds} embeds this order into \(\No_{<\kappa}\);
singleton words embed \(\No_{<\kappa}\) back.
When \(2^{<\kappa}=\kappa\), both orders have cardinality
\(\kappa\), and the ordinary back-and-forth of length \(\kappa\)
makes them isomorphic.

A related saturation statement applies to full permutations.

\begin{proposition}\label{swo:sw:permutationeta}
Let \(X\) be an \(\eta_\kappa\)-order of cardinality \(\kappa\),
where \(\kappa\) is regular infinite.
Then \(\WO_\kappa(X)\) is \(\eta_\kappa\).
Its global and local cofinalities and coinitialities are
all exactly \(\kappa\).
\end{proposition}

\begin{proof}
For a cut \(A<B\) of total cardinality below \(\kappa\),
apply the common-coordinate search to its permutations.
If both sides are nonempty, the first-disagreement indices
of the pairs in \(A\times B\) have supremum below \(\kappa\).
The search cannot continue sharing a coordinate past this
bound; if either active side becomes empty, it stops sooner.
The common prefix therefore has length below \(\kappa\).
The \(\eta_\kappa\) property chooses an unused coordinate
strictly between the active coordinate sets.
Complete the resulting prefix to a \(\kappa\)-permutation:
its remaining entries and remaining positions both have
cardinality \(\kappa\), and its remaining interval of positions
has order type \(\kappa\).
The completed permutation fills the cut.

Saturation excludes cofinal or coinitial subsets of size below
\(\kappa\), both globally and on either side of a point.
Conversely, Theorem~\ref{swo:sw:setspectrum} excludes monotone
sequences of length \(\kappa^+\).
If one of these segments had no cofinal subset of size at
most \(\kappa\), recursively choosing a point above all earlier
choices would produce such an increasing sequence.
The decreasing argument is identical.
Thus each of the stated characters equals \(\kappa\).
\end{proof}

In particular this proposition applies to \(\No_{<\kappa}\)
when \(2^{<\kappa}=\kappa\).
The resulting permutation order has cardinality \(2^\kappa\),
so it is substantially larger than its alphabet while retaining
the same local characters.  It is not
\(\eta_{\kappa^+}\): that stronger cut property would construct
an increasing \(\kappa^+\)-sequence, contrary to
Theorem~\ref{swo:sw:setspectrum}.

% ---- hand chunk part3
\part{Global choice and all class well-orders}\label{swo:part:classes}


% ---- source 11, lines 1014-1125
\section{Global choice and well-ordering the surreal class}
\label{swo:cf:existence}

The distinction between a set of surreal numbers and the entire surreal class
changes the existence question.  Here \(\GB\) has the sets of \(\mathrm{ZF}\),
without a choice axiom, and \(\GC\) asserts the existence of a class function
choosing an element of every nonempty set.  A class well-order is a strict
linear class relation for which every nonempty available subclass has a least
element.  Its underlying objects are still sets.  The relation itself can be
a proper class, and its predecessor classes need not be sets.  These are the
size distinctions emphasized in the repository's foundations discussion
\cite{ProveItFoundations}.

Throughout this section, \(\Hall\), \(\Hsl\), and \(\WO(X)\) when
\(X\) is a proper class are metanotation for the indicated collections
of class relations.  We quantify over their members as class variables;
we do not assert that those collections are GB classes.  When \(X\) is
a set, \(\WO(X)\) retains its ordinary set meaning from the preceding
section.  All constructions below are uniform formulas in the stated class
parameters.  In particular, restricting such a construction to an already
given uniformly indexed family produces an ordinary class relation on the
indexing domain.

Concretely, a family indexed by a class \(I\) is supplied by one class
\(\mathcal R\subseteq I\times\No\times\No\), whose slice at \(i\)
defines the relation \(R_i\).  Writing this family does not make \(R_i\)
a set or a member of \(I\): the index \(i\) is a set, and the slice is
a class defined with that parameter.  This distinction permits uniform
families of proper class relations while preserving the usual GB typing
rules.  No exhaustive indexing family is assumed in this section.

\begin{theorem}[Existence and global choice]
\label{swo:cf:choice}
Over \(\GB\), the following statements are equivalent:
\begin{enumerate}
  \item \(\GC\) holds;
  \item there is a class well-order of \(\No\);
  \item there is a set-like class well-order of \(\No\);
  \item there is a class bijection \(\Ord\to\No\).
\end{enumerate}
Here set-like means that the predecessors of every element form a set.
\end{theorem}

\begin{proof}
We give the implication from an arbitrary well-order to global choice in
detail.  It also explains why set choice need not be assumed in the base
theory.  Let \(W\) be a class well-order of the sign-sequence presentation of
\(\No\).  By ordinary set-valued transfinite recursion, with the fixed class
parameter \(W\), construct a set relation \(w_\alpha\) well-ordering
\(V_\alpha\), for every ordinal \(\alpha\).

Set \(w_0=\varnothing\).  Suppose \(w_\beta\) has been constructed.  There is
a unique ordinal \(\theta\) and a unique order isomorphism
\[
 e:\theta\longrightarrow(V_\beta,w_\beta).
\]
For \(A\subseteq V_\beta\), let \(s_A\) be the sign sequence of length
\(\theta\) whose value at \(\xi\) is \(+\) exactly when \(e(\xi)\in A\).
The map \(A\mapsto s_A\) is a bijection between \(\mathcal P(V_\beta)\)
and the set of all sign sequences of length \(\theta\).  Define
\[
 A\;w_{\beta+1}\;B\quad\Longleftrightarrow\quad s_A\;W\;s_B.
\]
This is a set relation by separation with a class parameter, and it
well-orders \(V_{\beta+1}=\mathcal P(V_\beta)\).  No selection of an
enumeration is involved: the enumeration of an already specified set
well-order is unique.

At a nonzero limit \(\lambda\), order \(V_\lambda\) first by rank, and
within rank \(\gamma\) by the restriction of \(w_{\gamma+1}\):
\[
 x\;w_\lambda\;y
 \quad\Longleftrightarrow\quad
 \rk(x)<\rk(y)\quad\text{or}\quad
 \bigl(\rk(x)=\rk(y)=\gamma\ \text{ and }\ x\;w_{\gamma+1}\;y\bigr).
\]
Since \(\lambda\) is a limit, every required \(\gamma+1\) is below
\(\lambda\).  A nonempty subset of \(V_\lambda\) has a least occurring
rank; the specified well-order on that rank then supplies its least member.
Thus the recursion continues.  Its stages are sets, even though their union
over all ordinals is a class.  This is precisely the ordinary GB recursion
theorem, rather than a recursion whose individual values are proper classes.

For a nonempty set \(A\), every element of \(A\) belongs to
\(V_{\rk(A)}\).  Choose the \(w_{\rk(A)}\)-least element of \(A\).
The recursion and this prescription define a single class choice function,
establishing \(\GC\).

Conversely, global choice implies choice for sets.  For each ordinal
\(\alpha\), the set of well-orders of the birthday layer
\(S_\alpha=\{-,+\}^{\alpha}\) is nonempty.  Apply the global choice
function to this set, and order \(\No\) first by birthday and then by the
selected order within the birthday layer.  This is a class well-order.
Every predecessor class is contained in
\(\bigcup_{\beta\leq\alpha}S_\beta\) for some \(\alpha\), hence is a
set.  This proves the set-like version directly.

Finally, for any set-like well-order \(R\) of the proper class \(\No\),
assign to \(x\) the ordinal order type of its predecessor set.  These ranks
form an initial subclass of \(\Ord\).  That subclass cannot be a set:
otherwise the inverse rank map, by class replacement on a set domain, would
make \(\No\) a set.  It is therefore all of \(\Ord\), and the inverse
rank map is the required bijection.  Conversely, a class bijection from
\(\Ord\) transports its well-order to a set-like well-order of \(\No\).
\end{proof}

The connection between surreal enumerations and global choice is discussed
by Hamkins \cite{HamkinsChoice}; the general equivalences for global choice
are collected in \cite{HamkinsGC}.  The explicit rank recursion above proves
the assertion over the stated choiceless GB base.  It does not assert that a
global choice class is definable without parameters in the first-order
language of set theory.  Definability is an additional requirement.

% ---- hand chunk choice12_head

\merge\ \emph{Second proof over $\GB+\mathrm{AC}$ (source~12).} Source~12
proves the equivalence by coding transitive closures on ordinals; set choice
is used only to code each $\operatorname{TC}(\{x\})$ on an ordinal. Its
equivalence includes non-set-like well-orders, as source~11's does. Its
subsection on birthday order follows its theorem.

% ---- source 12, lines 429-494
\subsection{Source 12's route over GB with set choice}
\label{swo:st:sec:choice}

\begin{theorem}[Choice equivalence]\label{swo:st:thm:choice}
Over GB plus set-level AC, the following are equivalent:
\begin{enumerate}[label=(\roman*)]
\item global choice;
\item a class bijection $\Ord\to\No$;
\item a set-like class well-order of $\No$;
\item a class well-order of $\No$, not necessarily set-like.
\end{enumerate}
\end{theorem}
\begin{proof}
Global choice supplies a set-like global well-order of the universe, and
hence of the proper class $\No$. A proper set-like well-order has type
$\Ord$, giving (i)$\Rightarrow$(iii)$\Rightarrow$(ii). A bijection in (ii)
pulls back ordinal order, giving (iii), and (iii)$\Rightarrow$(iv) is
immediate.

We prove the substantive reverse implication by coding sets. Suppose
$R$ well-orders $\No$. Let $x$ be a nonempty set and let
$T_x=\operatorname{TC}(\{x\})$; it contains $x$ and all its elements.
By set choice, there is a least initial ordinal $\lambda$ equinumerous
with $T_x$. There is a fixed definable ordinal $\tau_\lambda$ large
enough to code a binary relation on $\lambda$, together with one
distinguished index, by a binary sequence of length $\tau_\lambda$.
For example, concatenate a $\lambda\cdot\lambda$-long relation table
and a $\lambda$-long singleton marker.

Consider the \emph{set} of codes for well-founded extensional relations
on $\lambda$ whose Mostowski collapse has range $T_x$, and whose marked
index collapses to $x$. This set is nonempty. Convert $0,1$ to $-,+$,
view the codes as surreals, and take its $R$-least member. Decode it, and
choose the element of $x$ represented by the least ordinal index that
belongs to the marked node in the coded membership relation.

The result is an element of $x$, uniquely determined by $x$ and the class
parameter $R$. Elementary class comprehension gives the graph of this
choice operation on all nonempty sets. It is a global choice function.
Notice that only a set of codes for each $x$ was minimized: set-likeness
of $R$ was not needed.
\end{proof}

\begin{remark}[The explicit set-choice hypothesis]\label{swo:st:n4.2}
The converse proof uses set choice to code $T_x$ on an ordinal. We do not
claim the same equivalence over GB with no set-level choice. A purported
stronger choice equivalence must address that missing coding step.
\end{remark}

\subsubsection{Birthday is not a canonical global well-order}

Ordering first by birthday and then by numerical value does not work.
For $n<\omega$, let $s_n$ be the length-$\omega$ sign sequence consisting
of $n$ minuses followed by pluses forever. Then
\[
 \bd(s_n)=\omega,\qquad s_{n+1}<s_n.
\]
The whole descending sequence lies in one birthday level. Simplicity
itself is well-founded but is only a partial order; two incompatible sign
sequences need not be comparable by simplicity.

With global choice one may choose a well-order on every birthday level
and concatenate them in birthday order. The result is set-like because
all earlier levels, and every initial part of the current level, form a
set. The noncanonical choices have not disappeared: they are exactly the
kind of information accounted for by \cref{swo:st:thm:choice}.

% ---- hand chunk choice12_answered

\begin{remark}[\merge; source~12's question answered]\label{swo:rem:answered-choice}
Dated note, 2 October 2026. Source~12's \cref{swo:st:n4.2} declines the
choiceless case, and its research question on the exact weak-choice strength
(\cref{swo:st:n14.5}) asks which principle is equivalent to a class well-order
of the sign-sequence surreals over $\GB$ without set choice. Source~11's
\cref{swo:cf:choice} answers it: global choice itself. The ``missing coding
step'' that \cref{swo:st:n4.2} names is avoided, not supplied: the well-orders
$w_\alpha$ of $V_\alpha$ are built by set-valued recursion from the given
well-order of the sign sequences of length $\theta=\otp(V_\beta,w_\beta)$, a
subset $A\subseteq V_\beta$ being coded by the sign sequence $s_A$ of length
exactly $\theta$, so no ordinal coding of a transitive closure has to be
chosen. The proof thus produces a coherent sequence of well-orders of the sets
$V_\alpha$, in particular of the power sets of ordinals, which is the
comparison the question suggests. The same mechanism, at the level of
sets, is the ordinal-graph coding of the birthday-cutoff report
(\code{hset:thm:choice}, whose range over $\ZF$ is the class of sets with
well-orderable transitive closure) \cite{ProveItHset}; that theorem concerns
a different statement and is unchanged. Hamkins's question whether surreal
universality implies global choice \cite{HamkinsUnivChoice} is not affected.
\end{remark}

% ---- source 11, lines 1127-1249
\section{Lexicographic comparison without class ordinals}
\label{swo:cf:comparison}

The set-sized enumeration method studied for the real numbers
\cite{ProveItReals} suggests comparing the first unequal labels.  Arbitrary
class well-orders require a formulation that also permits a proper class of
common earlier labels.  No ordinal enumeration of that entire common prefix
is needed.

For a strict class well-order \(R\) on a fixed class \(X\), write
\[
 P_R(x)=\{y\in X:yRx\}.
\]
A common labeled prefix of \(R,S\) is a subclass initial in both orders on
which their restrictions agree.  The adjective labeled matters: agreement
here is equality of the actual elements and relations, rather than the
existence of an abstract order isomorphism.

\begin{lemma}[Unique first disagreement]
\label{swo:cf:first}
If \(R\ne S\) are class well-orders of the same class \(X\), there are
unique distinct \(a,b\in X\) such that
\[
 P_R(a)=P_S(b)=C,
 \qquad R\mathbin{\upharpoonright}C=S\mathbin{\upharpoonright}C.
\]
The subclass \(C\) is their maximal common labeled prefix.
\end{lemma}

\begin{proof}
The class \(D=\{x:P_R(x)\ne P_S(x)\}\) is nonempty, since predecessor
classes determine a strict order.  Let \(a\) be its \(R\)-least element
and put \(C=P_R(a)\).  For every \(y\in C\),
\(P_R(y)=P_S(y)\).  In particular, \(ySa\): otherwise \(aSy\), which
would imply \(aRy\), contradicting \(yRa\).  Thus \(C\subseteq P_S(a)\).
Also, if \(zSy\) with \(y\in C\), then \(zRy\), so \(z\in C\).
Therefore \(C\) is \(S\)-initial as well as \(R\)-initial, and the
restrictions of the orders to it agree.

The inclusion \(C\subseteq P_S(a)\) is strict because \(a\in D\).
Let \(b\) be the \(S\)-least element outside \(C\).  The strict
inclusion shows that \(bSa\), hence \(b\ne a\).  Initiality and the
minimality of \(b\) give \(P_S(b)=C\).

For uniqueness, suppose \(a',b',C'\) satisfy the displayed conditions with
\(a'\ne b'\).  Every element of \(C'\) has identical predecessor
classes in the two orders.  But \(b'Sa'\), so the predecessor classes of
\(a'\) differ.  Consequently \(a'\) is exactly the \(R\)-least element
of \(D\); hence \(a'=a\), \(C'=C\), and \(b'=b\).  Any common
prefix properly extending \(C\) would have to include incompatible next
labels, so \(C\) is maximal.
\end{proof}

Fix now a native linear order \(<_X\) on the labels.  For distinct \(R,S\),
define \(R<_{\lex}S\) when the labels supplied by
Lemma~\ref{swo:cf:first} satisfy \(a<_Xb\).

\begin{theorem}[Raw lexicographic order]
\label{swo:cf:lex}
This rule is a strict linear order on the class well-orders of \(X\),
understood as a definable relation on class variables.  In particular it
defines \(\Hall\) for \(X=\No\) with its native numerical order.
\end{theorem}

\begin{proof}
Existence and uniqueness of first disagreement give totality and asymmetry;
the relation is irreflexive by definition.  To verify transitivity, suppose
\(R<_{\lex}S<_{\lex}T\).  Write \((a,b)\) for the first unequal
labels of \(R,S\), and \((c,d)\) for those of \(S,T\).  Their common
prefixes are \(C=P_S(b)\) and \(D=P_S(c)\), respectively, and
\(a<_Xb\), \(c<_Xd\).

If \(bSc\), then \(C\subsetneq D\).  The orders \(S,T\) agree
through the label \(b\), so the first unequal labels of \(R,T\) are
\(a,b\), with common prefix \(C\).  Thus \(R<_{\lex}T\).
If \(cSb\), then \(D\subsetneq C\); now \(R,S\) agree through
\(c\), and the first unequal labels of \(R,T\) are \(c,d\), again
giving the conclusion.  If \(b=c\), the prefixes coincide and the next
labels of \(R,T\) are \(a,d\).  Since
\(a<_Xb=c<_Xd\), the conclusion follows in this case as well.
\end{proof}

This proof compares labeled relations directly.  It does not compare
arbitrary class well-orders up to isomorphism, and makes no appeal to a
collection of class ordinal types.  The defining conditions use only set
quantifiers with the given class relations as parameters: equality of
predecessor classes means a universal statement about their set elements.

For example, the predicate that \(a,b\) are the first unequal labels can
be written without quantifying over a common-prefix class:
\[
 \begin{split}
 a\ne b\quad&\land\quad\forall z\,(zRa\leftrightarrow zSb)\\
 &\land\quad\forall x\,\forall y\,
 \bigl((xRa\land yRa)\longrightarrow(xRy\leftrightarrow xSy)\bigr).
 \end{split}
\]
Adding \(a<_Xb\) and existentially quantifying the two labels gives the
lexicographic predicate.  Consequently its restriction to the indexed
family \(\mathcal R\) described above is a GB class of pairs of indices.
If the indexing contains duplicate orders, this relation orders the distinct
relations, and becomes a strict linear order on indices whenever the
indexing is injective.  Passing instead to abstract isomorphism types would
discard exactly the label information on which the construction depends.

\begin{proposition}[Convex prefix cones]
\label{swo:cf:cones}
Fix a well-ordered subclass \((C,P)\) of \(X\).  The collection of all
class well-orders of \(X\) extending \((C,P)\) as an initial labeled
segment is convex in the raw lexicographic order.
\end{proposition}

\begin{proof}
Suppose \(R,S\) extend the prefix and \(R<_{\lex}T<_{\lex}S\).
Let \(D\) be the first-disagreement prefix of \(R,T\).  The subclasses
\(C,D\) are initial in \(R\), and therefore inclusion-comparable.
If \(C\subseteq D\), then \(T\) extends \((C,P)\), as required.
If \(D\subsetneq C\), the next \(R\)-label lies in \(C\); it is
also the next \(S\)-label after \(D\).  Accordingly, \(T\) compares
with \(R\) and \(S\) in the same direction, contradicting the strict
inequalities.  This reasoning permits \(C\) and \(D\) to be proper
classes.
\end{proof}

% ---- hand chunk prefix12_head

\merge\ \emph{Second route (source~12): the largest common labelled prefix.}
Source~12's Section~3 proves \cref{swo:cf:first,swo:cf:lex} differently.
Instead of the class $D$ of points whose predecessor classes differ, it
defines the largest common labelled prefix $C(R,S)$ by an elementary formula
and proves directly that it is initial in both orders and contains every
common labelled initial segment. Source~11's shipped check compares its
predecessor rule with finite lexicographic order on all 266,272 pairs of
permutations of at most six points
(\path{data/11-raw-orders-finite_checks_results.txt}); source~12's checks the
common-prefix formula on 15,017 finite pairs
(\path{data/12-singular-finite_checks.json}).

% ---- source 12, lines 317-427
\subsection{Source 12's route: the largest common labelled prefix}
\label{swo:st:sec:commonprefix}

The central observation is that comparing two \emph{labelled}
well-orders is easier than comparing arbitrary abstract class order types.
Fix a class $X$ equipped with an ambient strict linear order $<$, and let
$R,S$ be class well-orders of the same carrier $X$.

\begin{definition}[Common labelled prefix]\label{swo:st:n3.1}
Define $C(R,S)\subseteq X$ by putting $x\in C(R,S)$ if
\begin{align*}
 \pred_R(x)&=\pred_S(x),\\
 \forall u,v\in\pred_R(x),\qquad uRv&\ \Longleftrightarrow\ uSv.
\end{align*}
Equality of the predecessor classes is expanded as a quantification over
sets $u\in X$. Thus this is an elementary formula with class parameters
$R,S$; it is not a quantification over all possible common prefixes.
\end{definition}

\begin{lemma}[Maximal-prefix lemma]\label[lemma]{swo:st:lem:common}
$C(R,S)$ is an initial segment of both well-orders, their restrictions to
it agree, and it contains every common labelled initial segment. If
$R\ne S$, it is proper, and the least elements
\[
 a=\min_R(X\setminus C(R,S)),\qquad
 b=\min_S(X\setminus C(R,S))
\]
exist and are distinct.
\end{lemma}
\begin{proof}
Suppose $x\in C(R,S)$ and $yRx$. Equality of the predecessor classes puts
$y$ before $x$ in $S$ as well. For $zRy$, transitivity gives $zRx$;
agreement of $R,S$ below $x$ then gives $zSy$. Conversely, $zSy$ implies
$zSx$, hence $zRx$, and agreement below $x$ yields $zRy$. The orders agree
below $y$ because they agree below $x$. Therefore $y\in C(R,S)$.
The symmetric argument proves initiality in $S$.

If $u,v\in C(R,S)$ and $uRv$, equality of the predecessor classes of $v$
gives $uSv$. Thus the restrictions agree. If $D$ is any common labelled
initial segment and $x\in D$, both predecessor classes of $x$ lie in $D$
and agree, as do the restricted relations. Hence $D\subseteq C(R,S)$.
If $C(R,S)=X$, equality of every predecessor class implies $R=S$.

For $R\ne S$, class well-foundedness supplies $a,b$. By initiality and
minimality, $\pred_R(a)=C(R,S)$ and $\pred_S(b)=C(R,S)$. If $a=b$, these
predecessor classes agree and the two relations agree on them. The
defining formula would put $a$ in $C(R,S)$, a contradiction.
\end{proof}

\begin{definition}[Full lexicographic comparison]\label{swo:st:n3.3}
For distinct $R,S$, put $R<_{\lex}S$ precisely when $a<b$, with $a,b$
from \cref{swo:st:lem:common}. Identical relations are equal in this order.
\end{definition}

\begin{theorem}[Full class lexicographic order]\label{swo:st:thm:fulllex}
In GBC, the preceding comparison is a strict linear ordering of the
predicate of all class well-orders of a fixed linearly ordered class
$X$. Its definition uses only elementary class comprehension. In
particular, ETR and the comparability of abstract class well-order types
are not required.
\end{theorem}
\begin{proof}
Irreflexivity and trichotomy follow from the maximal-prefix lemma and
trichotomy in $X$. For transitivity suppose $R<_{\lex}S<_{\lex}T$.
The classes $C(R,S)$ and $C(S,T)$ are initial segments of the same
well-order $S$, so they are comparable by inclusion.

If $C(R,S)$ is a proper subset of $C(S,T)$, the first $S$-element $b$
after $C(R,S)$ belongs to $C(S,T)$. Thus $T$ agrees with $S$ through that
element. The first disagreement of $R,T$ is exactly the earlier
$R,S$ disagreement, with the same two labels $a<b$. Consequently
$R<_{\lex}T$. The opposite proper-inclusion case uses the earlier
$S,T$ disagreement and is identical in form.

If the two common prefixes are equal to $C$, let $a,b,c$ be the first
labels after $C$ in $R,S,T$. The hypotheses give $a<b<c$. Hence $R,T$
also first disagree after $C$, and $R<_{\lex}T$. Every construction of
$C$, and every assertion about its next elements, expands into set
quantifiers with the given class relations as parameters.
\end{proof}

\begin{remark}[Why this does not prove class order-type comparability]\label{swo:st:n3.5}
In general one asks whether two arbitrary class well-orders admit an
initial-segment isomorphism. ETR is a sufficient hypothesis for that
principle \cite[Theorem 6]{HamkinsWoodin}. Here no such isomorphism is constructed.
Only their \emph{identical labelled} portions are compared. As soon as
the labels disagree, the fixed ambient order decides the answer.
Confusing these two tasks would impose an unnecessary stronger axiom.
\end{remark}

\begin{example}[A disagreement after every ordinary ordinal]\label{swo:st:n3.6}
Fix $a<b$ in $\No$. Well-order $\No\setminus\{a,b\}$ in type $\Ord$.
Append first $a,b$, and in a second order append $b,a$.
The orders agree throughout their initial $\Ord$-block and differ only
after it. No ordinary ordinal indexes that first disagreement, but
\cref{swo:st:thm:fulllex} compares them immediately: their common labelled
prefix is $\No\setminus\{a,b\}$.
\end{example}

\begin{corollary}[Agreement with enumeration lexicography]\label{swo:st:n3.7}
For set-sized exhaustive ordinal enumerations, and for class bijections
$\Ord\to X$, the common-prefix definition agrees with first-difference
lexicography.
\end{corollary}
\begin{proof}
Before the first unequal entry, exactly the same labels occur in exactly
the same order. Their ranges form the common labelled initial segment.
No later label can belong to that segment by initiality and the distinct
next labels. In the $\Ord$ case the disagreement is a nonempty class of
ordinals; restricting it below any one witness gives its least member.
\end{proof}

% ---- hand chunk prefix12_after

\merge\ Source~12's \cref{swo:st:n3.5} is source~11's remark after
\cref{swo:cf:lex} and in \cref{swo:sec:formalization}: comparability of
abstract class well-orders by initial-segment isomorphisms follows from $\ETR$
\cite{HamkinsWoodin}, and neither labelled comparison needs it. Its
\cref{swo:st:n3.6} is source~11's \cref{swo:cf:ordtwo}, and its
\cref{swo:st:n3.7} is the corresponding sentence after
\cref{swo:gl:uniform}.

% ---- source 11, lines 1251-1362
\section{The canonical initial set-like part}
\label{swo:cf:decomposition}

Every class well-order of \(\No\) has a canonical initial segment of
type \(\Ord\), even when its entire order is not set-like.  This gives
a structural description of all well-orders using ordinary class functions.

\begin{theorem}[Canonical initial part]
\label{swo:cf:core}
For \(R\in\Hall\), let
\[
 I_R=\{x\in\No:P_R(x)\text{ is a set}\}.
\]
Then \(I_R\) is a proper initial subclass of the domain or is the entire
domain; in either case it is a proper class.  Its induced order is set-like
and has a unique order enumeration
\[
 F_R:\Ord\longrightarrow I_R.
\]
The order \(R\) is the concatenation of this enumeration with its induced
well-order on \(\No\setminus I_R\).
\end{theorem}

\begin{proof}
If \(x\in I_R\) and \(yRx\), then \(P_R(y)\subseteq P_R(x)\),
so separation makes \(P_R(y)\) a set.  Thus \(I_R\) is initial.
Suppose it were a set.  Since \(\No\) is proper, there would be an
\(R\)-least element \(x\) outside it.  Every predecessor of \(x\)
would then belong to \(I_R\), and initiality would place every member
of \(I_R\) before \(x\).  Hence \(P_R(x)=I_R\), a contradiction.

The induced order on \(I_R\) has the same predecessor sets as \(R\).
The ordinal-rank map therefore identifies it with an initial subclass of
\(\Ord\).  As in Theorem~\ref{swo:cf:choice}, a set-sized range would make
\(I_R\) a set by replacement.  Its range is consequently all of
\(\Ord\).  The inverse map is unique, and initiality gives the asserted
concatenation.
\end{proof}

\begin{example}[Two proper blocks]
\label{swo:cf:two_blocks}
Under \(\GBC\), fix \(e:\Ord\to\No\) bijective.  The disjoint classes
\(A=\{e(2\alpha):\alpha\in\Ord\}\) and
\(B=\{e(2\alpha+1):\alpha\in\Ord\}\) partition \(\No\), where the
product is ordinal multiplication \(2\cdot\alpha\).  Order each block
by its displayed ordinal index and place all of \(A\) before all of
\(B\).  The resulting well-order has class type \(\Ord+\Ord\).
Its canonical initial part is exactly \(A\): a point of \(B\) has
the whole proper class \(A\) among its predecessors.  Hence an
\(\Ord\)-enumeration of the canonical part is not, in general, an
enumeration of the entire well-order.  Both blocks consist of ordinary
surreal labels; their size affects their positions in the chosen well-order,
not the status of individual surreal numbers as sets.
\end{example}

Let \(\Jinj\) denote the metacollection of injective class functions
\(F:\Ord\to\No\), with lexicographic comparison at their first ordinal
of disagreement.  Such a first disagreement exists for distinct functions:
after choosing any witness, minimize within the set of earlier ordinals.
Put \(Y_F=\No\setminus\operatorname{ran}(F)\).

\begin{proposition}[Convex fibers]
\label{swo:cf:fibers}
Assume \(\GBC\).  Every \(F\in\Jinj\) occurs as \(F_R\) for some
\(R\in\Hall\).  The fiber over \(F\) is a convex copy of
\(\WO(Y_F)\).  Distinct fibers are ordered by the first unequal values
of their functions \(F\).
\end{proposition}

\begin{proof}
The range of an injection from \(\Ord\) is proper, since a set could
not contain an injective image of all ordinals.  Choose a well-order \(Q\)
of \(Y_F\), using the restriction of any fixed global well-order, and
put the sequence \(F\) before \(Q\).  Each \(F(\alpha)\) has exactly
the set \(F``\alpha\) of predecessors.  Each residual point, however,
has the proper class \(\operatorname{ran}(F)\) among its predecessors.
Thus the resulting order has canonical function exactly \(F\).

All orders with that function extend the same initial labeled order on
\(\operatorname{ran}(F)\), and conversely every extension has that
canonical function.  Proposition~\ref{swo:cf:cones} proves convexity.
Restriction to \(Y_F\) and concatenation with \(F\) are inverse
operations; first disagreement inside the residual orders proves that they
preserve lexicographic comparison.  Finally, if \(F\ne G\), their first
ordinal disagreement occurs before either residual order begins.  It
therefore determines the comparison of every member of one fiber with
every member of the other.
\end{proof}

We may summarize this result as the lexicographic-sum description
\[
 \Hall\ \cong
 \sum_{F\in(\Jinj,<_{\lex})}\WO(Y_F).
\]
This is uniform metanotation for the preceding theorem, not the assertion
that its indexing hyperclass or all its fibers have been collected into a
GB class.  The set-like fragment \(\Hsl\) consists exactly of the fibers
with \(Y_F=\varnothing\), or equivalently the exhaustive enumerations.
For finite \(Y_F\) of size \(m\), its fiber is a convex block of
\(m!\) orders, with the usual lexicographic order on permutations.

\begin{remark}[Singleton fibers and iteration]
\label{swo:cf:singleton}
A singleton fiber need not be set-like.  If \(F\) omits exactly one
surreal \(z\), its unique completion puts \(z\) after the entire
\(\Ord\)-sequence; it has class order type \(\Ord+1\).  Conversely,
every exhaustive \(F\) gives a singleton set-like fiber.  Extracting the
initial part once also does not justify transfinite iteration of this
operation through class-sized tails.  Such an iteration would require a
separate specification of its indexing order, its objects, and its
recursion principle; none is asserted here.
\end{remark}

% ---- source 11, lines 1364-1480
\section{All jumps arise from finite residual permutations}
\label{swo:cf:jumps}

The full order has jumps despite the extreme density of the native surreal
line.  A proper common prefix can already have used every label between the
two labels at which comparison is decided.

\begin{lemma}[Extreme enumerations]
\label{swo:cf:extreme}
For a class \(Z\) with a native linear order, a class well-order is
lexicographically greatest among all well-orders of \(Z\) exactly when
it is the native reverse order.  It is lexicographically least exactly
when it is the native forward order.
\end{lemma}

\begin{proof}
If a well-order places \(y\) before \(z\) with \(y<z\), transpose
these two labels throughout the order.  The result is another class
well-order.  Its first changed label is \(z\) in place of \(y\), so it
is lexicographically greater.  A greatest order must therefore be strictly
decreasing in the native order, which determines the entire relation.
Conversely, in the reverse native order the next label at any first
disagreement is the greatest remaining label, so every different
well-order is smaller.  The least-order assertion follows by reversing
the native order.  The argument includes empty and singleton domains.
\end{proof}

The lemma is an existence criterion as well as a characterization: a
greatest enumeration exists only when the native reverse order itself is
well-founded.  It is therefore not permissible to use ``list the remaining
labels in decreasing order'' as an unconditional completion rule.  For
example, on a copy of the natural numbers with its usual native order, the
increasing enumeration is least, but there is no greatest well-order.
Finite residual domains have both extrema, and the adjacency theorem shows
that precisely such domains can support a jump after a common prefix.

\begin{theorem}[Exact adjacency criterion]
\label{swo:cf:adjacency}
Let \(R<_{\lex}S\) be class well-orders of \(X\).  Let \(C\) be their
first-disagreement prefix, let \(a<b\) be their respective next labels,
and put \(Y=X\setminus C\).  They are adjacent if and only if all
three conditions hold:
\begin{enumerate}
  \item there is no \(y\in Y\) with \(a<y<b\);
  \item the \(R\)-tail on \(Y\setminus\{a\}\) is the native reverse
        order;
  \item the \(S\)-tail on \(Y\setminus\{b\}\) is the native forward
        order.
\end{enumerate}
These conditions force \(Y\) to be finite.  Thus adjacency is exactly
ordinary finite permutation adjacency after a possibly proper class
common prefix.
\end{theorem}

\begin{proof}
Suppose first that the pair is adjacent.  A label \(y\in Y\) strictly
between \(a,b\) could be placed immediately after \(C\), followed by
the restriction of \(R\) to the other remaining labels.  This would
give a well-order strictly between \(R,S\); thus condition~1 holds.

Any lexicographically greater replacement of the \(R\)-tail after
\(a\) would also give an intermediate order: its next label after
\(C\) would still be \(a<b\).  Hence this tail is greatest, and
Lemma~\ref{swo:cf:extreme} gives condition~2.  A lexicographically smaller
replacement of the \(S\)-tail after \(b\) similarly gives
condition~3.  These modifications use restrictions and finite
transpositions of existing class relations; they require no new choice
principle.

The overlap \(Z=Y\setminus\{a,b\}\) is now well-ordered both forward
and backward in the native order.  If \(Z\) were not finite, repeatedly
take its least element after deleting the finitely many already selected.
Ordinary recursion on \(\omega\) and replacement produce a set of
strictly increasing selected elements.  That set has no greatest member,
contradicting reverse well-ordering.  Thus \(Z\), and hence \(Y\),
is finite.  This argument does not assume that \(Z\) was a set in
advance.

Conversely, suppose the three conditions hold and let
\(R<_{\lex}T<_{\lex}S\).  Convexity of the prefix cone forces \(T\)
to extend \(C\).  Let \(t\in Y\) be its next label.  A label below
\(a\) makes \(T<R\), and a label above \(b\) makes \(T>S\);
condition~1 excludes intermediate labels.  If \(t=a\), its remaining
tail cannot exceed the greatest tail of \(R\), so \(T\leq R\).
If \(t=b\), its tail cannot precede the least tail of \(S\), so
\(T\geq S\).  Every possibility contradicts strict betweenness.
\end{proof}

\begin{example}[An adjacent pair of type \(\Ord+2\)]
\label{swo:cf:ordtwo}
Assume \(\GBC\), choose surreal labels \(a<b\), and give
\(C=\No\setminus\{a,b\}\) a set-like well-order \(P\).  Consider
\[
 R=P+(a,b),\qquad S=P+(b,a).
\]
Each has class order type \(\Ord+2\), and \(R,S\) are adjacent.
They have the same canonical initial enumeration, and constitute its
entire two-point fiber.  The notation \(\Ord+2\) describes a class
concatenation; it is not a set ordinal.
\end{example}

For \(X=\No\), every adjacent pair has a cofinite, hence proper class,
common prefix.  Neither endpoint can be set-like.  Moreover, if two
set-like orders are given, their common prefix is a set.  The deletion
property of the surreal line leaves an unused label strictly between their
next labels, and that set prefix can be completed to a set-like well-order.
For completeness, this extension does not require an additional recursion
principle.  Retain the prescribed set prefix, put the new intermediate
label next, and then append the restriction of a fixed set-like global
well-order of \(\No\) to the unused labels.  A predecessor class in the
appended part is contained in the union of the prescribed set prefix and
a predecessor set of the fixed order, together with the one inserted
label.  It is therefore a set.  Every label occurs once, and the
concatenation is a well-order, so the completion belongs to \(\Hsl\).
Thus the induced order \(\Hsl\) is dense, while \(\Hall\) has the
finite terminal blocks classified above.  Allowing all well-orders changes
the local order structure in a concrete and unavoidable way.

% ---- hand chunk geometry12_head

\merge\ \emph{Second route (source~12): residual decomposition and
adjacency.} Source~12's Section~8 proves \cref{swo:cf:adjacency} in the form
``the residual class is finite and the residual permutations are consecutive'';
the two criteria are equivalent, since a tail that is both well-ordered and
reverse well-ordered is finite. Its \cref{swo:st:prop:residual} is the
description of one prefix cone given by \cref{swo:cf:cones,swo:cf:fibers}. Its
\cref{swo:st:n8.4} adds that interior points of the finite blocks with
$n\ge3$ are isolated in the order topology and that the set-like suborder is
not dense in the full order, and its last paragraph adds the order-reversing
involution induced by negation.

% ---- source 12, lines 773-861
\subsection{Source 12's route: unrestricted class well-orders and finite-tail geometry}
\label{swo:st:sec:fullgeometry}

The full order $\Hall(\No)$ contains $\Hsl(\No)$ as an ordered
subcollection, but it is not dense. The effect is not a failure of density
of the surreal alphabet; it is exhaustion of almost all labels by a
proper-class prefix.

\begin{definition}[A prefix cylinder in the full order]\label{swo:st:n8.1}
Let $C\subsetneq X$ carry a fixed well-order $P$. Consider all class
well-orders of $X$ that have $(C,P)$ as a labelled initial segment.
Write this cylinder as $[P]_X$ and put $D=X\setminus C$.
\end{definition}

\begin{proposition}[Residual-order decomposition]\label[proposition]{swo:st:prop:residual}
$[P]_X$ is convex and is order-isomorphic to the lexicographic order of
all well-orders of the residual carrier $D$, with its inherited ambient
order. If $D$ is a set, this is the ordinary set order $\WO(D)$,
with the same notation extended to finite carriers.
If $|D|=n<\omega$, the cylinder has exactly $n!$ elements.
\end{proposition}
\begin{proof}
Restriction to $D$ and ordered concatenation with $P$ are inverse
operations. Concatenation of two class well-orders is a class well-order:
a nonempty subclass has its minimum in the first part if it meets that
part, and otherwise in the second. The maximal common prefix after $P$
is exactly the common prefix of the tails, so comparisons are preserved.

If a third order differs before completing $P$, its first discrepancy
with every extension of $P$ is the same and places it on the same side of
all of them. This proves convexity. The last assertion counts
permutations of a finite carrier.
\end{proof}

\begin{theorem}[Exact adjacency criterion]\label{swo:st:thm:adjacency}
Let $R<_{\lex}S$ be two class well-orders of $X$, and let
$C=C(R,S)$, $D=X\setminus C$. They are adjacent in the full order if and
only if $D$ is finite and their residual permutations are consecutive in
the ordinary finite lexicographic order.
\end{theorem}
\begin{proof}
Let $a<b$ be their first residual labels. If some $d\in D$ satisfies
$a<d<b$, begin the residual order with $d$ and well-order the remaining
labels by a restriction of $R$. This gives an order strictly between.
So an adjacent pair can have no such $d$.

If the tail of $R$ after $a$ has two labels $u$ before $v$ with $u<v$ in
the ambient order, interchange the labels $u,v$ at their two positions.
The new well-order is larger than $R$, and still smaller than $S$ because
its first residual label remains $a$. Thus adjacency forces the $R$ tail
to be numerically decreasing. Dually, the tail of $S$ after $b$ must be
numerically increasing.

On $D\setminus\{a,b\}$, both the ambient numerical order and its reverse
are therefore well-orders. Such a carrier is finite: an infinite
well-order contains an increasing $\omega$-sequence, contradicting
well-foundedness of its reverse. Hence $D$ is finite.

Conversely, if $D$ is finite, convexity in \cref{swo:st:prop:residual} says that
every order between $R,S$ must be another permutation of $D$ after $C$.
There is none exactly when the two permutations are consecutive.
\end{proof}

\begin{corollary}[Proper-class prefixes are responsible]\label{swo:st:n8.4}
Adjacent pairs in $\Hall(\No)$ share a proper-class prefix.
There are convex blocks of every factorial size $n!$. For $n\geq3$,
interior elements of such blocks are isolated in the order topology.
The set-like suborder is not dense in the full order.
\end{corollary}
\begin{proof}
A set prefix cannot exhaust all but finitely many surreals. Conversely,
well-order $\No\setminus D$ in type $\Ord$ for any finite $D$, and append
all its permutations. For $|D|\geq3$, choose two nonconsecutive
permutations in the resulting finite convex block. Their nonempty open
interval contains only orders with a proper-class predecessor segment,
so it contains no set-like order.
\end{proof}

The full order still has no endpoints and no set-indexed cofinal or
coinitial family, by the first-entry argument of
\cref{swo:st:prop:characters}. It contains every class linear order, by
\cref{swo:st:thm:local}. But neither global density nor interpolation for all
set-indexed cuts survives: an adjacent pair already gives a two-sided
cut that cannot be filled.

Pointwise surreal negation gives an order-reversing involution on both
global orders: transport a well-order by the bijection $x\mapsto -x$.
This preserves enumeration type and set-likeness while reversing the
ambient comparison of the first unequal labels.

% ---- source 12, lines 1299-1334
\section{Numerically well-ordered subclasses and other interpretations (source 12)}
\label{swo:st:sec:subclasses}

The positive surreal ordinals form a numerically increasing class copy
of $\Ord$, but this is not a well-order of all of $\No$. More generally,
\cref{swo:st:lem:universalNo} implies that every class well-order is isomorphic
to a numerically well-ordered subclass of $\No$. That subclass need not
be set-like in its inherited numerical order.

For an explicit example, the class
\[
 \left\{-\frac{1}{\alpha+1}:\alpha\in\Ord\right\}\cup\{0\}
\]
has numerical order type $\Ord+1$. The arithmetic here is surreal
arithmetic on the embedded ordinal values; the reciprocal reverses
positive order and the minus sign reverses it back. The final point has
a proper class of predecessors. Putting a positive ordinal block after
this negative block gives a numerical copy of $\Ord+\Ord$.

For a \emph{set} $A\subseteq\No$, numerical well-ordering has the usual
ordinal order type. Its birthday bound places it inside some $\No_{<\kappa}$.
The class of all such set-sized well-ordered subsets is a legitimate
class of sets. Its members are not exhaustive enumerations of a fixed
carrier, however, so one must specify what happens when one enumeration
is a prefix of another. The ``shorter first'' convention and the surreal
$-<\End<+$ convention are different choices; neither can be silently
borrowed from the prefix-free exhaustive setting.

Finally, ordering binary relation tables rather than enumerations is
another legitimate but different construction. If one fixes a set-like
well-order of $\No\times\No$, compares the characteristic functions of
relations at the first unequal table entry, and restricts to well-orders,
one gets a table-dependent order. A first changed relation bit is not
generally the first changed enumerated label. Even on three points the
two lexicographic orders can differ. The present article uses labelled
initial segments, not relation-table lexicography.

% ---- hand chunk subclasses_after

\merge\ Source~12's shipped check records a three-point example on which
relation-table lexicography and enumeration lexicography differ
(\path{data/12-singular-finite_checks.json}, key
\code{relation\_table\_vs\_enumeration\_example}).

% ---- hand chunk part4
\part{Set-like global well-orders}\label{swo:part:setlike}

% ---- source 11, lines 1486-1538
\section{The lexicographic structure of global enumerations}
\label{swo:gl:section}

The set-like case has a particularly strong structure theorem.  Although the
full collection of global enumerations is a higher-order object, one ordinary
class of set codes already meets every cut specified by set-indexed families.
That class, with its induced order, is isomorphic to the surreal line.  The
distinction between a set-indexed family and an ordinal-indexed class family is
essential: an explicit omitted-range construction shows exactly where the
interpolation theorem stops.

Throughout this section we work in $\GBC$ and fix a class bijection
\[
 e_0:\Ord\longrightarrow\No.
\]
The existence of such a baseline is supplied by global choice; the relevant
equivalences are discussed in \cite{HamkinsGC}.  We use the surreal
set-cut property, its preservation under deleting a set
(Lemma~\ref{swo:lem:delete}), and the universality and uniqueness statements of
Proposition~\ref{swo:prop:universal}.  These are order-theoretic properties; the
field operations on $\No$ play no role in the main construction.

\subsection{Class parameters and uniform families}

\begin{definition}\label{swo:gl:uniform}
For this section, $\Jinj$ denotes the collection of class injections
$F:\Ord\to\No$, and $\Hsl$ denotes the subcollection of class bijections.
For distinct $F,G\in\Jinj$, put
\[
 d(F,G)=\min\{\alpha\in\Ord:F(\alpha)\ne G(\alpha)\},
 \qquad
 F<_{\lex}G\ \Longleftrightarrow\ F(d(F,G))<G(d(F,G)).
\]
A \emph{uniformly set-indexed family} $(F_i)_{i\in I}$ is given by a set $I$
and one class relation whose $i$th row is the graph of $F_i$.
\end{definition}

These notations are interpreted relationally in two-sorted class theory.
An individual $F$ is a class, not a possible element of a GB class.  Thus
``for all $F\in\Hsl$'' abbreviates quantification over class functions with the
stated property.  Likewise, a uniformly indexed family is a single class
of triples $(i,\alpha,x)$, not a literal set containing proper classes.
The first-difference definition is legitimate: choose any coordinate of
disagreement and minimize within the set of preceding ordinals.  The usual
first-difference proof gives a strict linear comparison on these objects.

Uniformity matters whenever a set bound is taken.  For example, if two
uniformly set-indexed families have different rows across every pair, Class
Replacement makes their cross first-difference ordinals a set.  No principle
of choosing an arbitrary family of classes is being assumed.  Finite lists of
given class parameters are automatically uniform, and the families of
set-coded objects constructed below have a fixed evaluation formula, so their
set-indexed subfamilies are uniform as well.

% ---- source 11, lines 1540-1601
\subsection{A class of bounded-support codes}

\begin{definition}\label{swo:gl:bounded-codes}
A \emph{normalized bounded-support code} is a set permutation
$\sigma:\lambda\to\lambda$ of an ordinal $\lambda$ satisfying
\[
 \lambda=\sup\{\xi+1:\xi<\lambda\text{ and }\sigma(\xi)\ne\xi\}.
\]
The empty permutation is allowed.  Extend $\sigma$ to a class permutation
$\widehat\sigma$ of $\Ord$ by the identity outside $\lambda$, and evaluate
the code as
\[
 E_\sigma=e_0\circ\widehat\sigma.
\]
Write $\Pbd=\Pbd(e_0)$ for the ordinary class of these codes, ordered by the
lexicographic comparison of their evaluations.
\end{definition}

Normalization removes the otherwise harmless duplication obtained by adding
fixed points to the end of a code.  Every class permutation with bounded
support has exactly one normalized code: restrict it to the supremum of one
plus its moved coordinates.  Its support is invariant under the permutation,
so this restriction is a permutation of the indicated ordinal.  Consequently
distinct normalized codes evaluate to distinct enumerations, and the induced
order on $\Pbd$ is a genuine class linear order.  The class is proper, since
the transpositions of $0$ with arbitrary nonzero ordinals have distinct codes.

Here and below, identifying a code with its evaluated enumeration is a
notational convenience.  It does not convert the evaluated proper-class graph
into a set.  In particular, the assertion that $\Pbd$ is an ordinary class
refers to the codes, while assertions comparing its members with arbitrary
members of $\Hsl$ use the common evaluation relation.

\begin{lemma}[Completion of an arbitrary set prefix]\label{swo:gl:completion}
Every injective set function $p:\delta\to\No$, with $\delta$ an ordinal,
is the initial segment of an enumeration represented in $\Pbd$.
The completion can be chosen uniformly in $p$.
\end{lemma}

\begin{proof}
Let $q=e_0^{-1}\circ p$.  This is a set injection from $\delta$ into
$\Ord$.  Choose an infinite initial cardinal $\gamma$ strictly larger than
$\delta$ and than $\sup\ran(q)$.  Both complements
\[
 \gamma\setminus\delta
 \quad\hbox{and}\quad
 \gamma\setminus\ran(q)
\]
have cardinality $\gamma$, since fewer than $\gamma$ elements were removed.
A bijection between the complements therefore extends $q$ to a permutation
of $\gamma$.  Extend that permutation by the identity above $\gamma$, and
normalize its code.  Its evaluation agrees with $p$ at every coordinate below
$\delta$.  To make the construction uniform, fix the least permissible
$\gamma$ and then the least complement bijection under a fixed global
well-order.  These choices select sets; they do not require a choice principle
for families of proper classes.
\end{proof}

The lemma is stronger than merely extending $p$ to some global well-order.
Even an arbitrarily long set-sized prefix can be realized by changing the
baseline only on a bounded set of coordinates.  The bound need not be small:
it must contain the baseline indices of all values prescribed by $p$.

% ---- hand chunk extend12

\paragraph{Source 12's prefix extension.}\label{swo:st:sec:setlike}
\merge\ Source~12's Section~6 (``Set-like global well-orders: interpolation
and universality'') proves the results of this section by the same arguments,
except for its prefix extension, printed next, which is a different
construction: it appends the restriction of the baseline order to the
complement of the prefix, so its completion need not have bounded support.
Source~12's interpolation theorem, local characters and local universality are
recorded in notes below; its avoidance lemma and its universality lemmas are
printed in \cref{swo:sec:prelim}.

% ---- source 12, lines 563-578
We now restrict to $\Hsl(\No)$, the class-bijection predicate. All statements
about set- or class-indexed collections of its members mean uniform
families, unless a larger ambient universe is explicitly specified.

\begin{lemma}[Extending a set prefix]\label[lemma]{swo:st:lem:extend}
Every injective set sequence $p:\alpha\to\No$ extends to a class
bijection $\Ord\to\No$. The extension can be chosen uniformly from $p$
and one fixed baseline bijection $B$.
\end{lemma}
\begin{proof}
Let $A=\ran(p)$. Restrict the baseline order to $\No\setminus A$.
It is a proper-class set-like well-order, so it has its canonical
increasing enumeration in type $\Ord$. Append that enumeration to $p$.
The class order $\alpha+\Ord$ is canonically isomorphic to $\Ord$.
The construction is set-valued recursion with $B$ as a parameter.
\end{proof}

% ---- source 11, lines 1603-1668
\subsection{Interpolation of every uniformly set-indexed cut}

\begin{theorem}[Prefix interpolation]\label{swo:gl:interpolation}
Let $(A_i)_{i\in I}$ and $(B_j)_{j\in J}$ be uniformly set-indexed families
in $\Jinj$, with
\[
 A_i<_{\lex}B_j\qquad(i\in I,\ j\in J).
\]
Either index set may be empty.  There is an injective set prefix $p$ such that
every member of $\Jinj$ extending $p$ is strictly above every $A_i$ and
strictly below every $B_j$.  In particular, a member of $\Pbd$ lies between
the two families.
\end{theorem}

\begin{proof}
If both index sets are nonempty, Class Replacement gives a set of ordinals
$d(A_i,B_j)$.  Choose $\beta$ so that
\[
 d(A_i,B_j)<\beta\qquad(i\in I,\ j\in J).
\]
If one side is empty, put $\beta=0$.  We build an injective prefix $p$, keeping
track of the rows that still agree with it.  Call these rows \emph{active}.
The invariant is that every discarded $A_i$ is below every extension of $p$,
and every discarded $B_j$ is above every extension of $p$.

Suppose $p$ has length $\alpha$ and the construction has not stopped.  Let
\[
 L=\{A_i(\alpha):A_i\text{ is active}\},\qquad
 R=\{B_j(\alpha):B_j\text{ is active}\}.
\]
These are sets.  Since active cross pairs agree before $\alpha$, the hypothesis
implies $\ell\le r$ whenever $\ell\in L$ and $r\in R$.  Thus $L\cap R$
has at most one member: two distinct common values, taken in reverse order
as a left and a right value, would contradict this inequality.

If $L\cap R$ is empty, all cross inequalities are strict.  By
Lemma~\ref{swo:lem:delete}, choose
\[
 z\in\No\setminus\ran(p),\qquad L<z<R.
\]
This notation includes an empty left or right side; if both are empty, choose
any fresh value.  Append $z$ and stop.  Every previously active left row is
now strictly below every extension of the new prefix, and every previously
active right row is strictly above it.  The invariant handles the already
discarded rows, proving the desired conclusion.

If $L\cap R=\{z\}$, append $z$ and retain as active exactly the rows whose
next value is $z$.  The new value is fresh because each surviving row is
injective.  A discarded left value is strictly below $z$, since $z\in R$;
a discarded right value is strictly above $z$, since $z\in L$.  Hence the
invariant is preserved.  At a limit stage take the union of the preceding
prefixes; its length is a set ordinal and it remains injective.

The procedure cannot continue with both sides active through a prefix of
length $\beta$.  Such active rows would agree below $\beta$, contrary to the
bound on every cross first difference.  Therefore it stops no later than the
next step, giving a set prefix.  Lemma~\ref{swo:gl:completion} supplies a
bounded-support completion.  Notice that surjectivity of the input rows was
never used; this is why the theorem was stated for $\Jinj$.
\end{proof}

All recursion histories in this proof are sets.  Even when the input rows are
proper classes, their restrictions to a set ordinal are sets, and their next
values form sets because the index sets are sets.  The proof therefore uses
ordinary set-valued transfinite recursion with class parameters, not $\ETR$.
The same observation applies to the back-and-forth argument below.

% ---- hand chunk n_interp12

\begin{note}[Source~12, Theorem~6.3, ``Set-indexed interpolation'']\label{swo:st:thm:etasl}
\emph{If $I,J$ are sets and $(E_i)_{i\in I}$, $(F_j)_{j\in J}$ are uniform
families in $\Hsl(\No)$ with $E_i<_\lex F_j$ for all $i,j$, there is
$G\in\Hsl(\No)$ with $E_i<_\lex G<_\lex F_j$; empty sides are allowed, so
$\Hsl(\No)$ is dense without endpoints in the indexed-family sense.} This is
\cref{swo:gl:interpolation}, by the same forced-prefix construction; source~12
completes the prefix by \cref{swo:st:lem:extend}, so its interpolant need not
lie in the bounded-support core.
\end{note}

% ---- source 11, lines 1670-1695
\begin{corollary}\label{swo:gl:pbd-no}
The class order $\Pbd$ is an $\eta_{\Ord}$ order and is order-isomorphic to
$(\No,<)$.  Moreover, it meets every strict cut in $\Hsl$ specified by
uniformly set-indexed families.
\end{corollary}

\begin{proof}
A set of codes is uniformly evaluable, so Theorem~\ref{swo:gl:interpolation}
applies to every separated pair of sets in $\Pbd$.  Empty sides give the
absence of endpoints, and singleton sides give density.  We already know
that $\Pbd$ is a proper class.  Proposition~\ref{swo:prop:universal} therefore
identifies it with $\No$.

Explicitly, global choice enumerates each of the two classes by ordinals.
Alternately extend a set-sized partial isomorphism to cover the next element
of one enumeration and then the next element of the other.  The lower and
upper images at either step are sets, so the cut property supplies the new
point.  The union of this ordinary set-valued recursion is an isomorphism.
The last assertion is the direct specialization of
Theorem~\ref{swo:gl:interpolation} to bijective input rows.
\end{proof}

This is an isomorphism of the class of bounded-support codes with the surreal
line.  It is not an isomorphism between $\No$ and the full higher-order
collection $\Hsl$.  Different baselines produce classes with the same abstract
order type, but the argument does not single out a canonical isomorphism.

% ---- source 11, lines 1697-1728
\begin{corollary}[The class of realizations of a set cut]
\label{swo:gl:realizers}
Under the hypotheses of Theorem~\ref{swo:gl:interpolation}, the ordinary class
\[
 K=\{p\in\Pbd:(\forall i\in I\ A_i<_{\lex}p)
                    \text{ and }(\forall j\in J\ p<_{\lex}B_j)\}
\]
is itself order-isomorphic to $\No$.
\end{corollary}

\begin{proof}
It is an ordinary class by comprehension with the two uniform input relations
as parameters.  Suppose $C,D\subseteq K$ are sets with $C<D$.  Add the
evaluated codes of $C$ to the lower input family and those of $D$ to the upper
input family.  All cross inequalities remain strict: members of $K$ already
lie between the original sides.  Theorem~\ref{swo:gl:interpolation} supplies a
member of $K$ between $C$ and $D$.  Thus $K$ has the set-cut property,
including empty sides.  It cannot be a set, since applying that property
with lower side all of $K$ would give a member above itself.  The uniqueness
statement of Proposition~\ref{swo:prop:universal} now gives the isomorphism.
\end{proof}

Thus the interpolation theorem produces a full proper-class order of
available choices, even when the input constraints are numerous and occur
at arbitrarily complicated set-ordinal coordinates.  In particular, no set
of codes can be order-dense in $\Hsl$.  Given such a set, choose a point
strictly above all its evaluations, and then a second, larger point.  The
open interval between the two contains no code from the proposed dense set.
The same argument applies within $K$: a set cannot be dense even in the
class of realizations of one specified set cut.  The ordinary class
$\Pbd$ is therefore a natural scale for density here, while an exhaustive
class of set codes will be ruled out by a separate diagonal argument.

% ---- source 11, lines 1730-1753
\subsection{Universality, cylinders, and local cofinality}

\begin{proposition}\label{swo:gl:universality}
The surreal line embeds explicitly into $\Pbd$.  Consequently every ordinary
class linear order embeds into $\Pbd$ and hence into $\Hsl$, in the uniform
evaluation sense.
\end{proposition}

\begin{proof}
For $x\in\No$, transpose the values $e_0(0)$ and $x$ in $e_0$, and denote
the resulting enumeration by $E_x$.  Its coordinate permutation has support
contained in $\{0,e_0^{-1}(x)\}$, so it has a code in $\Pbd$, and
\[
 E_x(0)=x,\qquad x<y\ \Longleftrightarrow\ E_x<_{\lex}E_y.
\]
For the general assertion, compose this embedding with the class-order
universality of $\No$ from Proposition~\ref{swo:prop:universal}.
\end{proof}

Thus all set ordinals, the class order $\Ord$, its reverse, and class orders
such as $\Ord+\Ord$ occur as suborders of the lexicographic comparison.
This statement concerns the order \emph{between} enumerations.  Each
individual enumeration in $\Hsl$ still presents a well-order of type
$\Ord$; its lexicographic environment need not be well-ordered.

% ---- hand chunk n_univ12

\begin{note}[Source~12, Theorem~6.7, ``Local class universality and cylinders'']\label{swo:st:thm:local}
\emph{Every class linear order has a uniform order-embedding into $\Hsl(\No)$
and into every nonempty open interval of $\Hsl(\No)$; for a set-sized
injective prefix $p$, the cylinder $[p]$ is convex and isomorphic to
$\Hsl(\No)$ by uniform transformations.} This is \cref{swo:gl:universality}
and \cref{swo:gl:cylinders}, by the same constructions ($x\mapsto E_x$
composed with \cref{swo:st:lem:universalNo}; a cylinder through
$\No\setminus\ran(p)\cong\No$, \cref{swo:st:lem:delete}). Source~12 lists
the examples ``all set ordinals, their reverses, $\Ord+1$, $\Ord+\Ord$, all
set-sized lexicographic powers of $\No$ formed as class orders'' and adds that
``every class linear order'' still means an order whose elements are sets; it
does not include an unrestricted higher collection such as all class
well-orders themselves.
\end{note}

% ---- source 11, lines 1755-1785
\begin{proposition}[Prefix cylinders]\label{swo:gl:cylinders}
For an injective set prefix $p$, let
\[
 C_p=\{E\in\Hsl:E\upharpoonright\dom(p)=p\}.
\]
The cylinder $C_p$ is convex and order-isomorphic to $\Hsl$, with the
isomorphism interpreted as a uniform transformation of class parameters.
Every nonempty open interval in $\Hsl$, and every interval specified by a
uniformly set-indexed strict cut, contains a whole such cylinder.
\end{proposition}

\begin{proof}
An enumeration that first departs from $p$ at a coordinate of $p$ lies on the
same side of all its extensions.  This proves convexity.  Put
$S=\ran(p)$ and $\delta=\dom(p)$.  By Lemma~\ref{swo:lem:delete} and
Proposition~\ref{swo:prop:universal}, fix an order-isomorphism
$h:\No\to\No\setminus S$.  The ordinal tail $[\delta,\Ord)$ has order
type $\Ord$, with its canonical increasing parametrization
$\xi\mapsto\delta+\xi$.  Given $E\in\Hsl$, prepend $p$ to the enumeration
$h\circ E$ on this tail.  The result is a bijection onto $\No$, since its
prefix covers $S$ and its tail covers exactly the complement.  Conversely,
restrict an element of $C_p$ to its tail and apply $h^{-1}$.  These
transformations are inverse and preserve first-difference comparisons.
The final assertion follows from Theorem~\ref{swo:gl:interpolation}.
\end{proof}

In particular, every such interval contains a copy of the entire comparison
structure $\Hsl$, not merely a copy of an arbitrarily large set order.  This
self-similarity does not invoke a class whose members are the cylinders;
for each set prefix the displayed construction gives the required
transformation, using the indicated class parameters.

% ---- source 11, lines 1787-1821
\begin{proposition}[Approximants at every point]\label{swo:gl:local}
For each $F\in\Jinj$ there are uniformly ordinal-indexed families
$(a_\alpha)$ and $(b_\alpha)$ of codes in $\Pbd$ such that, whenever
$\alpha<\beta$,
\[
 a_\alpha<_{\lex}a_\beta<_{\lex}F
 <_{\lex}b_\beta<_{\lex}b_\alpha.
\]
The lower family is cofinal among the global enumerations below $F$, and the
upper family is coinitial among those above $F$.  Neither side admits a
uniformly set-indexed cofinal or coinitial family, respectively.
\end{proposition}

\begin{proof}
Choose a fresh value below $F(\alpha)$, append it to
$F\upharpoonright\alpha$, and use Lemma~\ref{swo:gl:completion} to obtain
$a_\alpha$.  Choose a fresh value above $F(\alpha)$ to obtain $b_\alpha$.
Lemma~\ref{swo:lem:delete} supplies the fresh values, and the fixed global
well-order makes all choices uniform.  For $\alpha<\beta$, the comparison
between $a_\alpha$ and $a_\beta$ is decided at $\alpha$, giving the asserted
inequality; the other comparisons are identical arguments.

If $E<_{\lex}F$, choose $\alpha>d(E,F)$.  Then $a_\alpha$ agrees with $F$
at the decisive coordinate, so $E<_{\lex}a_\alpha<_{\lex}F$.  This proves
cofinality, and the upper assertion is symmetric.  Finally, given any
uniformly set-indexed family below $F$, Theorem~\ref{swo:gl:interpolation}
applied with upper side $\{F\}$ supplies a larger lower approximant in
$\Pbd$.  Hence the original family was not cofinal.  The dual argument
proves the coinitiality assertion.
\end{proof}

For $F\in\Hsl$, these are the local cofinality and coinitiality statements
at an actual point of the order.  It is appropriate to say that ordinal-length
class families suffice and set-sized families do not; assigning an ordinary
set cardinal to these proper-class cofinalities would obscure the distinction.

% ---- hand chunk n_local12

\begin{note}[Source~12, Proposition~6.4, ``Cofinality and local character'']\label{swo:st:prop:characters}
\emph{No set-indexed uniform family is cofinal or coinitial in $\Hsl(\No)$;
there are increasing cofinal and decreasing coinitial families indexed by
$\Ord$; at every point neither side has a set-indexed cofinal, respectively
coinitial, family, but both have $\Ord$-indexed ones.} This is
\cref{swo:gl:local} together with the remark that first entries of a
set-indexed family are bounded, by the same argument; the cofinal family is
$(E_\alpha)_{\alpha\in\Ord}$ for the positive ordinals $\alpha$.
\end{note}

% ---- source 11, lines 1823-1850
\begin{corollary}[Set saturation is not Dedekind completeness]
\label{swo:gl:no-suprema}
A nonempty uniformly set-indexed family in $\Hsl$ with no greatest member
has no supremum in $\Hsl$.  Dually, a family with no least member has no
infimum.  Both statements also hold in the ordinary class order $\Pbd$.
\end{corollary}

\begin{proof}
Suppose that $E$ is the supremum of a family $(A_i)$ with no greatest member.
Every $A_i$ is then strictly below $E$.  By
Theorem~\ref{swo:gl:interpolation}, there is a bounded-support enumeration $p$
with $A_i<p<E$ for every $i$.  This makes $p$ a smaller upper bound, a
contradiction.  The infimum argument reverses the inequalities, and the same
bounded-support interpolant proves the assertions internal to $\Pbd$.
\end{proof}

For example, take the enumerations $E_n$ from the transposition construction,
so that $E_n(0)=n$ for $n\in\mathbb N$.  They form an increasing sequence
of bounded-support codes.  The enumeration whose first value is the surreal
ordinal $\omega$ is an upper bound, but there is no least upper bound in
either $\Pbd$ or $\Hsl$.  The obstruction is not a failure to find bounds:
Theorem~\ref{swo:gl:interpolation}, with one input side empty, bounds every
uniformly set-indexed family above and below.  Rather, whenever a proposed
bound is strictly beyond every member, there remains room to move the bound
closer.  This distinction will also matter when the full order of Dedekind
cuts is embedded below.  An order embedding preserves comparisons, but does
not automatically preserve the role of a point as the least upper bound of
a specified family.

% ---- source 11, lines 1852-1894
\subsection{Binary class words, cuts, and diagonalization}

\begin{theorem}\label{swo:gl:binary}
The lexicographic orders $\Hsl$ and $\{0,1\}^{\Ord}$ are order-bi-embeddable,
where a binary word of length $\Ord$ means a class characteristic function.
Both embeddings have uniform relational definitions in $\GBC$.
\end{theorem}

\begin{proof}
Partition the coordinate ordinals into successive pairs
$2\cdot\xi,2\cdot\xi+1$.  Here the multiplication is ordinal multiplication
with $2$ on the left; every ordinal has a unique such pair position.  Sort
the corresponding baseline values as
\[
 u_\xi=\min\{e_0(2\cdot\xi),e_0(2\cdot\xi+1)\},\qquad
 v_\xi=\max\{e_0(2\cdot\xi),e_0(2\cdot\xi+1)\}.
\]
For a class $D\subseteq\Ord$, put these two values in the order
$(u_\xi,v_\xi)$ if $\xi\notin D$, and in the order $(v_\xi,u_\xi)$ if
$\xi\in D$.  Denote the resulting enumeration by $T_D$.  Each pair retains
exactly its original two values, so $T_D$ is a bijection onto $\No$.  At the
least differing bit of two characteristic functions, $0$ chooses the smaller
first value and $1$ the larger one.  Thus $D\mapsto T_D$ preserves and
reflects lexicographic comparison.

For the converse, fix an ordinal enumeration $(p_\xi)_{\xi\in\Ord}$ of
the set codes in $\Pbd$.  Associate with $E\in\Hsl$ the class
\[
 D_E=\{\xi\in\Ord:p_\xi<_{\lex}E\}.
\]
If $E<_{\lex}F$, then $D_E\subseteq D_F$, and the inclusion is strict:
Corollary~\ref{swo:gl:pbd-no} supplies a code strictly between $E$ and $F$.
At the least ordinal belonging to $D_F\setminus D_E$, the characteristic
functions therefore have values $0$ and $1$, respectively.  This proves the
second order embedding.  Both constructions use elementary class
comprehension with the input word or enumeration as a class parameter.
\end{proof}

Bi-embeddability must not be replaced by isomorphism.  Binary class words have
a least and a greatest word and have adjacent pairs: after a set prefix,
the word with next bit $0$ and all subsequent bits $1$ immediately precedes
the word with next bit $1$ and all subsequent bits $0$.  By contrast,
$\Hsl$ is dense and has no endpoints.

% ---- hand chunk binary12_head

\merge\ \emph{Second route (source~12).} Source~12's Section~7 embeds
$\Hsl(\No)$ into binary class words by concatenating self-delimiting codes of
the entries, instead of source~11's characteristic function of a dense class;
the other embedding is source~11's pair swap.

% ---- source 12, lines 721-771
\subsection{Source 12's route to the binary-class representation}
\label{swo:st:sec:binaryclass}

Let $\BOrd$ denote the predicate of class subsets of $\Ord$, regarded as
binary functions on $\Ord$ and compared at their least disagreement.
This is again notation for a predicate and relation on classes, not a
powerclass object.

\begin{lemma}[A self-delimiting numerical code]\label[lemma]{swo:st:lem:code}
Replace every minus in a sign sequence by $00$, every plus by $11$, and
append the terminal block $01$. The resulting binary ordinal word $c(s)$
has length $2\cdot \bd(s)+2$. Distinct codewords are prefix-incomparable,
and their first-difference comparison agrees with numerical comparison
of the surreals.
\end{lemma}
\begin{proof}
The three two-bit blocks have order $00<01<11$, exactly matching
$-<\End<+$. At a sign disagreement the first different block decides.
When one sign sequence ends, its terminal block differs from the next
sign block of the other. Hence neither complete codeword can be a prefix
of another, and the numerical comparison is preserved.
\end{proof}

\begin{theorem}[Binary-class equimorphism]\label{swo:st:thm:binaryclass}
In the sense of uniform class transformations,
\[
 \Hsl(\No)\equim\BOrd.
\]
\end{theorem}
\begin{proof}
For $E:\Ord\to\No$, concatenate $c(E(\alpha))$ over all ordinals.
The cumulative length before each stage is a set ordinal. All blocks
are nonempty, so the cumulative lengths are unbounded in $\Ord$ and the
result is a binary class on all of $\Ord$. If $E,F$ first disagree at
$\alpha$, their earlier block lengths and bits agree, and the code of
that entry decides the same inequality by \cref{swo:st:lem:code}.

For the converse, use the disjoint pairs $a_\alpha<b_\alpha$ in the proof
of \cref{swo:st:thm:diagonal}. For a binary class $A$, list the pair in ascending
order if its bit is $0$, and descending order if its bit is $1$.
Put pair blocks in ordinal order. The first different bit gives the
first different pair orientation and the correct lexical inequality.
All labels occur exactly once.
\end{proof}

This is not an isomorphism assertion. The binary-class order has least
and greatest elements and adjacent pairs, such as a zero followed by all
ones versus a one followed by all zeros. By
\cref{swo:st:thm:etasl}, $\Hsl(\No)$ has no endpoints or adjacent pairs.
The distinction between equimorphism and isomorphism is already visible
at this level.

% ---- source 11, lines 1896-1936
\begin{corollary}\label{swo:gl:cuts}
The higher-order order of all class initial segments of $\No$, ordered by
inclusion, embeds into $\Hsl$.  In particular, this holds for any specified
convention of Dedekind cuts.  Conversely, each $E\in\Hsl$ determines the
initial segment $\{p\in\Pbd:p<_{\lex}E\}$, and these initial segments
separate distinct enumerations.
\end{corollary}

\begin{proof}
For a class initial segment $C\subseteq\No$, form
$e_0^{-1}[C]\subseteq\Ord$ and apply the pair-swap embedding of
Theorem~\ref{swo:gl:binary}.  Proper inclusion of initial segments gives proper
inclusion of their inverse images, which implies the correct lex comparison
of their characteristic functions.  The converse is the dense-class
construction in the second half of that theorem.
\end{proof}

This corollary does not form a GB class of all Dedekind cuts of $\No$.
Each individual cut is an allowed class parameter; the rule associates an
enumeration to it and compares any two such results.  An encompassing object
belongs to a higher type or to an explicitly larger ambient universe.  Nor
does the corollary assert that the natural inclusion of $\Pbd$ in $\Hsl$
is a Dedekind completion: an order embedding of cuts need not preserve their
intended suprema or their original positions relative to the embedded line.

There is a useful description of the natural cuts produced by an enumeration
outside the bounded-support family.  If $E\in\Hsl$ has no code in $\Pbd$,
then
\[
 L_E=\{p\in\Pbd:p<_{\lex}E\},\qquad
 U_E=\{p\in\Pbd:E<_{\lex}p\}
\]
partition $\Pbd$ into two nonempty classes.  Proposition~\ref{swo:gl:local}
shows that $L_E$ has no set-sized cofinal subset and $U_E$ has no set-sized
coinitial subset, while ordinal-indexed class families suffice on both
sides.  This cut is unfilled in $\Pbd$: a putative filling code and $E$
would be separated by another bounded-support code.  Transporting the cut
through an isomorphism $\Pbd\cong\No$ therefore exhibits a surreal class
cut with neither side controlled by a cofinal or coinitial set.  This is a
description of individual additional enumerations relative to the fixed
dense class, rather than a classification of all surreal class cuts.

% ---- source 11, lines 1938-1965
\begin{theorem}[No exhaustive uniform class of codes]\label{swo:gl:diagonal}
There is no ordinary class of set codes with a single class evaluation
relation whose evaluated rows include every member of $\Hsl$.
\end{theorem}

\begin{proof}
If such a coding existed, global choice would enumerate its codes by
ordinals, with repetitions if the class of codes were a set.  Write the
resulting row relation as $\operatorname{Eval}(\xi,\alpha,x)$.
Use the sorted pairs $u_\xi<v_\xi$ from Theorem~\ref{swo:gl:binary},
and define the class
\[
 D=\{\xi:\operatorname{Eval}(\xi,2\cdot\xi,u_\xi)\}.
\]
The associated pair-swap enumeration $T_D$ differs from row $\xi$
on the pair $(2\cdot\xi,u_\xi)$.  If that pair belongs to the old row,
it does not belong to the graph of $T_D$, whose value there is
$v_\xi$.  If it does not belong to the old row, it belongs to the
new graph, whose value there is $u_\xi$.
Thus $T_D$ is a member of $\Hsl$ omitted by the proposed exhaustive list.
The argument even permits invalid or nonfunctional rows in the evaluator.
\end{proof}

The theorem explains why class universality of $\No$ cannot simply be applied
to the full collection $\Hsl$.  It also pinpoints the limitation of coding
only bounded modifications, or any other one uniformly evaluable family of
set codes: such a family may be order-dense and may realize all set-indexed
cuts, while still omitting global enumerations by diagonalization.

% ---- hand chunk diag12_head

\merge\ \emph{Strengthening (source~12).} Source~12's Section~5 proves that no
uniform family of class well-orders indexed by a class of sets contains every
class well-order of $\No$, and that the missing one can be chosen set-like;
its \cref{swo:st:n5.2} is \cref{swo:gl:diagonal}. Its subsection on definable
classes explains why enumerating formulas does not evade the theorem.

% ---- source 12, lines 496-558
\subsection{Source 12's strengthening: why the collection of all global well-orders is not a class}
\label{swo:st:sec:diagonal}

The formal obstacle is not that individual class relations are illegal.
They are legitimate class variables. The obstacle is attempting to make
them the elements of an ordinary class of the same theory.

A useful legal replacement is a \emph{uniform family}. For example,
$F\subseteq I\times\Ord\times\No$ is a class of triples of sets, and its
section $F_i$ can be a class bijection for every $i\in I$. Similarly a
class of triples $(i,x,y)$ can code a family of class order relations.
The index class $I$ may be proper. These are single classes, not functions
whose values are class objects.

\begin{theorem}[No exhaustive class-coded family]\label{swo:st:thm:diagonal}
In GBC, no uniform family of class well-orders of $\No$ indexed by a
class of sets contains every class well-order of $\No$. In fact, for
every such family there is a \emph{set-like} class well-order not among
its sections.
\end{theorem}
\begin{proof}
Global choice permits the family to be reindexed by an ordinal or by
$\Ord$. Extend an ordinal-indexed family to an $\Ord$-indexed one by
repeating an arbitrary default well-order; an empty family is trivial.
Write the resulting relations as $R_\alpha$.

Fix a bijection $B:\Ord\to\No$ and partition its range into the pairs
\[
 \{B(2\cdot\alpha),B(2\cdot\alpha+1)\}
   =\{a_\alpha,b_\alpha\},\qquad a_\alpha<b_\alpha.
\]
Here multiplication is ordinal multiplication. The blocks of length two
partition $\Ord$. Define a new enumeration $E$ by placing the two labels
of pair $\alpha$ in the order \emph{opposite} to their relative order in
$R_\alpha$. Keep the pair blocks in increasing $\alpha$ order.

This is a class bijection $\Ord\to\No$, defined elementarily from the
family and $B$. Its induced well-order disagrees with $R_\alpha$ on the
relative order of $a_\alpha,b_\alpha$. Thus it is not any section of the
family.
\end{proof}

\begin{corollary}[No uniform set-code solution]\label{swo:st:n5.2}
There is no class $D$ of set codes together with a uniform class decoder
whose decoded relations are exactly all class well-orders of $\No$.
\end{corollary}

This diagonal theorem applies equally to an attempted exhaustive coding
of all class subsets of $\Ord$ through the pair construction. It does not
say that an externally specified set model has too many classes to form
an external set. It says that the corresponding exhaustive decoder is
not available \emph{internally at the same size level}.

\subsubsection{Definable classes and the truth-predicate trap}

Formula codes and set parameters are sets, but the semantic map
``formula plus parameters $\mapsto$ its extension in the universe'' is
not a uniformly available truth predicate for the universe. For each
fixed formula one may define a class. One may also work externally with
all definable classes of a chosen model. Neither maneuver supplies an
internal universal decoder defeating \cref{swo:st:thm:diagonal}.
The same caution applies to quantifying over all possible Lean universe
levels: universe polymorphism is not an internal type of all universes.

% ---- source 11, lines 1967-2004
\subsection{Omitted ranges and the transition to raw well-orders}

\begin{theorem}[An omitted-range cut]\label{swo:gl:omitted}
Fix $r\in\No$ and a class bijection
$F:\Ord\to\No\setminus\{r\}$.  There are increasing and decreasing
ordinal-indexed families $(a_\alpha)$ and $(b_\alpha)$ in $\Pbd$ such that
\[
 a_\alpha<_{\lex}b_\beta\qquad(\alpha,\beta\in\Ord),
\]
but no member of $\Hsl$ lies strictly between both entire families.
The raw well-order that lists $F$ and then $r$ has type $\Ord+1$ and is the
unique member of $\Hall$ between these families.
\end{theorem}

\begin{proof}
Use the approximants from Proposition~\ref{swo:gl:local}.  They satisfy
$a_\alpha<_{\lex}F<_{\lex}b_\beta$ for every pair of indices.  If a global
bijection $E$ lay between both families and differed from $F$ first at
$\delta$, an approximant with index greater than $\delta$ would place $E$
on the wrong side.  Therefore $E$ would equal $F$ at every ordinal, which
contradicts its surjectivity.

Let $R$ be the raw order obtained by appending $r$ after the enumeration $F$.
Its set-like initial part is precisely $F$: every point listed by $F$ has
set-sized predecessors, whereas $r$ has a proper class of predecessors.
Its comparisons with the approximants are decided at their prescribed
ordinal first differences, so it lies between both families.  Conversely,
the same first-difference argument forces the core $F_S$ of any raw order
$S$ between the families to equal $F$.  By the core decomposition of
Theorem~\ref{swo:cf:core}, its residual order must be the unique order on the
singleton $\{r\}$.  Thus $S=R$.
\end{proof}

The lower and upper families here are indexed by all ordinals, not by sets.
Accordingly, the missing interpolant does not contradict
Theorem~\ref{swo:gl:interpolation}.  It gives a concrete boundary between the
set-cut saturation of the set-like comparison and the additional points
provided by non-set-like global well-orders.

% ---- source 11, lines 2006-2050
\begin{example}[Two omitted points produce adjacent raw orders]
\label{swo:gl:omitted-two}
Let $r<s$ and choose $F:\Ord\to\No\setminus\{r,s\}$ bijectively.  The
approximants to $F$ have exactly two raw interpolants: append $r,s$ in that
order, or append $s,r$.  Both raw orders have type $\Ord+2$, and the first
is lexicographically smaller.  They are adjacent in $\Hall$.  Indeed, a raw
order between them cannot have a different core, since an ordinal first
difference would put it on the same side of both.  With core $F$ fixed, its
residual order is a well-order of a two-element set, and these are exactly
the two displayed possibilities.  Thus the density theorem for $\Hsl$
does not extend to all raw well-orders.
\end{example}

\begin{corollary}[Exactly which raw set cuts admit bounded completions]
\label{swo:gl:raw-cut}
Let $(A_i)_{i\in I}<_{\lex}(B_j)_{j\in J}$ be a uniformly set-indexed
strict cut in $\Hall$.  A member of $\Pbd$ lies between the two sides if
and only if
\[
 F_{A_i}\ne F_{B_j}\qquad(i\in I,\ j\in J),
\]
where $F_R$ is the canonical ordinal enumeration of the set-like initial
part of a raw well-order $R$.  The condition is vacuous when a side is empty.
\end{corollary}

\begin{proof}
If the cross cores are different, their comparisons are decided at ordinal
first differences and have the same direction as the given raw comparisons.
Apply Theorem~\ref{swo:gl:interpolation} to these uniform families in $\Jinj$.
Every comparison of the resulting completion with a raw input is already
decided within its core, giving the required raw inequalities.

Conversely, suppose a cross pair shares a core $F$.  This core cannot be
surjective: a surjective core determines the entire raw order, contradicting
the strict comparison of the pair.  Any $E\in\Hsl$ therefore differs from
$F$ at an ordinal coordinate.  That first difference places $E$ on the same
side of both raw orders, so $E$ cannot lie between them.  This excludes a
bounded-support completion in particular.
\end{proof}

The criterion isolates the obstruction without appealing merely to size.
Distinct ordinal cores can be separated by one bounded change of the
baseline, even for a whole set-indexed cut.  Two opposite sides with the same
core require information in the residual non-set-like order, beyond every
ordinal coordinate available to a global set-like enumeration.

% ---- hand chunk part5
\part{Foundations, models and formalization}\label{swo:part:models}

% ---- source 11, lines 2055-2383
\section{Formalizing the subject without changing its meaning}
\label{swo:sec:formalization}

\subsection{A bounded development in ordinary set theory}

For any ordinal \(\theta\), use the set \(X=\No_{<\theta}\).
The well-order relations on \(X\) form a subset of
\(\mathcal P(X\times X)\). Their lexicographic comparison can be
defined either by unique ordinal enumeration or by the predecessor
formula. This is ordinary \(\ZFC\) mathematics. No inaccessible
cardinal, class theory, or new logical sort is required.

This approach is particularly good for proving the general alphabet
theorem, investigating cardinal spectra, and checking how a proposed
construction changes birthdays. One should distinguish \(\WO(X)\),
which permits every possible ordinal type of cardinality \(\card X\),
from its minimal-length stratum. Their cardinalities agree, but
their jumps and finer order structure need not.

A family of results for every \(\theta\) does not automatically give
a coherent global map. To take a union, the maps must agree on overlaps.
The repository's bounded-surreal nonembedding example makes this
restriction concrete.

\subsection{Second-order class theory}

In \(\GB\) or \(\GBC\), the formula
\[
 \operatorname{ClassWO}(R)\ \wedge\ \operatorname{ClassWO}(S)
 \ \Longrightarrow\
 R=S\ \vee\ \operatorname{Lex}(R,S)\ \vee\ \operatorname{Lex}(S,R)
\]
is meaningful with class variables. One can prove trichotomy,
transitivity, existence theorems, and uniform-family embedding
theorems in this language. It is unnecessary to form a class
whose members are \(R,S\).

When a family is indexed by a class \(I\) of sets, one class
relation \(T\subseteq I\times\No\times\No\) supplies its rows.
If the rows are pairwise distinct, substituting them into the
elementary predecessor formula produces a class linear order on
\(I\). With repeated rows, one instead has the corresponding
comparison on codes modulo equal evaluations. This is the correct format for
most class-sized examples and applications.

In contrast, an expression such as
\[
 \{R:R\text{ is a class well-order of }\No\}
\]
does not define an ordinary \(\GB\) class of members. Our fraktur
notation abbreviates a domain of class quantification. The
diagonal theorem shows why a single uniform set-code replacement
cannot recover the full domain either.

\subsection{Why \texorpdfstring{\(\KM\)}{KM} does not erase the extra level}

\(\KM\) strengthens comprehension by allowing class quantifiers,
but its classes still have \emph{sets} as members. It therefore
does not turn the totality of class well-orders into an ordinary
class. A genuine object containing arbitrary classes as members
requires another sort, or a containing universe in which those
classes become sets.

This distinction matters when one wants an arbitrary topology
of higher-order open collections, a measure on all global
well-orders, or a power object of their totality. Such notions
require an explicit ambient theory. Hyperclass forcing provides
one example of higher-level foundational practice
\cite{AntosFriedman}; it is not a reason to silently ignore the
level distinction in bare \(\KM\).

There is also a separate recursion issue. Extracting the
canonical initial \(\Ord\)-part of one raw well-order is an
elementary class construction. Iteratively deleting proper-class
tails through an arbitrary transfinite sequence would have
class-valued stages and needs its own analysis. General
comparability of abstract class well-orders by initial-segment
isomorphisms is another stronger topic; \(\ETR\) implies it
\cite{HamkinsWoodin}. None of that machinery is needed for
the labeled comparison proved in this article.

\subsection{An explicit containing-universe model}

A useful concrete model assumes a strongly inaccessible cardinal
\(\kappa\) in an ambient universe and takes
\[
 U=V_\kappa,\qquad \mathcal C=\mathcal P(V_\kappa).
\]
This is a sufficient setting for interpreting a full class
universe. The internal surreal class becomes the ambient set
\[
 \No^U=\No_{<\kappa},\qquad \card{\No^U}=\kappa.
\]
All its ambient well-orders form a set of cardinality \(2^\kappa\).
Those that are internally set-like have ambient order type
\(\kappa\); unrestricted well-orders may have any ambient
ordinal type of cardinality \(\kappa\).

To see the set-like assertion, every initial segment of a
\(\kappa\)-enumeration has fewer than \(\kappa\) members, with
ranks bounded below \(\kappa\); regularity puts it in \(U\).
Conversely, an ordinal type strictly exceeding \(\kappa\) has
an initial segment of cardinality \(\kappa\), which cannot be
a set of \(U\). A smaller type cannot enumerate the entire
carrier.

The bounded-support class becomes an ambient order of cardinality
\(\kappa\), while the full \(\kappa\)-permutation order has
cardinality \(2^\kappa\). Thus a \(\kappa\)-sized dense suborder
can coexist with a larger order of all enumerations. The exact
ambient ordinal cutoff remains \(\kappa^+\).

Every ordinal \emph{of \(U\)} is below \(\kappa\). There is no
contradiction between internally embedding all ordinals and
having an ambient successor-cardinal barrier. Inaccessibility
is an additional assumption for this full universe example;
it is not required for an individual birthday fragment.

\subsection{A strict nonembedding result}

The containing-universe interpretation makes a further comparison
decidable by ordinary set-order invariants.

\begin{theorem}[Nonembedding in a full universe]
\label{swo:thm:universe-nonembedding}
In the preceding full interpretation
\(U=V_\kappa\), \(\mathcal C=\mathcal P(V_\kappa)\), with ambient
\(\ZFC\) and \(\kappa\) strongly inaccessible,
\[
 \Hsl(\No^U)\emb\Hall(\No^U),
 \qquad
 \Hall(\No^U)\not\emb\Hsl(\No^U).
\]
These are assertions about the actual ambient sets of internal
class well-orders and arbitrary ambient order embeddings.
\end{theorem}

\begin{proof}
The first embedding is inclusion. Put \(X=\No^U\).
The target of the proposed reverse embedding is the dense order
of all bijections \(\kappa\to X\). Its bounded-support suborder
\(D=\Pbd^U\) is dense and has ambient cardinality \(\kappa\).
For the upper bound, codes with domain \(\lambda<\kappa\) are
subsets of \(\lambda\times\lambda\); strong inaccessibility bounds
their number below \(\kappa\), and there are \(\kappa\) possible
domains. Transpositions give the matching lower bound.

Fix labels \(a<b\) in \(X\). There are \(2^\kappa\) bijections
\(F:\kappa\to X\setminus\{a,b\}\), by the fixed-alphabet
cardinality theorem. Each determines a convex adjacent pair
\[
 R_F=F+(a,b)<_{\lex}S_F=F+(b,a)
\]
in the unrestricted order. Distinct \(F\)'s determine distinct,
disjoint convex fibers; one entire pair precedes the other.

If \(\Phi:\Hall(X)\emb\Hsl(X)\) existed, the intervals
\[
       \bigl(\Phi(R_F),\Phi(S_F)\bigr)
\]
would be \(2^\kappa\) pairwise disjoint nonempty open intervals
in the target. Each contains a point of \(D\). Choosing the
least such point under a fixed well-order of \(D\) would inject
\(2^\kappa\) into \(\kappa\), a contradiction.
\end{proof}

Thus equal cardinality \(2^\kappa\) and equal ordinal embedding
cutoff \(\kappa^+\) do not imply equimorphism. The useful
invariant here is the size of a family of disjoint open
intervals relative to a dense suborder. The repository already
uses related separability and cellularity arguments for the
real alphabet \cite[Remark
\texttt{lwo:sp:rem:conditional}]{ProveItReals}; the argument
above supplies the precise surreal-universe version.

There is also a direct class-theoretic obstruction, provided the
inverse evaluation relation can be formed by impredicative
comprehension.

\begin{theorem}[No definable uniform embedding in
  \(\KM+\GC\)]\label{swo:thm:km-nonembedding}
For every fixed formula, with fixed set and class parameters,
that defines a total transformation \(R\mapsto\Phi(R)\) from
\(\Hall(\No)\) to \(\Hsl(\No)\), the theory \(\KM+\GC\)
proves that this transformation is not an order embedding.
The transformation is specified relationally: the defining
formula assigns a unique output class graph to each input
class well-order. This is a schema for fixed formulas, not
quantification over a third sort of arbitrary transformations.
\end{theorem}

\begin{proof}
Suppose such an order embedding were defined. Fix \(a<b\)
and a global well-order of the set codes in \(\Pbd\).
For each class bijection
\(F:\Ord\to\No\setminus\{a,b\}\), form the adjacent raw pair
\(R_F,S_F\) from the preceding proof. Density of \(\Pbd\)
in \(\Hsl\) supplies codes strictly between
\(\Phi(R_F)\) and \(\Phi(S_F)\). Let \(q_F\) be the
globally least such code. Comprehension in \(\KM\) makes
this a legitimate least-element definition, even if the
formula specifying \(\Phi\) contains class quantifiers.
Disjointness of the image intervals implies
\[
                 q_F=q_G\quad\Longrightarrow\quad F=G.
\]

Now use \(\KM\) comprehension to form the ordinary class
relation
\[
 \operatorname{Eval}(q,\alpha,x)
 \ \Longleftrightarrow\
 \exists\hbox{ class }F\,
 \bigl[
 F:\Ord\overset{\sim}{\longrightarrow}\No\setminus\{a,b\}
 \ \wedge\ q=q_F\ \wedge\ F(\alpha)=x
 \bigr].
\]
Its rows are empty or the graphs of unique such bijections,
and every such bijection occurs. Compose all rows with one
fixed class bijection
\(\No\setminus\{a,b\}\to\No\). The resulting single
class evaluator, with the ordinary set codes \(q\in\Pbd\),
contains every global enumeration of \(\No\).
This contradicts Theorem~\ref{swo:gl:diagonal}.
\end{proof}

The inverse evaluator is the decisive extra comprehension
step: its definition quantifies over an arbitrary class
\(F\). The preceding proof does not claim that elementary
\(\GB\) comprehension suffices.
Nor does it exclude arbitrary external maps between the
available-class collections of a nonfull model. In particular,
for an externally countable model these collections can both
be countable, and its externally dense set-like order can
admit an external copy of its unrestricted order. Such a map
need not be representable by any uniform definition in the
model. This is why the stated interpretation of an embedding
matters.

\subsection{Internal and external well-foundedness}

For a model \((M,\mathcal X)\), an internally well-founded
class relation need not be externally well-founded. An external
descending sequence might not be represented by a set in \(M\).
Thus comparing the internal theory with an external collection
of all relations requires care.

Over \(\GBC\), well-foundedness of a class relation can be
characterized by the absence of a set-coded descending
\(\omega\)-sequence \cite[Section 3]{HamkinsWoodin}. Keeping
the sets and an existing relation fixed while changing only
the available classes does not arbitrarily change that
relation's internal well-foundedness in suitable models.
What can change is which class relations are available at all.
External correctness can be ensured by an appropriate
well-founded full universe interpretation, as above.

\subsection{A practical Lean architecture}

The repository already has a universe-indexed sign carrier
\cite{ProveItSigns}. Its type structure can be displayed schematically as:
\begin{lstlisting}
structure SignSequence : Type (u + 1) where
  birthday : Ordinal.{u}
  signAt : Ordinal.{u} -> SignType
  signAt_ne_zero_iff : forall i, signAt i != 0 <-> i < birthday
\end{lstlisting}
This is mathematical pseudocode using ASCII logical notation,
not a new executable Lean declaration. The key point
is the level separation: the full carrier is one universe
above the types indexing its small sign lengths.

The following existing declarations were inspected:
\begin{center}
\small
\begin{tabular}{@{}>{\raggedright\arraybackslash}p{0.43\textwidth}>{\raggedright\arraybackslash}p{0.49\textwidth}@{}}
\toprule
Module or declaration & Role in the proposed development\\
\midrule
\code{SignSequence.ordinalOrderEmbedding} &
The all-plus ordinal embedding\\[3pt]
\code{birthdays\_bounded}, \code{small\_bounded} &
Control of small parameter families and bounded fragments\\[3pt]
\code{SignSequence.not\_small} &
The carrier cannot be treated as small at its ordinal level\\[3pt]
\code{SignSequenceCutOperation.cut} and \code{cut\_realizes} &
Actual separators for small cut data\\[3pt]
\code{SizeObstructions.noUnrestrictedCuts} &
The obstruction to unrestricted cut filling\\
\bottomrule
\end{tabular}
\end{center}

These are source-level observations, not results of a fresh
build. The current repository contains later arithmetic
modules as well; historical proposal notes should not be
mistaken for the present implementation status.

A sensible formalization proceeds in four parts.

\begin{enumerate}
\item \textbf{Set carriers.} For \(X:\code{Type u}\), represent
well-orders as relations on \(X\) together with their laws.
Prove the predecessor-based comparison and its equivalence
with ordinal enumeration. This avoids transporting across
ordinal representatives at every comparison.
\item \textbf{Ordinal words.} Define words with small
ordinal domains, an injective subtype, ordinal concatenation,
and the sign code. Formalize prefix-freeness and first
disagreement before proving the cardinal spectrum.
\item \textbf{Relative set-likeness.} If \(X\) lives in
\(\code{Type (u+1)}\), express the condition that each
predecessor subtype is \(\code{Small.\{u\}}\). A well-order
of this larger \(X\) has its unrestricted order type at
the larger ordinal universe. A same-level enumeration must
not be inferred without the relative smallness hypothesis.
\item \textbf{Uniform families.} Keep an explicit index type
and an evaluation relation. The bounded-support codes are
sets of ordinal pairs at the selected level. Formulate the
interpolation theorem only for small index families and
the no-code theorem at its correct larger level.
\end{enumerate}

Lean universes and Mathlib ordinals provide precisely the
needed separation \cite{LeanUniverses,MathlibOrdinal}. Universe
polymorphism is a schema of well-typed constructions, not a
term denoting the union of all universes. This article supplies
no new Lean file and makes no machine-verification claim.

% ---- hand chunk model12_head

\merge\ Source~12's Section~11 identifies both global orders in the full-class
model $(V_\kappa,\mathcal P(V_\kappa))$ of \cref{swo:thm:universe-nonembedding}.
Its clause~(iv) gives a second proof of that theorem's non-embedding: the full
order $\WO(\No_{<\kappa})$ has binary coding length $\kappa^+$
(\cref{swo:st:thm:setspectrum}), whereas the set-like order
$\WO_\kappa(\No_{<\kappa})$ has coding length $\kappa$
(\cref{swo:st:thm:transition}, regular case), and coding length cannot
decrease under embeddings (\cref{swo:st:lem:strictcube}). Source~11's proof
counts $2^\kappa$ disjoint intervals against a dense suborder of size
$\kappa$.

% ---- source 12, lines 1207-1297
\section{A full-class model that connects the two size levels (source 12)}
\label{swo:st:sec:model}

Assume in an ambient ZFC universe that $\kappa$ is strongly inaccessible.
Consider the two-sorted structure
\[
 \mathcal M=(V_\kappa,\mathcal P(V_\kappa)),
\]
with all external subclasses of $V_\kappa$ available as its classes.
This is a familiar full-class interpretation of strong class set theory:
regularity and the strong-limit property ensure the required closure of
set operations and class Replacement. It satisfies KM and global choice.
We use this as an explicit relative model, not as a claim that inaccessible
cardinals are necessary for the earlier GBC theorems.

\begin{theorem}[Exact external identification]\label{swo:st:thm:model}
For the full-class interpretation above:
\begin{enumerate}[label=(\roman*)]
\item the internal surreal carrier is externally $\No_{<\kappa}$, and its
internal ordinals are the external ordinals below $\kappa$;
\item the internal set-like global surreal well-orders, viewed as an
external set order, are exactly $\WO_\kappa(\No_{<\kappa})$;
\item all internal global surreal well-orders, viewed externally, form
exactly $\WO(\No_{<\kappa})$;
\item both external orders have cardinality $2^\kappa$ and ordinal
spectrum $\beta<\kappa^+$, but their binary coding ranks are respectively
$\kappa$ and $\kappa^+$. Therefore the full order does not embed into the
set-like order.
\end{enumerate}
\end{theorem}
\begin{proof}
A sign sequence belongs to $V_\kappa$ precisely when its ordinal domain
has length less than $\kappa$. Inaccessibility gives
$|\No_{<\kappa}|=\kappa$.

Every external relation on $\No_{<\kappa}$ is a permitted internal class:
each ordered pair of its labels is an element of $V_\kappa$. External
well-foundedness implies internal well-foundedness. Conversely, if a
relation is externally ill-founded, ambient choice supplies a descending
$\omega$-sequence. Its countably many labels have ranks bounded below
$\kappa$ by regularity, so the sequence is an element of $V_\kappa$ and
witnesses internal ill-foundedness. Thus the well-orders agree.

An internal set has external cardinality less than $\kappa$.
Conversely, a subset of $\No_{<\kappa}$ of cardinality less than $\kappa$ has
bounded ranks and is an element of $V_\kappa$. Therefore internal
set-likeness means that every proper predecessor segment has cardinality
less than $\kappa$. An external well-order of $\kappa$ labels has that
property exactly when its order type is $\kappa$, proving (ii).
Arbitrary external well-orders on those labels have types exactly the
ordinals of cardinality $\kappa$, proving (iii). Comparisons agree by
\cref{swo:st:thm:fulllex}. Finally apply
\cref{swo:st:thm:setspectrum,swo:st:thm:transition} and the monotonicity of coding rank
under embeddings.
\end{proof}

\subsection{Why this does not contradict surreal universality}

Internally, each class linear order has elements that are sets of
$V_\kappa$ and hence externally at most $\kappa$ elements.
The external collection of all class well-orders has $2^\kappa$
elements and is not an internal class carrier. Thus it is outside the
scope of the internal class-universality theorem.

If $\lambda>\kappa$ is another inaccessible cardinal, the entire
external set $\WO(\No_{<\kappa})$ can be studied as an ordinary set in
$V_\lambda$. It then embeds into the larger surreal carrier there by
set-order universality. Its labels have moved up one foundational level;
no same-level universal object has been created.

The singular cutoff $\beth_\omega$ in \cref{swo:st:sec:transition} must not be
used as though $V_{\beth_\omega}$ were this model. That rank lacks the
regularity used in Replacement and in the well-foundedness argument.
The singular theorem is an ordinary ZFC theorem about a bounded set,
not a full-class model construction at a singular height.

\subsection{Henkin classes, countable models, and relativity}

If the class collection is some $\mathcal X\subsetneq\mathcal P(M)$,
only those well-orders whose relations lie in $\mathcal X$ are admitted.
External cardinality and saturation can then differ from the full-class
calculation. A countable model may have an externally countable collection
of admitted classes and external countable cofinal families, while
internally having no \emph{coded set-indexed} cofinal family. The missing
external families are simply not internal objects.

For nonstandard or non-$\beta$ class models, one must additionally
separate internal and external well-foundedness. Our exact identification
in \cref{swo:st:thm:model} used transitivity, full classes, and regularity; it
was not asserted for arbitrary models. The general subtleties of class
well-foundedness are discussed in \cite[\S3]{HamkinsWoodin}.

% ---- hand chunk model12_after

\merge\ Clause~(ii) of \cref{swo:st:thm:model} is source~11's statement that
the internally set-like well-orders have ambient type $\kappa$, and
clause~(iii) its statement that unrestricted ones have every ambient type of
cardinality $\kappa$ (\cref{swo:sec:formalization}). Source~12's remark that
the singular cutoff $\beth_\omega$ must not be used as though
$V_{\beth_\omega}$ were such a model is the counterpart of source~11's remark
that inaccessibility is needed for the full universe example but not for an
individual birthday fragment.

% ---- hand chunk formal12_head@1

\merge\ The next section is source~12's formalization proposal. The modules it
names are proposals and do not exist; the existing declarations it cites exist
at the current revision, and none states a result of this report
(\cref{swo:sec:formal}).

% ---- source 12, lines 1336-1445
\section{A formalization architecture with explicit size boundaries (source 12)}
\label{swo:st:sec:formalization}

\subsection{Start with a universe-polymorphic set theorem}

The most efficient verified starting point is independent of surreal
arithmetic. For a type $X$ with a linear order, represent a well-order on
$X$ by its binary relation and a proof of well-foundedness and linearity.
Define the common prefix by equality of predecessor sets and equality of
the restricted relations, exactly as in \cref{swo:st:sec:commonprefix}.
Prove maximality, distinct next labels, and transitivity of comparison.
This gives a generic theorem reusable for reals, surreal cutoffs, and
other ordered alphabets.

For a fixed ordinal domain, represent an enumeration as an equivalence
between an ordinal's carrier and $X$. For all well-orders of a small $X$,
using relations on $X$ often avoids unnecessarily quantifying over every
ordinal in a larger universe. The ordinal enumeration can be recovered
later from the well-order type. This is a representation choice, not an
identification of all ordinals with a small type.

\subsection{Use the existing sign-sequence foundations selectively}

The inspected repository file
\begin{quote}\small
\path{Algebra/SurrealNumbers/Surreal/Foundations/SignSequenceSimplicity.lean}
\end{quote}
contains, among others,
\begin{quote}\small\ttfamily
isPrefix\_of\_minimum\_birthday\par
existsUnique\_prefix\_of\_ordConnected\par
existsUnique\_minimum\_birthday\_of\_ordConnected\par
existsUnique\_simplest\_separator
\end{quote}
\cite{ProveItSimplicity}. Its comments explicitly distinguish uniqueness of a
simplest separator from existence of separators for arbitrary small cuts.
The last theorem assumes a separator exists. It must not be imported as
though it independently established all cut-filling instances.

A suitable interface for the order-theoretic development needs a birthday
map, a numerical linear order, sign concatenation and truncation, the
first-difference comparison, and small-cut existence with a birthday
bound. Field operations can be omitted except in illustrative examples.

\subsection{Suggested proof modules}

\begin{longtable}{@{}p{0.27\textwidth}p{0.65\textwidth}@{}}
\toprule
Module & Principal obligations\\
\midrule
\endhead
\texttt{CommonPrefix} & The elementary definition; initiality; maximality;
next-label inequality; generic full lexical linearity.\\
\texttt{ExhaustiveLex} & Prefix-freeness; equivalence with ordinal
first-difference comparison; transport between relations and enumerations.\\
\texttt{BinaryCubes} & Cylinder fusion; strict hierarchy;
characteristic-cut embeddings; precise ordinal concatenation.\\
\texttt{PermutationCoding} & Pair orientation, product scheduler,
transpositions, and exhaustive output.\\
\texttt{BirthdayCutoffs} & Size of $\No_{<\kappa}$; prefix-free numerical code;
cone partition; regular and singular sums.\\
\texttt{ClassFamilies} & A formal set/class model, sections of class
relations, baseline enumeration, and diagonal non-enumerability.\\
\texttt{UniverseModels} & A chosen inaccessible-rank interpretation;
internal/external well-foundedness and identification of the two orders.\\
\bottomrule
\end{longtable}

The finite common-prefix theorem should be proved directly for relations,
not by first developing all of ordinal arithmetic. The scheduler proof
should record both surjectivity and the first-row-appearance invariant;
a bare cardinal pairing equivalence is insufficient. The singular theorem
must use ordinal sums, never a cardinal cast that silently loses order.

\subsection{Universe accounting}

Schematically, a sign-sequence type built from small ordinal lengths has
one universe level for its ordinal index and potentially a higher level
for the total carrier. One should retain explicit levels rather than
promise a concrete declaration before checking the chosen ordinal API.
A bounded carrier with a specified small birthday bound is the cleanest
first implementation.

Three operations must remain visibly different: taking a subtype of one
fixed surreal carrier; lifting an existing carrier to a larger universe;
and constructing all surreals of the larger universe. The third adds new
birthdays and is not merely a type lift. Likewise, a type of relations on
a fixed carrier can exist at a higher level without becoming a same-level
class of all class relations.

For global theorems, two safe routes are available. One may formalize GB
or GBC models and express theorems about their class predicates. Or one
may formalize the bounded analogue inside a hierarchy of explicit
universes and separately prove its interpretation in a set/class model.
Neither route requires an inconsistent type of all types.

\subsection{Axioms, kernels, and honest validation}

A full formalization should record the uses of classical logic, set
choice, global choice, and any ambient model-existence assumptions.
The common-prefix comparison should not acquire an ETR assumption merely
because abstract class order comparison sometimes uses one.
The inaccessible-rank example is conditional on that ambient cardinal,
whereas the bounded singular theorem needs no inaccessible cardinal.

The present package includes exact executable tests of finite shadows:
sign coding, pair coding, the scheduler, the common-prefix formula, and
finite adjacency. These are useful implementation checks and
counterexample detectors. They are not kernel-checked proofs of any
infinite or class theorem, and no supplied file is advertised as such.

% ---- hand chunk formal_after

\merge\ The four plans agree on a first milestone that is independent of
surreal arithmetic: the predecessor or common-prefix comparison of
well-orders on a small carrier, prefix-freeness and the
successor-cardinal obstruction, then bounded-support prefix completion and
small-family interpolation, reusing the existing small-cut interface. This
is a proposal recorded for the collection's formalization work; no such
module exists.

% ---- hand chunk questions_head@1
\cleardoublepage\phantomsection
\addcontentsline{toc}{part}{Questions, conclusions and status}
\part*{Questions, conclusions and status}
\section{Research questions}\label{swo:sec:questions}

\merge\ The sources' questions are printed per source in the following
subsections, as proposed continuations, not as a census of open problems.
Source~12's question on the weak-choice strength (\cref{swo:st:n14.5}) is
answered by \cref{swo:cf:choice} and printed as an answered remark. Source~11's
\cref{swo:n16.7} and source~12's \cref{swo:st:n14.3} ask related questions on
invariants at cutoffs; source~11's \cref{swo:n16.2} and source~12's
\cref{swo:st:n14.6} on minimal axioms; source~11's \cref{swo:n16.3} and
source~12's \cref{swo:st:n14.4} on proper-class blocks; source~11's
\cref{swo:n16.4} and source~12's \cref{swo:st:n14.7} on automorphisms; and
source~11's \cref{swo:n16.8} and source~12's \cref{swo:st:n14.8} on the class
realization. Source~12's \cref{swo:st:n14.1} continues the reals report's
question ``Other ground orders'', which \cref{swo:st:thm:transition} answers
in part.

% ---- source 11, lines 2385-2596
\subsection{Source 11's research questions and topics}\label{swo:sec:research}

The questions below are proposed extensions of the present
development. Except where a source is explicitly named, they
are not asserted to be established open problems in the
literature. Several ask for a stronger classification than
the proved results supply.

\subsubsection{A precise description of the represented cuts}

Fix the dense ordinary class \(\Pbd\cong\No\). Every
set-like global enumeration determines a lower cut of
\(\Pbd\), and so does every injective nonexhaustive
\(\Ord\)-word. The omitted-range example shows that the
first collection does not fill every cut generated by the
second.

\begin{question}\label{swo:n16.1}
Characterize, using the prefix tree and its numerical
labels, exactly which class cuts of \(\Pbd\) arise from
injective \(\Ord\)-words, and which of these words are
exhaustive.
\end{question}

A useful first target is a necessary-and-sufficient
compatibility condition on the intersections of a cut with
prefix cylinders. The challenge is to separate missing
labels, impossible repeated labels, and cuts that escape
every coherent branch. This should first be stated for a
fixed uniform class cut, before quantifying over an entire
hyperclass of cuts.

\subsubsection{The embedding relationship between the two global orders}

The set-like fragment and binary class words are
equimorphic, while the unrestricted order has additional
convex fibers. Theorem~\ref{swo:thm:universe-nonembedding}
rules out a reverse embedding in the full universe
interpretation. Theorem~\ref{swo:thm:km-nonembedding} also
rules out a definable uniform reverse embedding in
\(\KM+\GC\).

\begin{question}\label{swo:n16.2}
What is the weakest comprehension scheme over \(\GBC\)
that proves the uniform nonembedding result? Can the
inverse-evaluation step be replaced by an elementary
argument, or is a genuine strengthening necessary?
\end{question}

A precise project should specify the syntactic complexity
of the proposed transformation and of the comprehension
used to invert its code assignment. In the bounded
setting one can instead compare the density and
cellularity of all well-orders of \(\No_{<\kappa}\)
with their minimum-length stratum, especially when
\(\kappa\) is not inaccessible. The full-universe case
is already settled above; the remaining questions
concern strength and wider hypotheses.

\subsubsection{Iterated residual decomposition}

One extraction of the canonical initial \(\Ord\)-part
requires only elementary class comprehension. A residual
proper class has another such part, suggesting a hierarchy
of successive blocks.

\begin{question}\label{swo:n16.3}
What is the exact recursion strength needed to organize
iterated residual decompositions uniformly, and what
class order types can occur as their block-index orders?
\end{question}

An initial project could analyze finitely many blocks and
specified set-length constructions, recording where the
recursion state becomes a class. The general case should
be compared with known \(\ETR\) results on class
well-order comparability, rather than assumed to follow
from the one-step theorem.

\subsubsection{Recovering the prefix structure from the bare order}

The definitions exhibit a labeled tree of prefixes.
An automorphism of the bare lexicographic order need not
obviously preserve its displayed labels or cylinders.

\begin{question}\label{swo:n16.4}
Which prefix cones, canonical-core fibers, and finite
permutation blocks are definable from the linear order
alone? What extra predicates make the reconstruction
unique?
\end{question}

Adjacency directly detects finite stretches, but a
singleton fiber can be exhaustive or can omit one label.
Thus local block size alone cannot recover set-likeness.
An analysis of the automorphism group may clarify what
the order remembers about the original enumeration.

\subsubsection{The bounded-support choice of baseline}

Every baseline enumeration gives an ordinary class
\(\Pbd\) abstractly isomorphic to \(\No\). If two
baselines differ on a proper class of positions, their
bounded-support subfamilies can be disjoint.

\begin{question}\label{swo:n16.5}
How do the two dense classes sit inside the same
\(\Hsl(\No)\), and when does an automorphism of that
higher-order order carry one onto the other?
\end{question}

The abstract isomorphism of the two dense orders does
not automatically extend to all global enumerations.
The represented-cut problem is likely relevant. A
restricted first case is a baseline change obtained by
an order automorphism of the numerical surreal line.

\subsubsection{Sharper birthday bounds and minimal codes}

The explicit code gives a usable bound
\[
 \bd(E(f))=\sum_{i<\dom f}
              \bigl(1+2\cdot\bd(f(i))+2\bigr).
\]
It is designed for a transparent proof, not asserted
to be optimal.

\begin{question}\label{swo:n16.6}
Which birthday bounds are optimal for uniform embeddings
of injective surreal words? Can shorter prefix-free
coefficient codes preserve useful simplicity relations?
\end{question}

Regular cutoffs permit same-stage compression; singular
cutoffs show why a cardinal bound alone is insufficient.
A sharpened result should specify whether it preserves
only numerical order, prefix extension, simplicity, or
all of these.

\subsubsection{Saturation of more general bounded alphabets}

If \(X\) is an \(\eta_\kappa\) order of cardinality
\(\kappa\), its minimal-length permutation order inherits
\(\eta_\kappa\). Birthday fragments often have cardinality
\(2^{<\kappa}\), which may exceed \(\kappa\).

\begin{question}\label{swo:n16.7}
Determine the exact cut spectrum and point characters of
\(\WO_\lambda(\No_{<\kappa})\), where
\(\lambda=2^{<\kappa}\), without assuming
\(2^{<\kappa}=\kappa\).
\end{question}

The issue combines the size of a used prefix with the
number and birthdays of possible numerical separators.
The real-alphabet formulas in the precursor report
cannot settle this case by substitution.

\subsubsection{Class extensions and available global orders}

Two models may have the same sets and hence the same
individual surreal numbers, but different available
classes and different global well-orders.

\begin{question}\label{swo:n16.8}
Which features of the lexicographic theory are invariant
under extensions of the class part, and how does the
represented-cut collection change?
\end{question}

Uniform formulas such as first divergence remain the
same for the existing relations. Existence of additional
enumerations changes the larger totality. This offers a
way to distinguish properties of the labeled surreal
order from properties of the surrounding class theory.

\subsubsection{A verified first milestone for ProveIt}

\begin{question}\label{swo:n16.9}
Can the explicit word compression and the direct
predecessor comparison be formalized first, with a small,
reusable interface independent of surreal arithmetic?
\end{question}

These are concrete, finite-dependency targets:
ordinal concatenation, numerical first difference,
injectivity, and well-order predecessor sets. The next
milestone is the general successor-cardinal obstruction,
followed by bounded-support completion and the
small-family interpolation theorem. This ordering keeps
the central algorithms separate from a formalization of
an entire higher-order set theory.

\subsubsection{Uniform mathematical constructions indexed by well-orders}

The Kanovei--Shelah method suggests using a definable
linear order of codes to index constructions without
selecting one privileged code. The surreal setting
raises additional smallness conditions.

\begin{question}\label{swo:n16.10}
Which finite-support or bounded-support constructions
can be indexed by a uniform class family of global
well-orders, and which require a genuine higher-order
index object?
\end{question}

A first investigation should specify its support sets,
evaluation maps, and quantifier domains before adding
algebra or model theory. The diagonal theorem prevents
an exhaustive class-code shortcut, but leaves many
useful uniform families available.

% ---- hand chunk q11_after

\merge\ Status of source~11's questions after the merge: all remain open.
Source~12's dichotomy (\cref{swo:st:thm:transition}) computes the binary
coding length of $\WO_\mu(\No_{<\kappa})$ in every case, but not the cut
spectrum or the point characters that \cref{swo:n16.7} asks for.

% ---- source 12, lines 1447-1549
\subsection{Source 12's research questions and prioritized program}
\label{swo:st:sec:questions}

The following are proposed continuations. Except where an older question
is explicitly attributed, they are not claimed to be recognized open
problems in the literature.

\begin{question}[All fixed enumeration types]\label{swo:st:n14.1}
Determine $\bl(\WO_\alpha(\No_{<\kappa}))$ for every ordinal $\alpha$ of
cardinality $2^{<\kappa}$, not only the minimal type. Can the answer be
expressed by a precise ordinal-sum invariant of admissible entry-code
lengths, together with the finite terminal permutation factor?
\end{question}
The singular theorem identifies the first obstruction that any such
formula must capture. A useful initial target is
$\alpha=\mu\cdot2$ and then $\mu\cdot\gamma+n$ for finite $n$.

\begin{question}[Noncardinal birthday bounds]\label{swo:st:n14.2}
For an arbitrary infinite ordinal $\theta$, classify the minimal
permutation order of $\No_{<\theta}$. Which features of the ordinal tail of
$\theta$, beyond $|\theta|$ and $|\No_{<\theta}|$, survive in its coding
rank and local character?
\end{question}
The cone argument for an initial cardinal relies on closure under
$i+1+\alpha$ below the cutoff. At a noncardinal cutoff that closure can
fail and should be replaced by a stratified family of cones.

\begin{question}[Isomorphism, not merely equimorphism]\label{swo:st:n14.3}
Classify the order topology and isomorphism invariants of
$\WO_\mu(\No_{<\kappa})$. At regular cutoffs, and at singular strong limits,
what are its density, cellularity, cut-character spectra, and degrees of
homogeneity? Which distinguish it from other orders of the same binary
rank?
\end{question}
The binary cube itself cannot be the literal answer: it has endpoints
and adjacent pairs, while the surreal permutation order has different
local behavior. A dense suborder equimorphic with a cube can retain
substantial information not measured by rank.

\begin{question}[The hierarchy of proper-class blocks]\label{swo:st:n14.4}
Study the subpredicates of $\Hall(\No)$ with types
$\Ord\cdot\gamma+\alpha$ for fixed set ordinals $\gamma,\alpha$.
Give explicit, size-safe embedding and nonembedding theorems between
these subpredicates and the set-like order.
\end{question}
The inaccessible-model calculation provides concrete bounded targets and
shows that the full order is strictly more complex than the set-like
order there. Extending binary fusion to uniform transformations on
class predicates may require stronger comprehension; that dependence
should be proved, not assumed away.

\begin{remark}[Source 12's question ``The exact weak-choice strength'', answered]\label{swo:st:n14.5}
\emph{Source 12's text:} Remove set-level AC from the base theory in \cref{swo:st:thm:choice}. What
principle is equivalent to a class well-order of the sign-sequence
surreals then? Compare it with coherent well-ordering of power sets of
ordinals and with uniform ordinal presentations of transitive closures.

\merge\ \emph{Answered (2 October 2026).} Source 11's \cref{swo:cf:choice} proves the
equivalence over GB with no set-level choice: a class well-order of the
sign-sequence surreals is equivalent to global choice. Its proof is
precisely a coherent well-ordering of the power sets $V_{\beta+1}$,
obtained by set-valued recursion from the given well-order on the sign
sequences of each length $\theta$; see \cref{swo:rem:answered-choice}.
\end{remark}
This is separate from the question whether surreal universality for
class linear orders itself implies global choice, raised by Hamkins
\cite{HamkinsUnivChoice}. The present equivalence concerns a well-order of
all surreal labels and does not settle that universality question.

\begin{question}[Elementary common-prefix theory]\label{swo:st:n14.6}
Formalize \cref{swo:st:thm:fulllex} over a precisely chosen weak class theory.
How little comprehension and class well-foundedness are needed? Can one
characterize the same relation directly by a formula avoiding explicit
mention of the common-prefix class?
\end{question}
The definition already expands to elementary quantifiers. A mechanized
proof would make the separation from abstract class comparability
particularly transparent and yield a reusable theorem for other carriers.

\begin{question}[Automorphisms and finite condensation]\label{swo:st:n14.7}
Determine natural automorphism groups of the full lexicographic order,
and of the order after collapsing finite residual blocks. Which
operations arise from numerical automorphisms or permutations of the
underlying surreal carrier, and which genuinely act on enumeration
structure?
\end{question}
Finite-tail blocks are now characterized intrinsically by adjacency.
That makes condensation a plausible first invariant, but no complete
classification of its quotient is claimed here.

\begin{question}[Dependence on the class realization]\label{swo:st:n14.8}
Fix a transitive set universe $M$ and vary its admissible class collection
$\mathcal X$. Which external order invariants of the admitted global
well-orders detect this change? Can the order reconstruct useful
information about available truth predicates or class choice functions?
\end{question}
The diagonal obstruction persists internally, while external cardinality
can vary. This makes the distinction between full and Henkin semantics a
source of mathematical questions rather than merely a formal warning.

\begin{question}[Compatibility between birthday cutoffs]\label{swo:st:n14.9}
Construct and analyze coherent embeddings from the well-order spaces at
smaller birthday cutoffs into those at larger cutoffs. Which choices of
reserve labels and tail enumerations are functorial, and which preserve
finite condensation, coding ranks, or canonical cylinders?
\end{question}
Naively extending every enumeration by an arbitrary tail can change
minimal length, and fixed choices need not be compatible across three
cutoffs. A coherence theorem would make universe-lifting arguments much
more useful in a proof assistant.

% ---- hand chunk q12_after

\merge\ Status of source~12's questions after the merge: \cref{swo:st:n14.5} is
answered (it is printed above as a remark). \Cref{swo:st:n14.1} continues the
reals report's question ``Other ground orders'' for the surreal cutoffs, whose
minimal type \cref{swo:st:thm:transition} settles. The others remain open.

% ---- hand chunk conclusions_head@1
\section{Source 12's conclusion}\label{swo:sec:conclusions}

\merge\ Source~11 has no conclusion section; its summary is
\cref{swo:sec:status}.

% ---- source 12, lines 1551-1577
\subsection{Source 12's conclusion}

There is a rigorous full lexicographic order of all global surreal
well-orders, provided it is understood as a relation on class objects,
not as an ordinary class containing those objects. Its comparison needs
no universal enumeration of abstract class order types: the largest
common \emph{labelled} prefix is elementary and sufficient.

The set-like suborder has a particularly uniform geometry. It fills
set-indexed cuts, contains every class linear order even locally, and is
equimorphic with binary classes on $\Ord$. The unrestricted order adds
new proper-class prefix phenomena, including finite adjacent blocks.
The diagonal theorem prevents either construction from being exhausted
by a same-level class of set codes.

At set-sized birthday bounds the theory becomes quantitative. The full
well-order space of a carrier of cardinality $\mu$ has exact ordinal
spectrum below $\mu^+$ and binary rank $\mu^+$. The minimal surreal slice
has rank $\mu$, except at singular strong-limit cutoffs, where it has
rank $\kappa\cdot\kappa$. This is an order-theoretic effect of cofinal
blocks of code lengths, invisible to cardinal counting alone.

The inaccessible-rank interpretation shows exactly how these facts can
coexist: internal classes and the external collection of all such
classes occupy different levels. Respecting those levels removes the
apparent paradox without weakening the embedding and universality
questions into trivialities.

% ---- source 11, lines 2598-2658
\section{Results, assumptions, and verification status}
\label{swo:sec:status}

\begin{center}
\small
\begin{tabular}{@{}>{\raggedright\arraybackslash}p{0.54\textwidth}>{\raggedright\arraybackslash}p{0.38\textwidth}@{}}
\toprule
Result & Foundation used here\\
\midrule
Fixed-alphabet cardinality and ordinal spectrum &
\(\ZFC\); existing repository theorem\\[3pt]
Explicit set-length word compression &
\(\ZF\), as a definable class construction\\[3pt]
Local word characters and set-cut arguments &
Surreal set cuts; \(\ZFC\) for cardinal-valued characters\\[3pt]
Any class well-order of \(\No\) implies global choice &
\(\GB\) with \(\ZF\) set part\\[3pt]
Raw comparison, canonical initial part, finite-jump criterion &
\(\GB\), with the given class orders\\[3pt]
Bounded-support completion and interpolation &
\(\GBC\), with a fixed global enumeration\\[3pt]
\(\Pbd\cong\No\) and class-order universality &
\(\GBC\), by set-valued back-and-forth\\[3pt]
Binary class coding and no exhaustive uniform evaluator &
\(\GBC\), interpreted with class parameters\\[3pt]
Unrestricted order does not embed into its set-like fragment &
Full inaccessible universe; or \(\KM+\GC\) for definable uniform maps\\[3pt]
All higher-order collections treated as actual objects &
An explicitly chosen higher-order or containing-universe setting\\
\bottomrule
\end{tabular}
\end{center}

The proof dependencies are short. Set cuts yield fresh
coefficients after deleting a set. Fresh coefficients and
bounded-support completion yield strong interpolation.
Strong interpolation yields density, local cofinality
statements, the \(\No\)-isomorphism, and cut coding.
Independently, the predecessor comparison yields convex
cones, the canonical initial-part decomposition, and the
finite-residual jump theorem. Pair swaps supply both
positive binary embeddings and the diagonal obstruction.

The bundled Python script checks finite instances of
the two comparison rules and of the sign encoding. It
checks 465 coefficient-code comparisons, 79,800
word-code comparisons, and 266,272 pairs of finite
permutations, with the predecessor rule checked in both
directions and the adjacency criterion compared with the
finite lexicographic order. These checks help detect a
reversed delimiter or an incorrect finite comparison.
They do not verify an ordinal induction, a class
comprehension argument, or a cardinal nonembedding
theorem.

The article contains full mathematical proofs for the
claims it promotes to theorems. It does not claim that
the new results have been checked in Lean, that the
whole repository has been audited, or that every
research question listed here is a previously published
open problem.

% ---- hand chunk ledgers_head

\merge\ The other sources' ledgers follow; each records the hypotheses its
source used, not that each is necessary.

% ---- source 12, lines 1580-1606
\subsection{Source 12's proof-dependency and assumption ledger}

\begin{longtable}{@{}p{0.32\textwidth}p{0.24\textwidth}p{0.36\textwidth}@{}}
\toprule
Result & Stated setting & Essential input\\
\midrule
\endhead
Common labelled prefix and full lexical comparison & GBC; comparison is
elementary & Existing well-orders; class minima; ambient linear order.\\
Global-choice equivalence & GB + set AC & Ordinal coding of transitive
closures; restriction of a surreal well-order to a set of codes.\\
No exhaustive class-coded family & GBC & A baseline $\Ord$-bijection;
elementary diagonal pair orientation.\\
Set-like cut interpolation & GBC & Set-sized discrepancy bound; surreal
small-cut filling; prefix extension.\\
Local class universality & GBC & Set-valued forth and back-and-forth
recursions; global enumeration.\\
Finite-tail adjacency & GBC & Label transpositions; finite residual
cylinders; no infinite order well-ordered in both directions.\\
Set-sized full spectra & ZFC & Pair coding; regularity of $\mu^+$;
cylinder-fusion nonembedding.\\
Minimal-slice singular transition & ZFC & Cone partition; scheduler;
ordinal block sums; regularity or singular cofinality as appropriate.\\
Full-class rank model & Ambient ZFC + an inaccessible $\kappa$ & Closure
of $V_\kappa$; all external classes; regularity for descending sequences.\\
\bottomrule
\end{longtable}

% ---- hand chunk appendix_head
\appendix

% ---- source 11, lines 2661-2700
\section{Source crosswalk and reproducibility}\label{swo:app:sources}

The immutable sources used in the targeted audit are
recorded in the accompanying source audit. For convenient
verification, the most relevant mappings are:

\begin{longtable}{@{}>{\raggedright\arraybackslash}p{0.39\textwidth}>{\raggedright\arraybackslash}p{0.53\textwidth}@{}}
\toprule
Source & Relevant content\\
\midrule
\endhead
\cite{ProveItReals} &
Labels \code{lwo:lem:prefix-free},
\code{lwo:lem:cylinder}, \code{lwo:prop:cardinality},
\code{lwo:thm:no-long},
\code{lwo:thm:ordinal-spectrum}; bounded sign orders in
\code{lwo:sp:thm:surreals}.\\[4pt]
\cite{ProveItFoundations} &
Set-coded surreals and bounded birthdays; unrestricted
cut obstruction; class-quantification and universe
distinctions; set-valued recursion.\\[4pt]
\cite{ProveItSigns} &
Actual sign carrier, ordinal embedding, smallness
obstruction, and bounded fragments.\\[4pt]
\cite{ProveItCuts} &
Small cut operation, separator properties, and
birthday bounds.\\[4pt]
\cite{ProveItVector} &
A nearby use of a set-like global well-order and
ordinary class-parameter recursion.\\
\bottomrule
\end{longtable}

To compile the delivered standalone source, run
\code{latexmk -pdf surreal\_well\_orders.tex} from its directory.
All article sections and the bibliography are included in that
single file; no separate bibliography processor is needed.
Run \code{python3 finite\_checks.py} for the finite checks.
The supplied PDF was compiled from the included sources
and inspected after rendering.

% ---- hand chunk app11_after

\merge\ \emph{Delivery names.} The paragraph above describes source~11's
delivered package. Its \code{surreal\_well\_orders.tex} is shipped as
\path{article.tex}, which this merged text replaces; its PDF is not shipped
(\path{article.pdf} is a build of this text); its \code{finite\_checks.py} and
\code{finite\_checks\_results.txt} are
\path{code/11-raw-orders-finite_checks.py} and
\path{data/11-raw-orders-finite_checks_results.txt}; the ``accompanying source
audit'' is \path{11-raw-orders-repository_audit.md}. Build this report with
\code{latexmk -pdf -interaction=nonstopmode -halt-on-error article.tex}, and
rerun the checks as in \cref{swo:app:checks}. The same applies to the
reproducibility records of the other sources that follow.

% ---- source 12, lines 1608-1623
\subsection{Source 12: reproducibility and inspection record}

The package contains \texttt{article.tex}, the compiled
\texttt{article.pdf}, a build script, an exact finite-check program,
its recorded JSON output, and a research-status and source-audit file.
The Python program uses only the standard library. The recorded run
passed all 180,905 finite comparison cases. Its comparisons are exact; there are no floating-point tolerances or inferred transfinite
results. The PDF was compiled and rendered for visual inspection.

The repository audit directly inspected the merged real-well-order
README and the sign-sequence simplicity Lean source at the stated commit.
It did not independently compile the repository or verify the predecessor
report's proofs. The needed predecessor patterns have therefore been
reproved in the present manuscript. Literature checking was focused on
surreal foundations, Kanovei--Shelah indexing, and class recursion and
well-orders; it was not sufficient to certify historical originality.

% ---- hand chunk app_checks@1
\section{The finite checks as shipped}\label{swo:app:checks}

\merge\ The programs of sources~11 and~12 use only the Python standard library
and check finite instances; they verify no transfinite, class-theoretic or
cardinal statement. Rerun on copies on 2 October 2026 (Python~3.14.4,
Windows), each reproduced its recorded output exactly apart from line endings
(both write CRLF on Windows; compare with \code{diff --strip-trailing-cr}):
\path{code/11-raw-orders-finite_checks.py} prints the content of
\path{data/11-raw-orders-finite_checks_results.txt} (465 sign-code pairs,
79,800 word-code pairs, 266,272 pairs of permutations, 867 adjacent), and
\path{code/12-singular-finite_checks.py --output} writes the content of
\path{data/12-singular-finite_checks.json} (180,905 cases). Always pass an
output path outside the report directory. The delivered
\path{code/12-singular-build.sh} changes into \path{code/} and then calls
\code{code/finite\_checks.py}; it builds nothing in this layout.

% ---- hand chunk app_abstracts_head
\section{The other sources' title pages and abstracts}\label{swo:app:abstracts}

\merge\ The results these summaries announce are printed in the body; the
summaries are kept for the record.

\subsection*{Source 12}
\paragraph{Title-page summary.}

% ---- source 12, lines 86-111
\textbf{Main distinction.} An arbitrary class well-order of $\No$ need not
have order type $\Ord$. Nevertheless, two such well-orders have an
\emph{elementarily definable largest common labelled initial segment},
possibly a proper class. This supplies a lexicographic comparison in GBC
without invoking general class well-order comparability or ETR.

\medskip
\textbf{Principal quantitative result.} Put
$\No_{<\kappa}=\{s:\alpha\to\{-,+\}:\alpha<\kappa\}$ and
$\mu=|\No_{<\kappa}|=2^{<\kappa}$, for an infinite initial cardinal $\kappa$.
The minimal-length permutation order satisfies
\[
\WO_\mu(\No_{<\kappa})\equim
\begin{cases}
\BB_{\kappa\cdot\kappa},&\mu=\kappa\text{ and }\cf\kappa<\kappa,\\
\BB_\mu,&\text{otherwise},
\end{cases}
\]
where $\BB_\gamma=(2^\gamma,<_{\lex})$ and $\equim$ means mutual
order-embeddability, not isomorphism.

\medskip
\textbf{Status.} This is an AI-assisted, unrefereed research manuscript.
Proofs are supplied for the stated results; historical priority is not
claimed. The accompanying exact finite tests do not verify the transfinite
arguments. No new Lean formalization is claimed.

% ---- hand chunk abs12_mid

\paragraph{Abstract.}

% ---- source 12, lines 117-143
We study the lexicographic order of well-orderings of the surreal numbers,
carefully distinguishing bounded sets of surreals, set-like global
well-orders, and arbitrary global class well-orders. For the full class
problem, ordinal-indexed enumerations are insufficient. We construct the
largest common labelled initial segment of two well-orders directly from
their relations, prove that comparison of their next elements defines a
linear order, and show that this construction needs only elementary class
comprehension once the well-orders exist. Over GB with set choice, the
existence of any class well-order of the surreals is equivalent to global
choice. The collection of all these well-orders is not itself an ordinary
class of GBC or Kelley--Morse theory; a diagonal theorem rules out any
uniform exhaustive coding by a class of sets.

For set-like global well-orders we prove density, interpolation for
set-indexed cuts, local universality for class linear orders, and mutual
embeddability with the lexicographic comparison of binary classes on
$\Ord$. For unrestricted global well-orders we classify adjacent pairs:
they arise precisely from finite residual permutations after a common,
necessarily proper-class, prefix. At a birthday cutoff $\kappa$, we obtain
an exact regular/singular dichotomy for the binary coding rank of
minimal-length enumerations. We also establish the exact ordinal and
binary-cube embedding spectra for the full set-sized order. Finally, a
full-class model based on an inaccessible rank explains exactly how the
bounded and global theories fit together, and provides a rigorous example
in which the full order and the set-like order have the same cardinality
and ordinal spectrum but different coding ranks. Explicit proof-assistant
interfaces and further research questions are included.

% ---- hand chunk crosswalk_head
\section{Crosswalk of the sources' numbered statements}\label{swo:app:crosswalk}

\merge\ Each row gives a numbered statement of a source (numbered as in that
source's own PDF), and where it is printed or recorded here. A note records a
statement proved here by the same argument; a section labelled ``route''
prints a different argument. Labels are the source's own with its prefix, or
\code{swo:}\emph{sub}\code{:n}\emph{number} for unlabelled statements.

% ---- generated crosswalk
\subsection{Source 11 (base)}
\begingroup\small
\begin{longtable}{@{}p{0.11\textwidth}p{0.47\textwidth}p{0.34\textwidth}@{}}
\toprule
No. & Statement in source 11 & Printed here as\\
\midrule\endhead
3.1 & Definition & \Cref{swo:n3.1}, p.~\pageref{swo:n3.1}\\
3.2 & Lemma (Deleting a set) & \Cref{swo:lem:delete}, p.~\pageref{swo:lem:delete}\\
3.3 & Proposition & \Cref{swo:prop:universal}, p.~\pageref{swo:prop:universal}\\
4.1 & Theorem (Cardinality and exact ordinal spectrum) & \Cref{swo:sw:setspectrum}, p.~\pageref{swo:sw:setspectrum}\\
4.2 & Proposition (Exactly where jumps occur) & \Cref{swo:sw:adjacency}, p.~\pageref{swo:sw:adjacency}\\
5.1 & Corollary & \Cref{swo:sw:cutoffspectrum}, p.~\pageref{swo:sw:cutoffspectrum}\\
5.2 & Lemma & \Cref{swo:sw:stageeta}, p.~\pageref{swo:sw:stageeta}\\
6.1 & Theorem (Explicit sign coding) & \Cref{swo:sw:code}, p.~\pageref{swo:sw:code}\\
6.2 & Proposition (Birthday bounds) & \Cref{swo:sw:codebounds}, p.~\pageref{swo:sw:codebounds}\\
7.1 & Theorem (Exact local character) & \Cref{swo:sw:character}, p.~\pageref{swo:sw:character}\\
7.2 & Corollary & \Cref{swo:sw:notiso}, p.~\pageref{swo:sw:notiso}\\
8.1 & Theorem (Word-cut classification) & \Cref{swo:sw:cutclassification}, p.~\pageref{swo:sw:cutclassification}\\
8.2 & Corollary (Successor words) & \Cref{swo:sw:successors}, p.~\pageref{swo:sw:successors}\\
9.1 & Proposition & \Cref{swo:sw:permutationeta}, p.~\pageref{swo:sw:permutationeta}\\
10.1 & Theorem (Existence and global choice) & \Cref{swo:cf:choice}, p.~\pageref{swo:cf:choice}\\
11.1 & Lemma (Unique first disagreement) & \Cref{swo:cf:first}, p.~\pageref{swo:cf:first}\\
11.2 & Theorem (Raw lexicographic order) & \Cref{swo:cf:lex}, p.~\pageref{swo:cf:lex}\\
11.3 & Proposition (Convex prefix cones) & \Cref{swo:cf:cones}, p.~\pageref{swo:cf:cones}\\
12.1 & Theorem (Canonical initial part) & \Cref{swo:cf:core}, p.~\pageref{swo:cf:core}\\
12.2 & Example (Two proper blocks) & \Cref{swo:cf:two_blocks}, p.~\pageref{swo:cf:two_blocks}\\
12.3 & Proposition (Convex fibers) & \Cref{swo:cf:fibers}, p.~\pageref{swo:cf:fibers}\\
12.4 & Remark (Singleton fibers and iteration) & \Cref{swo:cf:singleton}, p.~\pageref{swo:cf:singleton}\\
13.1 & Lemma (Extreme enumerations) & \Cref{swo:cf:extreme}, p.~\pageref{swo:cf:extreme}\\
13.2 & Theorem (Exact adjacency criterion) & \Cref{swo:cf:adjacency}, p.~\pageref{swo:cf:adjacency}\\
13.3 & Example (An adjacent pair of type \(\Ord+2\)) & \Cref{swo:cf:ordtwo}, p.~\pageref{swo:cf:ordtwo}\\
14.1 & Definition & \Cref{swo:gl:uniform}, p.~\pageref{swo:gl:uniform}\\
14.2 & Definition & \Cref{swo:gl:bounded-codes}, p.~\pageref{swo:gl:bounded-codes}\\
14.3 & Lemma (Completion of an arbitrary set prefix) & \Cref{swo:gl:completion}, p.~\pageref{swo:gl:completion}\\
14.4 & Theorem (Prefix interpolation) & \Cref{swo:gl:interpolation}, p.~\pageref{swo:gl:interpolation}\\
14.5 & Corollary & \Cref{swo:gl:pbd-no}, p.~\pageref{swo:gl:pbd-no}\\
14.6 & Corollary (The class of realizations of a set cut) & \Cref{swo:gl:realizers}, p.~\pageref{swo:gl:realizers}\\
14.7 & Proposition & \Cref{swo:gl:universality}, p.~\pageref{swo:gl:universality}\\
14.8 & Proposition (Prefix cylinders) & \Cref{swo:gl:cylinders}, p.~\pageref{swo:gl:cylinders}\\
14.9 & Proposition (Approximants at every point) & \Cref{swo:gl:local}, p.~\pageref{swo:gl:local}\\
14.10 & Corollary (Set saturation is not Dedekind completeness) & \Cref{swo:gl:no-suprema}, p.~\pageref{swo:gl:no-suprema}\\
14.11 & Theorem & \Cref{swo:gl:binary}, p.~\pageref{swo:gl:binary}\\
14.12 & Corollary & \Cref{swo:gl:cuts}, p.~\pageref{swo:gl:cuts}\\
14.13 & Theorem (No exhaustive uniform class of codes) & \Cref{swo:gl:diagonal}, p.~\pageref{swo:gl:diagonal}\\
14.14 & Theorem (An omitted-range cut) & \Cref{swo:gl:omitted}, p.~\pageref{swo:gl:omitted}\\
14.15 & Example (Two omitted points produce adjacent raw orders) & \Cref{swo:gl:omitted-two}, p.~\pageref{swo:gl:omitted-two}\\
14.16 & Corollary (Exactly which raw set cuts admit bounded completions) & \Cref{swo:gl:raw-cut}, p.~\pageref{swo:gl:raw-cut}\\
15.1 & Theorem (Nonembedding in a full universe) & \Cref{swo:thm:universe-nonembedding}, p.~\pageref{swo:thm:universe-nonembedding}\\
15.2 & Theorem & \Cref{swo:thm:km-nonembedding}, p.~\pageref{swo:thm:km-nonembedding}\\
16.1 & Research question & \Cref{swo:n16.1}, p.~\pageref{swo:n16.1}\\
16.2 & Research question & \Cref{swo:n16.2}, p.~\pageref{swo:n16.2}\\
16.3 & Research question & \Cref{swo:n16.3}, p.~\pageref{swo:n16.3}\\
16.4 & Research question & \Cref{swo:n16.4}, p.~\pageref{swo:n16.4}\\
16.5 & Research question & \Cref{swo:n16.5}, p.~\pageref{swo:n16.5}\\
16.6 & Research question & \Cref{swo:n16.6}, p.~\pageref{swo:n16.6}\\
16.7 & Research question & \Cref{swo:n16.7}, p.~\pageref{swo:n16.7}\\
16.8 & Research question & \Cref{swo:n16.8}, p.~\pageref{swo:n16.8}\\
16.9 & Research question & \Cref{swo:n16.9}, p.~\pageref{swo:n16.9}\\
16.10 & Research question & \Cref{swo:n16.10}, p.~\pageref{swo:n16.10}\\
\bottomrule
\end{longtable}
\endgroup
\subsection{Source 12}
\begingroup\small
\begin{longtable}{@{}p{0.11\textwidth}p{0.47\textwidth}p{0.34\textwidth}@{}}
\toprule
No. & Statement in source 12 & Printed here as\\
\midrule\endhead
3.1 & Definition (Common labelled prefix) & \Cref{swo:st:n3.1}, p.~\pageref{swo:st:n3.1}\\
3.2 & Lemma (Maximal-prefix lemma) & \Cref{swo:st:lem:common}, p.~\pageref{swo:st:lem:common}\\
3.3 & Definition (Full lexicographic comparison) & \Cref{swo:st:n3.3}, p.~\pageref{swo:st:n3.3}\\
3.4 & Theorem (Full class lexicographic order) & \Cref{swo:st:thm:fulllex}, p.~\pageref{swo:st:thm:fulllex}\\
3.5 & Remark (Why this does not prove class order-type comparability) & \Cref{swo:st:n3.5}, p.~\pageref{swo:st:n3.5}\\
3.6 & Example (A disagreement after every ordinary ordinal) & \Cref{swo:st:n3.6}, p.~\pageref{swo:st:n3.6}\\
3.7 & Corollary (Agreement with enumeration lexicography) & \Cref{swo:st:n3.7}, p.~\pageref{swo:st:n3.7}\\
4.1 & Theorem (Choice equivalence) & \Cref{swo:st:thm:choice}, p.~\pageref{swo:st:thm:choice}\\
4.2 & Remark (The explicit set-choice hypothesis) & \Cref{swo:st:n4.2}, p.~\pageref{swo:st:n4.2}\\
5.1 & Theorem (No exhaustive class-coded family) & \Cref{swo:st:thm:diagonal}, p.~\pageref{swo:st:thm:diagonal}\\
5.2 & Corollary (No uniform set-code solution) & \Cref{swo:st:n5.2}, p.~\pageref{swo:st:n5.2}\\
6.1 & Lemma (Extending a set prefix) & \Cref{swo:st:lem:extend}, p.~\pageref{swo:st:lem:extend}\\
6.2 & Lemma (Avoiding a forbidden set) & \Cref{swo:st:lem:avoid}, p.~\pageref{swo:st:lem:avoid}\\
6.3 & Theorem (Set-indexed interpolation) & \Cref{swo:st:thm:etasl}, p.~\pageref{swo:st:thm:etasl}\\
6.4 & Proposition (Cofinality and local character) & \Cref{swo:st:prop:characters}, p.~\pageref{swo:st:prop:characters}\\
6.5 & Lemma (Surreal universality, with the size of the recursion exposed) & \Cref{swo:st:lem:universalNo}, p.~\pageref{swo:st:lem:universalNo}\\
6.6 & Lemma (Deleting a set does not change the numerical class order) & \Cref{swo:st:lem:delete}, p.~\pageref{swo:st:lem:delete}\\
6.7 & Theorem (Local class universality and cylinders) & \Cref{swo:st:thm:local}, p.~\pageref{swo:st:thm:local}\\
7.1 & Lemma (A self-delimiting numerical code) & \Cref{swo:st:lem:code}, p.~\pageref{swo:st:lem:code}\\
7.2 & Theorem (Binary-class equimorphism) & \Cref{swo:st:thm:binaryclass}, p.~\pageref{swo:st:thm:binaryclass}\\
8.1 & Definition (A prefix cylinder in the full order) & \Cref{swo:st:n8.1}, p.~\pageref{swo:st:n8.1}\\
8.2 & Proposition (Residual-order decomposition) & \Cref{swo:st:prop:residual}, p.~\pageref{swo:st:prop:residual}\\
8.3 & Theorem (Exact adjacency criterion) & \Cref{swo:st:thm:adjacency}, p.~\pageref{swo:st:thm:adjacency}\\
8.4 & Corollary (Proper-class prefixes are responsible) & \Cref{swo:st:n8.4}, p.~\pageref{swo:st:n8.4}\\
9.1 & Lemma (Prefix-freeness of exhaustive enumerations) & \Cref{swo:st:lem:prefixfree}, p.~\pageref{swo:st:lem:prefixfree}\\
9.2 & Lemma (Strict binary hierarchy) & \Cref{swo:st:lem:strictcube}, p.~\pageref{swo:st:lem:strictcube}\\
9.3 & Lemma (Characteristic-cut coding) & \Cref{swo:st:lem:cutcode}, p.~\pageref{swo:st:lem:cutcode}\\
9.4 & Lemma (Pair orientation coding) & \Cref{swo:st:lem:pairs}, p.~\pageref{swo:st:lem:pairs}\\
9.5 & Lemma (No $\mu^+$-long monotone chain) & \Cref{swo:st:lem:nochain}, p.~\pageref{swo:st:lem:nochain}\\
9.6 & Theorem (Cardinality and exact spectra) & \Cref{swo:st:thm:setspectrum}, p.~\pageref{swo:st:thm:setspectrum}\\
9.7 & Remark (What the rank does not say) & \Cref{swo:st:n9.7}, p.~\pageref{swo:st:n9.7}\\
10.1 & Lemma (Row scheduler) & \Cref{swo:st:lem:scheduler}, p.~\pageref{swo:st:lem:scheduler}\\
10.2 & Lemma (Surreal cone decomposition) & \Cref{swo:st:lem:cones}, p.~\pageref{swo:st:lem:cones}\\
10.3 & Theorem (Exact minimal-slice dichotomy) & \Cref{swo:st:thm:transition}, p.~\pageref{swo:st:thm:transition}\\
10.4 & Corollary (A complete cube test for the minimal slice) & \Cref{swo:st:n10.4}, p.~\pageref{swo:st:n10.4}\\
10.5 & Example (An unconditional singular example) & \Cref{swo:st:n10.5}, p.~\pageref{swo:st:n10.5}\\
10.6 & Example (Regular cutoffs) & \Cref{swo:st:n10.6}, p.~\pageref{swo:st:n10.6}\\
11.1 & Theorem (Exact external identification) & \Cref{swo:st:thm:model}, p.~\pageref{swo:st:thm:model}\\
14.1 & Research question (All fixed enumeration types) & \Cref{swo:st:n14.1}, p.~\pageref{swo:st:n14.1}\\
14.2 & Research question (Noncardinal birthday bounds) & \Cref{swo:st:n14.2}, p.~\pageref{swo:st:n14.2}\\
14.3 & Research question (Isomorphism, not merely equimorphism) & \Cref{swo:st:n14.3}, p.~\pageref{swo:st:n14.3}\\
14.4 & Research question (The hierarchy of proper-class blocks) & \Cref{swo:st:n14.4}, p.~\pageref{swo:st:n14.4}\\
14.5 & Research question (The exact weak-choice strength) & \Cref{swo:st:n14.5}, p.~\pageref{swo:st:n14.5}\\
14.6 & Research question (Elementary common-prefix theory) & \Cref{swo:st:n14.6}, p.~\pageref{swo:st:n14.6}\\
14.7 & Research question (Automorphisms and finite condensation) & \Cref{swo:st:n14.7}, p.~\pageref{swo:st:n14.7}\\
14.8 & Research question (Dependence on the class realization) & \Cref{swo:st:n14.8}, p.~\pageref{swo:st:n14.8}\\
14.9 & Research question (Compatibility between birthday cutoffs) & \Cref{swo:st:n14.9}, p.~\pageref{swo:st:n14.9}\\
\bottomrule
\end{longtable}
\endgroup

% ---- hand chunk bibliography

\begin{thebibliography}{99}
\small

\bibitem{ProveItReals}
VladimirReshetnikov/ProveIt repository,
\emph{Lexicographic orders of the well-orderings of the reals},
merged research report, 2 October 2026,
\path{SetTheory/Cardinals/docs/reports/ordinals-and-order-types/lexicographic-well-orderings-of-reals/}.
Source~11 read its article at revision \texttt{ae7c1e6aa}; sources~09 and~12
read its README at revision \texttt{63bb8b1d9}; source~08 did not read it
(see \cite{Batch76Source02}). The article has not changed since
\texttt{603f6ed31}.

\bibitem{Batch76Source02}
\emph{Lexicographic Orders of Well-Orderings of the Continuum: Ordinal
spectra, lexicographic complexity, finite-tail geometry, and self-absorption},
research manuscript prepared for Vladimir Reshetnikov, 2 October 2026,
24 pages; batch-76 manuscript~02, the base source~02 of \cite{ProveItReals}.
Source~08 cites it as a ``prior working manuscript in the user's library,
supplied there as \texttt{article(20261002-170841).pdf} and corresponding
\LaTeX{} source''; unrefereed.

\bibitem{ProveItFoundations}
VladimirReshetnikov/ProveIt repository,
\emph{Foundations for surreal and surcomplex mathematics}, research report,
\path{Algebra/SurrealNumbers/docs/foundations-and-computation/foundations/};
read by source~11 at revision \texttt{ae7c1e6aa}.

\bibitem{ProveItHset}
VladimirReshetnikov/ProveIt repository,
\emph{Birthday cutoffs and hereditary sets}, merged research report,
\path{Algebra/SurrealNumbers/docs/foundations-and-computation/birthday-cutoffs-and-hereditary-sets/}.

\bibitem{ProveItVector}
VladimirReshetnikov/ProveIt repository,
\emph{The surreal numbers as a real vector space},
research report,
\path{Algebra/SurrealNumbers/docs/surreal/real-vector-space-structure/};
read by source~11 at revision \texttt{ae7c1e6aa}.

\bibitem{ProveItSurrealREADME}
VladimirReshetnikov/ProveIt repository, surreal-number project README,
\path{Algebra/SurrealNumbers/README.md}; read by source~08 at revision
\texttt{ae7c1e6aa}.

\bibitem{ProveItSigns}
ProveIt contributors, \path{Algebra/SurrealNumbers/Surreal/Foundations/SignSequence.lean};
read by sources~11 and~08 at revision \texttt{ae7c1e6aa} (blob
\texttt{e2c093860f29b0b9ea312d50099fc4aa89a76a99}) and by source~09 at
\texttt{63bb8b1d9}.

\bibitem{ProveItCutFill}
ProveIt contributors, \path{Algebra/SurrealNumbers/Surreal/Foundations/SignSequenceCut.lean};
read by source~08 at revision \texttt{ae7c1e6aa} (blob
\texttt{c5c35afd98a5ad293687367ee9fb0244f94cc91c}).

\bibitem{ProveItCuts}
ProveIt contributors, \path{Algebra/SurrealNumbers/Surreal/Foundations/SignSequenceCutOperation.lean};
read by source~11 at revision \texttt{ae7c1e6aa} and by source~09 at
\texttt{63bb8b1d9}.

\bibitem{ProveItSimplicity}
ProveIt contributors, \path{Algebra/SurrealNumbers/Surreal/Foundations/SignSequenceSimplicity.lean};
read by source~12 at revision \texttt{63bb8b1d9} and by source~11 at
\texttt{ae7c1e6aa}. It proves minimum-birthday and simplest-separator
uniqueness; source~12 did not rebuild it.

\bibitem{Conway}
J.~H. Conway,
\emph{On Numbers and Games},
Academic Press, 1976; second edition, A~K Peters, 2001.

\bibitem{Gonshor}
Harry Gonshor,
\emph{An Introduction to the Theory of Surreal Numbers},
London Mathematical Society Lecture Note Series 110,
Cambridge University Press, 1986.
\href{https://doi.org/10.1017/CBO9780511629143}{doi:10.1017/CBO9780511629143}.

\bibitem{Alling}
N.~L. Alling,
Conway's field of surreal numbers,
\emph{Transactions of the American Mathematical Society} \textbf{287} (1985),
365--386.
\href{https://doi.org/10.1090/S0002-9947-1985-0766225-7}{doi:10.1090/S0002-9947-1985-0766225-7}.

\bibitem{BagayokoHoeven}
Vincent Bagayoko and Joris van der Hoeven,
\emph{Surreal substructures},
Fundamenta Mathematicae 266 (2024), 25--96.
\href{https://doi.org/10.4064/fm231020-23-2}{doi:10.4064/fm231020-23-2}.
\href{https://arxiv.org/abs/2305.02001}{arXiv:2305.02001}.

\bibitem{Ehrlich2012}
Philip Ehrlich,
\emph{The absolute arithmetic continuum and the unification
of all numbers great and small},
Bulletin of Symbolic Logic 18 (2012), no.~1, 1--45.
\href{https://doi.org/10.2178/bsl/1327328438}{doi:10.2178/bsl/1327328438}.

\bibitem{KanoveiShelah}
Vladimir Kanovei and Saharon Shelah,
\emph{A definable nonstandard model of the reals},
Journal of Symbolic Logic 69 (2004), no.~1, 159--164.
\href{https://doi.org/10.2178/jsl/1080938834}{doi:10.2178/jsl/1080938834}.
\href{https://arxiv.org/abs/math/0311165}{arXiv:math/0311165}; the index is
defined in Section~1, page~2 of the preprint.

\bibitem{RochaRangelMariano}
D. Rocha Rangel and H.~L. Mariano,
An algebraic (set) theory of surreal numbers, I, 2019.
\href{https://arxiv.org/abs/1911.12726}{arXiv:1911.12726}. Research preprint.

\bibitem{HamkinsChoice}
Joel David Hamkins,
\emph{Surreal and ordinal numbers}, author answer,
Mathematics Stack Exchange, 5 November 2010.
\url{https://math.stackexchange.com/a/9038}.

\bibitem{HamkinsGC}
Joel David Hamkins,
\emph{The global choice principle in G\"odel--Bernays set theory},
3 December 2014.
\url{https://jdh.hamkins.org/the-global-choice-principle-in-godel-bernays-set-theory/}.

\bibitem{HamkinsUniversal}
J.~D. Hamkins,
Answer to ``Universal order type'', MathOverflow, 6 March 2011,
\url{https://mathoverflow.net/questions/57597/universal-order-type}.
An author-given account of the forth construction for surreal universality.

\bibitem{HamkinsUnivChoice}
J.~D. Hamkins,
``Is the universality of the surreal number line a weak global choice
principle?'', MathOverflow, 7 January 2016,
\url{https://mathoverflow.net/questions/227849/}.
Cited as a historical question, not as a certified 2026 status claim.

\bibitem{GitmanHamkins}
Victoria Gitman and Joel David Hamkins,
\emph{Open determinacy for class games},
research paper.
\href{https://arxiv.org/abs/1509.01099}{arXiv:1509.01099}.

\bibitem{HamkinsOpenDet}
J.~D. Hamkins,
``Open determinacy for class games'',
author's research overview of joint work with V.~Gitman, 4 September 2015.
\href{https://jdh.hamkins.org/open-determinacy-for-class-games/}{Author's source page}.

\bibitem{GHHSW}
V. Gitman, J.~D. Hamkins, P. Holy, P. Schlicht, and K.~J. Williams,
``The exact strength of the class forcing theorem,''
\emph{Journal of Symbolic Logic} \textbf{85} (2020), 869--905.
\href{https://arxiv.org/abs/1707.03700}{arXiv:1707.03700}; version 2,
10 March 2021.

\bibitem{HamkinsClassForcing}
J.~D. Hamkins,
``The exact strength of the class forcing theorem,''
author's research overview of joint work with V.~Gitman, P.~Holy,
P.~Schlicht, and K.~J. Williams, 12 July 2017.
\href{https://jdh.hamkins.org/class-forcing-theorem/}{Author's source page}.

\bibitem{HamkinsRecursion}
J.~D. Hamkins, ``Second-order transfinite recursion is equivalent to
Kelley--Morse set theory over GBC,'' author's research exposition,
23 July 2017.
\href{https://jdh.hamkins.org/second-order-transfinite-recursion-is-equivalent-to-kelley-morse-set-theory/}{Author's webpage}.

\bibitem{HamkinsWoodin}
Joel David Hamkins and W.~Hugh Woodin,
\emph{Open class determinacy is preserved by forcing},
research paper, 2018.
\href{https://arxiv.org/abs/1806.11180}{arXiv:1806.11180}.
Especially Sections~2--3 and Theorem~6.

\bibitem{AntosFriedman}
Carolin Antos and Sy-David Friedman,
\emph{Hyperclass forcing in Morse--Kelley class theory},
research paper.
\href{https://arxiv.org/abs/1510.04082}{arXiv:1510.04082}.

\bibitem{LeanUniverses}
Lean language reference,
\emph{Universes}, official documentation, accessed October 2026.
\url{https://lean-lang.org/doc/reference/latest/The-Type-System/Universes/}.

\bibitem{MathlibOrdinal}
Mathlib contributors,
\emph{Mathlib.SetTheory.Ordinal.Basic}, official API documentation,
accessed October 2026.
\url{https://leanprover-community.github.io/mathlib4_docs/Mathlib/SetTheory/Ordinal/Basic.html}.

\end{thebibliography}
\end{document}

